\documentclass[11pt, a4paper]{article}
% GENERIERT von scripts/build_gesamtskript.py — nicht von Hand bearbeiten.
% Lehrtext kommt aus den Handouts, Aufgaben/Lösungen aus aufgaben-l*.tex / loesungen-l*.tex.
\input{bfw-style.tex}
\definecolor{bfwRot}{HTML}{EF4444}
\definecolor{bfwBgGruen}{HTML}{ECFDF5}
\kopfzeile{Gesamtskript VWL}
\setcounter{tocdepth}{2}

\begin{document}
\thispagestyle{empty}
\vspace*{3cm}
{\bfwlabelfont\color{bfwPrimary}WIRTSCHAFTSFACHWIRT (IHK) \,\textperiodcentered\, GESAMTSKRIPT\par}
\vspace{0.8em}
{\fontsize{34}{40}\selectfont\bfseries\color{bfwPrimaryDark}Volkswirtschaftslehre\par}
\vspace{0.8em}
{\Large\color{bfwInkSoft}Lehrtext, Übungsaufgaben auf IHK-Niveau und Musterlösungen\par}
\vspace{1.2em}
{\color{bfwPrimary}\rule{\linewidth}{2pt}}\par
\vspace{0.8em}
{\small\color{bfwInkSoft}Rahmenplan: 1.1 Volkswirtschaftliche Grundlagen\par
Kurszeitraum: 19.08.2026 -- 26.09.2026 \,\textperiodcentered\, Dozent: Can Siebert\par}
\vfill
\begin{merksatz}[title={\bfseries So arbeiten Sie mit diesem Skript}]
\textbf{Teil A} fasst den gesamten Lehrstoff aller Lektionen zusammen.
\textbf{Teil B} enthält neue Übungsaufgaben im Stil der IHK-Prüfung --- bearbeiten Sie sie
schriftlich und unter Zeitdruck (Richtwert: etwa eine Minute pro Punkt).
Vergleichen Sie erst danach mit den Musterlösungen in \textbf{Teil C}.
\end{merksatz}
\clearpage
\tableofcontents

\clearpage
\addcontentsline{toc}{part}{Teil A\quad Lehrtext}
\vspace*{5cm}
{\bfwlabelfont\color{bfwPrimary}TEIL A\par}
\vspace{0.5em}
{\Huge\bfseries\color{bfwPrimaryDark}Lehrtext\par}
\vspace{0.6em}
{\color{bfwPrimary}\rule{\linewidth}{2pt}}\par
\vspace{0.6em}
{\large\color{bfwInkSoft}Der gesamte Stoff aller Lektionen im Zusammenhang\par}
\clearpage

\clearpage
\section{Lektion 1: Grundbegriffe \& Preisbildung am Markt}
\begin{merksatz}[title={\bfseries Worum geht es in dieser Lektion?}]
Die Volkswirtschaftslehre beginnt bei einem Grundkonflikt: Die menschlichen Bedürfnisse
sind (nahezu) unbegrenzt, die Güter zu ihrer Befriedigung aber knapp. Aus dieser Knappheit
folgt die Notwendigkeit des Wirtschaftens. In dieser Lektion lernen Sie die Grundbegriffe
kennen (Wirtschaftssubjekte, ökonomisches Prinzip, Produktionsfaktoren), unterscheiden
Marktarten und Marktformen und verstehen, wie auf einem Markt aus Angebot und Nachfrage
ein Preis entsteht --- das Marktgleichgewicht bei vollständiger Konkurrenz sowie die
Preisbildung bei unvollständiger Konkurrenz. Das ist die Grundlage für den gesamten
Themenblock 1.1 des Rahmenplans und regelmäßig Gegenstand der IHK-Prüfung.
\end{merksatz}



\subsection{Warum wir wirtschaften: Knappheit und ökonomisches Prinzip}

Der Ausgangspunkt allen Wirtschaftens ist der Konflikt zwischen \textbf{Güterknappheit}
einerseits und der Tendenz \textbf{unbegrenzter Bedürfnisse} andererseits. Weil nicht alle
Wünsche mit den vorhandenen Mitteln erfüllt werden können, müssen Menschen, Unternehmen
und der Staat \emph{wählen} und \emph{haushalten} --- sie wirtschaften.

\begin{itemize}
  \item \textbf{Bedürfnis:} ein empfundener Mangel mit dem Wunsch, ihn zu beseitigen
    (z.\,B. Hunger, Sicherheit, Anerkennung).
  \item \textbf{Bedarf:} der Teil der Bedürfnisse, der mit \emph{Kaufkraft} ausgestattet ist.
  \item \textbf{Nachfrage:} der Bedarf, der am Markt tatsächlich \emph{wirksam} wird.
\end{itemize}

Wirtschaften heißt, nach dem \textbf{ökonomischen Prinzip} (Rationalprinzip) zu handeln.
Es tritt in zwei Ausprägungen auf:

\begin{itemize}
  \item \textbf{Maximalprinzip:} Mit \emph{gegebenen Mitteln} einen \emph{möglichst großen}
    Erfolg erzielen (z.\,B. mit festem Werbebudget maximale Reichweite).
  \item \textbf{Minimalprinzip:} Ein \emph{gegebenes Ziel} mit \emph{möglichst geringem}
    Mitteleinsatz erreichen (z.\,B. vorgegebene Stückzahl mit minimalen Kosten).
\end{itemize}

\begin{eselsbruecke}
\textbf{Ma}ximalprinzip: \textbf{M}ittel fest, \textbf{a}lles rausholen.
\textbf{Mi}nimalprinzip: Ziel fest, \textbf{mi}t wenig hinkommen.
Beides gleichzeitig („mit wenig Mitteln viel erreichen``) ist \emph{kein} ökonomisches
Prinzip --- ein beliebter Prüfungsfehler!
\end{eselsbruecke}

\subsubsection{Wirtschaftssubjekte und Wirtschaftssektoren}

Ein \textbf{Wirtschaftssubjekt} ist das kleinste wirtschaftende Grundelement, das
selbstständig wirtschaftliche Entscheidungen treffen kann --- z.\,B. ein Privathaushalt,
ein Unternehmen oder eine staatliche Körperschaft. Für gesamtwirtschaftliche Betrachtungen
werden die Wirtschaftssubjekte zu \emph{Wirtschaftssektoren} zusammengefasst: private
Haushalte, Unternehmen, Staat, (Banken) und das Ausland.

\subsubsection{Produktionsfaktoren}

Der Produktionsprozess (Leistungserstellung) erfolgt volkswirtschaftlich durch den Einsatz
der \textbf{Produktionsfaktoren}:

\begin{itemize}
  \item \textbf{Arbeit} --- jede auf Erwerb gerichtete körperliche und geistige Tätigkeit,
  \item \textbf{Boden} --- Anbau, Abbau (Rohstoffe) und Standort,
  \item \textbf{Kapital} --- produzierte Produktionsmittel (Maschinen, Anlagen, Gebäude),
  \item zunehmend auch \textbf{Wissen (Know-how)} als eigenständiger vierter Faktor.
\end{itemize}

\begin{merksatz}
Unterscheiden Sie die \emph{volkswirtschaftlichen} Produktionsfaktoren (Arbeit, Boden,
Kapital, Wissen) von den \emph{betriebswirtschaftlichen} (Elementarfaktoren und dispositiver
Faktor) --- Letztere behandeln wir im BWL-Teil des Kurses.
\end{merksatz}

\subsection{Der Markt: Arten und Formen}

\subsubsection{Was ist ein Markt?}

Der \textbf{Markt} ist der (reale oder virtuelle) Ort, an dem Angebot und Nachfrage
zusammentreffen und sich der \textbf{Preis} bildet. Seine Aufgaben: Verteilung der Güter,
bestmögliche Versorgung der Marktteilnehmer, Preisbildung. Ziel der Nachfrager ist das
\emph{Nutzenmaximum}, Ziel der Anbieter die \emph{Gewinnmaximierung} (mindestens
Kostendeckung). Für jede Gesellschaft stellen sich drei Grundfragen: \emph{Was} soll
\emph{wie} und \emph{für wen} produziert werden?

\subsubsection{Marktarten (nach dem Marktgegenstand)}

\begin{center}
\begin{tabular}{p{3.6cm} p{4.2cm} p{6.2cm}}
\textbf{Marktart} & \textbf{Marktgegenstand} & \textbf{Beispiele}\\
\hline
Faktormarkt & Produktionsfaktoren & Arbeitsmarkt, Immobilienmarkt, Sachkapitalmarkt\\
Gütermarkt & Güter & Konsumgüter, Investitionsgüter, Dienstleistungen\\
Geld- und Kapitalmarkt (Finanzmarkt) & Geld, Kredite, Kapitalanlagen & Devisenmarkt, Aktienmarkt\\
\end{tabular}
\end{center}

\subsubsection{Vollkommener und unvollkommener Markt}

Der \textbf{vollkommene Markt} ist ein \emph{Idealtyp}, der in der Realität nicht
anzutreffen ist (die Börse kommt ihm sehr nah). Bedingungen:

\begin{enumerate}
  \item \textbf{Homogenität der Güter} --- sachlich gleichwertige Güter, keine Unterschiede
    in Qualität, Aufmachung, Verpackung;
  \item \textbf{Fehlen von Präferenzen} --- keine persönlichen, räumlichen oder zeitlichen
    Vorlieben;
  \item \textbf{vollkommene Markttransparenz} --- alle Teilnehmer haben vollständigen
    Überblick über Preise und Konditionen;
  \item \textbf{unendlich schnelle Reaktionen} der Marktteilnehmer (Punktmarkt);
  \item \textbf{polypolistische Marktstruktur} --- viele Anbieter, viele Nachfrager, freier
    Marktzutritt.
\end{enumerate}

Ist \emph{mindestens eine} Bedingung nicht erfüllt, liegt ein \textbf{unvollkommener
Markt} vor.

\subsubsection{Marktformen (nach der Zahl der Marktteilnehmer)}

\begin{center}
\begin{tabular}{p{2.6cm}|p{3.6cm}|p{3.6cm}|p{3.9cm}}
\textbf{Nachfrager $\downarrow$ / Anbieter $\rightarrow$} & \textbf{viele} & \textbf{wenige} & \textbf{einer}\\
\hline
\textbf{viele} & (zweiseitiges) Polypol \newline \small z.\,B. Wochenmärkte & Angebotsoligopol \newline \small z.\,B. Mobilfunkanbieter & Angebotsmonopol \newline \small z.\,B. regionaler Wasserversorger\\
\hline
\textbf{wenige} & Nachfrageoligopol \newline \small z.\,B. Großmolkereien ggü. Milchbauern & (zweiseitiges) Oligopol \newline \small Spezialmaschinen & beschränktes Angebotsmonopol\\
\hline
\textbf{einer} & Nachfragemonopol \newline \small Staat bei Straßenbau & beschränktes Nachfragemonopol \newline \small Bundeswehr ggü. wenigen Anbietern & (zweiseitiges) Monopol\\
\end{tabular}
\end{center}

\begin{beispiel}
Benzinmarkt: \emph{wenige} Mineralölkonzerne stehen \emph{vielen} Autofahrern gegenüber
$\Rightarrow$ \textbf{Angebotsoligopol}. Apothekenmarkt: \emph{viele} Apotheken,
\emph{viele} Kunden $\Rightarrow$ \textbf{Polypol} (Monopole sind bei Apotheken
\emph{nicht} typisch --- beliebte Quiz-Falle).
\end{beispiel}

\subsection{Die Nachfrage}

\subsubsection{Bestimmungsfaktoren der Nachfrage}

Die Nachfrager (Haushalte, aber auch Unternehmen, Staat, Ausland) wollen mit begrenzten
Mitteln ihren Nutzen maximieren. Wichtige Bestimmungsfaktoren der Güternachfrage:

\begin{itemize}
  \item \textbf{Preis des Gutes:} Bei hohem Preis wird i.\,d.\,R. weniger nachgefragt.
  \item \textbf{Preise anderer Güter:} \emph{komplementäre} Güter (ergänzen sich, z.\,B.
    Drucker und Druckerpatronen) und \emph{substitutive} Güter (ersetzen sich, z.\,B.
    Butter und Margarine).
  \item \textbf{Verfügbares Einkommen} (Konsumsumme).
  \item \textbf{Bedürfnisstruktur} (Mode, Demografie, Technologie, Werte).
  \item \textbf{Zukunftserwartungen} bezüglich Preisen und Einkommen.
\end{itemize}

\subsubsection{Elastizitäten: Wie stark reagieren die Kunden?}
\label{sec:elastizitaet}

Die \textbf{Preiselastizität der Nachfrage} beantwortet eine ganz praktische Frage:
\emph{Wie stark reagieren die Kunden, wenn der Preis sich ändert?} Erhöht ein Bäcker
den Brötchenpreis um 10\,\%, kaufen die Kunden dann fast unverändert weiter (träge
Reaktion) --- oder laufen sie in Scharen zur Konkurrenz (heftige Reaktion)? Genau
dieses Reaktionsausmaß fasst die Elastizität in einer einzigen Zahl zusammen.

\medskip
\noindent
\begin{minipage}[c]{0.55\textwidth}
\begin{formelbox}
\[
\varepsilon \;=\; \frac{\text{Mengenänderung in \%}}{\text{Preisänderung in \%}}
\]
\end{formelbox}
\end{minipage}\hfill
\begin{minipage}[c]{0.41\textwidth}
\small
\textbf{In Worten:} Oben steht die \emph{Reaktion} der Kunden (um wie viel Prozent
ändert sich die nachgefragte \textbf{Menge}?), unten der \emph{Auslöser} (um wie viel
Prozent hat sich der \textbf{Preis} geändert?). Beides immer in Prozent --- nie in
Euro oder Stück!
\end{minipage}
\medskip

Da Preis und Menge normalerweise gegenläufig reagieren (Preis rauf, Menge runter),
ist $\varepsilon$ meist \textbf{negativ}. Für die Interpretation zählt allein der
\textbf{Betrag} $|\varepsilon|$, also die Zahl ohne Vorzeichen:

\begin{center}
\small
\begin{tabular}{p{2.7cm} p{4.0cm} p{3.3cm} p{4.0cm}}
\toprule
\rowcolor{bfwTabHead}
\textbf{Ergebnis} & \textbf{Bedeutung in einfachen Worten} & \textbf{Typische Güter} & \textbf{Preiserhöhung wirkt auf den Umsatz \ldots}\\
\midrule
$|\varepsilon| > 1$ \newline \textbf{elastisch} &
  Kunden reagieren \emph{stark}: Die Menge ändert sich prozentual stärker als der Preis. &
  Luxusgüter, Urlaubsreisen, Markenartikel --- alles, worauf man verzichten oder ausweichen kann &
  \ldots{} \textbf{senkend}: Der Mengenrückgang frisst die Preiserhöhung mehr als auf.\\
\addlinespace
$|\varepsilon| < 1$ \newline \textbf{unelastisch} &
  Kunden reagieren \emph{schwach}: Die Menge ändert sich prozentual weniger als der Preis. &
  Lebensnotwendiges: Grundnahrungsmittel, Medikamente, Strom, Benzin &
  \ldots{} \textbf{steigernd}: Fast alle kaufen weiter --- zum höheren Preis.\\
\addlinespace
$|\varepsilon| = 0$ \newline \textbf{vollkommen unelastisch} &
  Kunden reagieren \emph{gar nicht}: Die Menge bleibt exakt gleich. &
  lebenswichtige Güter ohne jeden Ersatz, z.\,B. Insulin für Diabetiker &
  \ldots{} \textbf{voll steigernd}: Der Umsatz wächst im selben Maß wie der Preis.\\
\bottomrule
\end{tabular}
\end{center}

\begin{beispiel}[title={SCHRITT FÜR SCHRITT: PREISELASTIZITÄT BERECHNEN}]
Ein Medikament kostet 2,00~€; der Preis steigt auf 2,40~€. Die verkaufte Menge sinkt
von 1.000 auf 980 Packungen.
\begin{enumerate}
  \item \textbf{Preisänderung in \% berechnen:} Änderung $= 2{,}40 - 2{,}00 = 0{,}40$~€.
    Bezogen auf den \emph{Ausgangspreis}: $0{,}40 \div 2{,}00 = \mathbf{+20\,\%}$.
  \item \textbf{Mengenänderung in \% berechnen:} Änderung $= 980 - 1.000 = -20$ Packungen.
    Bezogen auf die \emph{Ausgangsmenge}: $-20 \div 1.000 = \mathbf{-2\,\%}$.
  \item \textbf{In die Formel einsetzen:}
    $\varepsilon = -2\,\% \div (+20\,\%) = -0{,}1$, also $|\varepsilon| = 0{,}1$.
  \item \textbf{Interpretieren:} $|\varepsilon| = 0{,}1 < 1$ $\Rightarrow$
    \textbf{unelastisch} --- die Käufer können auf das lebensnotwendige Gut kaum
    verzichten.
  \item \textbf{Umsatz-Probe:} vorher $2{,}00~\text{€} \times 1.000 = 2.000$~€,
    nachher $2{,}40~\text{€} \times 980 = 2.352$~€. Der Umsatz \emph{steigt} --- genau
    wie es die Tabelle für unelastische Güter vorhersagt.
\end{enumerate}
Zum Vergleich: Bei einem Modeartikel mit Preis $+10\,\%$ und Menge $-25\,\%$ wäre
$|\varepsilon| = 2{,}5$ --- \emph{elastisch}, der Umsatz würde sinken.
\end{beispiel}

\begin{achtung}[title={TYPISCHE FEHLER BEI DER PREISELASTIZITÄT}]
\begin{itemize}
  \item \textbf{Vorzeichen falsch gedeutet:} Ein negatives $\varepsilon$ ist bei der
    Preiselastizität \emph{normal} (Preis rauf, Menge runter). Für
    elastisch/unelastisch zählt nur der Betrag $|\varepsilon|$ --- nicht das
    Minuszeichen.
  \item \textbf{Absolute statt prozentuale Änderungen:} „$0{,}40$~€ geteilt durch
    20~Stück`` ist falsch --- beide Änderungen müssen \emph{in Prozent}, jeweils
    bezogen auf den \emph{Ausgangswert}, in die Formel.
  \item \textbf{Zähler und Nenner vertauscht:} Die \emph{Menge} gehört nach oben,
    der \emph{Preis} nach unten. Merkhilfe: erst die Reaktion, dann der Auslöser.
  \item \textbf{Verwechslung mit der Kreuzpreiselastizität:} Die direkte
    Preiselastizität betrachtet \emph{ein} Gut (eigener Preis $\rightarrow$ eigene
    Menge). Sobald \emph{zwei} Güter im Spiel sind (Butter/Margarine), ist die
    Kreuzpreiselastizität gefragt.
\end{itemize}
\end{achtung}

\subsubsection{Kreuzpreiselastizität und Einkommenselastizität --- kurz erklärt}

Die \textbf{Kreuzpreiselastizität} fragt: Wie reagiert die Menge von Gut~A, wenn sich
der \emph{Preis eines anderen Gutes~B} ändert?
\[
\varepsilon_K \;=\; \frac{\text{Mengenänderung bei Gut A in \%}}{\text{Preisänderung bei Gut B in \%}}
\]
Hier trägt das \textbf{Vorzeichen} die Information: \textbf{positiv} $\Rightarrow$
\emph{Substitute} (Butter wird teurer $\rightarrow$ mehr Margarine); \textbf{negativ}
$\Rightarrow$ \emph{Komplemente} (Drucker wird teurer $\rightarrow$ weniger Patronen).

\begin{beispiel}[title={BEISPIEL: KREUZPREISELASTIZITÄT}]
Der Butterpreis (Gut~B) steigt um $10\,\%$; die nachgefragte Menge an Margarine
(Gut~A) steigt um $5\,\%$. Einsetzen:
$\varepsilon_K = +5\,\% \div (+10\,\%) = +0{,}5$. Das \emph{positive} Vorzeichen
zeigt: Butter und Margarine sind \textbf{Substitute}.
\end{beispiel}

Die \textbf{Einkommenselastizität} fragt: Wie reagiert die nachgefragte Menge, wenn
sich das \emph{Einkommen} der Haushalte ändert? Bei normalen Gütern ist sie positiv
(mehr Einkommen $\rightarrow$ mehr Nachfrage), bei Luxusgütern besonders hoch, bei
inferioren Gütern (z.\,B. einfachste Lebensmittel) negativ --- wer mehr verdient,
kauft davon weniger.

\subsubsection{Nachfragekurve: Bewegung und Verschiebung}

Die Nachfragekurve verläuft \textbf{fallend}: Steigt der Preis, sinkt die nachgefragte
Menge (Bewegung \emph{auf} der Kurve). Ändern sich \emph{andere} Einflussgrößen als der
Preis, \textbf{verschiebt} sich die gesamte Kurve:

\begin{itemize}
  \item \textbf{Linksverschiebung} (Nachfragesenkung): z.\,B. Gut kommt aus der Mode,
    Realeinkommen sinkt, Preis des Substituts fällt, Preis des Komplements steigt.
  \item \textbf{Rechtsverschiebung} (Nachfrageerhöhung): z.\,B. höhere Realeinkommen,
    optimistische Zukunftserwartungen.
\end{itemize}

\subsection{Das Angebot}

Bestimmungsfaktoren des Güterangebots:

\begin{itemize}
  \item \textbf{Preis des Gutes:} steigender Marktpreis $\Rightarrow$ Anreiz, die
    Produktionsmenge auszuweiten;
  \item \textbf{Kosten} der Produktionsfaktoren (Material-, Lohnkosten);
  \item \textbf{Stand der Technik} (Produktivität).
\end{itemize}

Die Angebotskurve verläuft \textbf{steigend}. Steigende Faktorkosten oder sinkende
Produktivität verschieben sie nach \emph{links} (Angebotssenkung), Kostensenkungen oder
technischer Fortschritt nach \emph{rechts} (Angebotserhöhung).

\subsection{Preisbildung im Marktgleichgewicht (vollständige Konkurrenz)}

Liegt eine Kombination aus \emph{vollkommenem Markt} und \emph{Polypol} vor, spricht man
vom Modell der \textbf{vollständigen Konkurrenz}. Kein einzelner Teilnehmer kann den Preis
beeinflussen --- alle sind \textbf{Preisnehmer} und \textbf{Mengenanpasser}.

\subsubsection{Der Preismechanismus}

\begin{itemize}
  \item \textbf{Angebotsüberschuss} (Nachfragedefizit): Bei hohem Preis übersteigt das
    Angebot die Nachfrage $\Rightarrow$ \emph{Käufermarkt}, Anbieter senken die Preise.
  \item \textbf{Nachfrageüberschuss} (Angebotsdefizit): Bei niedrigem Preis übersteigt die
    Nachfrage das Angebot $\Rightarrow$ \emph{Verkäufermarkt}, Preise steigen.
  \item Der Markt tendiert zum \textbf{Gleichgewicht}: Die angebotene Menge entspricht der
    nachgefragten Menge. Die \emph{Gleichgewichtsmenge} ($x^*$) wird zum
    \emph{Gleichgewichtspreis} ($p^*$) abgesetzt --- grafisch der \textbf{Schnittpunkt} von
    Angebots- und Nachfragekurve. Nur dort findet der mengenmäßig größte Umsatz statt.
\end{itemize}

Dieser Anpassungsprozess über Preis und Menge heißt \textbf{Preismechanismus} oder
\textbf{Marktmechanismus}.

\subsubsection{Konsumenten- und Produzentenrente}

\begin{itemize}
  \item \textbf{Konsumentenrente:} Nutzengewinn der Nachfrager, die bereit gewesen wären,
    \emph{mehr} als den Gleichgewichtspreis zu zahlen.
  \item \textbf{Produzentenrente:} Mehrerlös der Anbieter, die bereit gewesen wären, ihre
    Güter \emph{billiger} als zum Gleichgewichtspreis zu verkaufen.
\end{itemize}

\subsubsection{Funktionen des Marktpreises}

\begin{itemize}
  \item \textbf{Signalfunktion} (Informationsfunktion): Der Preis spiegelt die Knappheit
    eines Gutes wider.
  \item \textbf{Ausgleichsfunktion}: Preise gleichen Angebots- und Nachfragemenge
    tendenziell aus.
  \item \textbf{Allokationsfunktion} (Lenkungsfunktion): Preise lenken die
    Produktionsfaktoren in die Produktion marktfähiger Güter.
  \item \textbf{Selektionsfunktion} (Sanktionsfunktion): Anbieter mit zu hohen Kosten und
    Nachfrager, die den Preis nicht zahlen können oder wollen, scheiden aus dem Markt aus.
\end{itemize}

\begin{eselsbruecke}
\textbf{SAAS} --- \textbf{S}ignal, \textbf{A}usgleich, \textbf{A}llokation,
\textbf{S}elektion: die vier Funktionen des Marktpreises.
\end{eselsbruecke}

\subsection{Preisbildung bei unvollständiger Konkurrenz}

In der Realität existieren auf unvollkommenen Märkten unterschiedliche Preise, weil Güter
heterogen sind, Präferenzen bestehen und die Markttransparenz fehlt.

\subsubsection{Unvollkommenes Polypol: der monopolistische Absatzbereich}

Viele Anbieter stehen vielen Nachfragern auf einem \emph{unvollkommenen} Markt gegenüber.
Durch \textbf{Produktdifferenzierung} (Marken, Werbung, Qualität) schafft sich der Anbieter
einen \textbf{monopolistischen Absatzbereich}: Innerhalb dieses Preisspielraums kann er den
Preis nahezu frei festlegen, ohne größere Nachfrageeinbußen. Überschreitet er ihn, wandern
Kunden ab; unterschreitet er ihn, gewinnt er Kunden, deckt aber evtl. seine Kosten nicht
mehr.

\subsubsection{Oligopol}

Wenige Anbieter mit großen Marktanteilen (z.\,B. Waschmittel, Automobile, Mineralöl)
beobachten sich gegenseitig. Typische Verhaltensweisen:

\begin{itemize}
  \item \textbf{Marktverdrängungspolitik:} aggressive Preissenkungen, um Marktanteile zu
    gewinnen (ruinöser Preiskampf --- selten);
  \item \textbf{Parallelverhalten:} abgestimmtes Verhalten ohne Absprache, häufig
    \emph{Preisruhe} (Preisstarre) und \emph{Gleichschritt bei Preisänderungen};
  \item \textbf{Preisführerschaft:} ein Anbieter erhöht die Preise, die übrigen ziehen nach
    (z.\,B. Tankstellen vor Feiertagen). Der Wettbewerb verlagert sich auf Qualität, Marke,
    Werbung und Service.
  \item Gefahr verbotener \textbf{Absprachen} (Kartellbildung) --- dazu mehr in Lektion~2
    (Wettbewerbspolitik).
\end{itemize}

\subsubsection{Monopol}

Beim \textbf{Angebotsmonopol} beliefert nur ein Anbieter den Markt --- er ist theoretisch
\textbf{Preisfixierer}. Die Marktmacht ist besonders stark, wenn kaum Substitutionsgüter
existieren und das Produkt dringend benötigt wird (z.\,B. regionaler Wasserversorger).
Langfristig bestehen Monopole nur bei Marktzutrittsbeschränkungen.

\subsection{Formeln \& Rechenbeispiele}

Alle Berechnungen dieser Lektion folgen einfachen Grundformeln. Rechnen Sie die
Beispiele einmal selbst nach --- genau solche kleinen Rechnungen kommen in der
IHK-Prüfung vor.

\subsubsection{Preiselastizität, Kreuzpreiselastizität \& Einkommenselastizität}

Formeln, Interpretations-Tabelle und die durchgerechneten Schritt-für-Schritt-Beispiele
(Medikament: Preis $2{,}00 \rightarrow 2{,}40$~€, Menge $1.000 \rightarrow 980$
Packungen; Butter/Margarine) finden Sie gebündelt in
Abschnitt~\ref{sec:elastizitaet} --- inklusive der Box „Typische Fehler``.
Prüfen Sie sich selbst: Können Sie das Medikamenten-Beispiel ohne Blick in den
Lösungsweg nachrechnen und das Ergebnis interpretieren?

\subsubsection{Gleichgewichtspreis und -menge aus einer Tabelle bestimmen}

\begin{merksatz}[title={\bfseries Vorgehen}]
Gesucht ist der Preis, bei dem \textbf{angebotene Menge = nachgefragte Menge}. Diese
Zeile liefert Gleichgewichtspreis $p^*$ und Gleichgewichtsmenge $x^*$; der Umsatz im
Gleichgewicht beträgt $p^* \cdot x^*$.
\end{merksatz}

\begin{beispiel}
Für einen Markt liegen folgende Zahlen vor:

\begin{center}
\begin{tabular}{cccl}
\textbf{Preis} & \textbf{Nachfrage (Stück)} & \textbf{Angebot (Stück)} & \textbf{Marktlage}\\
\hline
2~€ & 180 & 60  & Nachfrageüberschuss $\Rightarrow$ Verkäufermarkt\\
3~€ & 140 & 100 & Nachfrageüberschuss $\Rightarrow$ Verkäufermarkt\\
\textbf{4~€} & \textbf{120} & \textbf{120} & \textbf{Gleichgewicht}\\
5~€ & 90  & 150 & Angebotsüberschuss $\Rightarrow$ Käufermarkt\\
\end{tabular}
\end{center}

\begin{enumerate}
  \item \textbf{Zeile suchen}, in der Angebot und Nachfrage übereinstimmen: bei
    \textbf{4~€} sind es jeweils 120 Stück.
  \item \textbf{Ablesen:} Gleichgewichtspreis $p^* = 4$~€, Gleichgewichtsmenge
    $x^* = 120$ Stück.
  \item \textbf{Umsatz im Gleichgewicht:} $4~\text{€} \times 120 = 480$~€.
\end{enumerate}
\textbf{Interpretation:} Oberhalb von $p^*$ (5~€) übersteigt das Angebot die Nachfrage
--- Käufermarkt, Preisdruck nach unten. Unterhalb von $p^*$ (2~€ oder 3~€) übersteigt
die Nachfrage das Angebot --- Verkäufermarkt, die Preise steigen. Der Preismechanismus
führt den Markt zum Gleichgewicht zurück.
\end{beispiel}

\begin{merksatz}[title={\bfseries Wichtig für die Prüfungsvorbereitung}]
Sie sollten sicher können: die Definition „Markt``, die \emph{fünf Bedingungen des
vollkommenen Marktes}, die \emph{Marktformen-Matrix} mit Beispielen, die Skizze der
\emph{Preisbildung} (Gleichgewicht, Überschüsse), die \emph{vier Preisfunktionen},
den Unterschied zwischen \emph{Bewegung auf} und \emph{Verschiebung von} Kurven sowie
die \emph{Rechenwege} für Preiselastizität, Kreuzpreiselastizität und Gleichgewichtspreis
(Abschnitte „Elastizitäten`` und „Formeln \& Rechenbeispiele``).
Üben Sie mit dem Aufgabenheft und dem Quiz dieser Lektion!
\end{merksatz}

\vspace{1em}
\hrule
\vspace{0.8em}
\noindent\textsf{\small Quellen: Rahmenplan Wirtschaftsbezogene Qualifikationen (DIHK), Abschnitt 1.1.1;
Textband Volkswirtschaft (DIHK-Bildungs-GmbH), Kap.~1, S.~1--11. Nächste Lektion (Sa, 12.09.2026):
Wettbewerbspolitik, Eingriffe des Staates in die Preisbildung \& Volkswirtschaftliche Gesamtrechnung.}
\clearpage
\section{Lektion 2: Wettbewerbspolitik, Staatseingriffe \& VGR}
\par\medskip{\bfwlabelfont\color{bfwPrimary}SESSION 1\par}\smallskip
\begin{merksatz}[title={\bfseries Worum geht es in Teil 1?}]
Basis der sozialen Marktwirtschaft ist ein funktionierender Markt mit ausreichendem
Wettbewerb. Da ein sich selbst überlassener Markt die Tendenz zeigt, Wettbewerb zu
reduzieren, greift der Staat regulierend ein: durch \emph{Wettbewerbspolitik} (GWB,
Bundeskartellamt) und durch \emph{Eingriffe in die Preisbildung} (Höchst- und
Mindestpreise, Subventionen, Steuern). Beide Themen sind Standardstoff der IHK-Prüfung.
In Teil~2 (Session~2) folgt darauf die Erfolgsmessung: Volkswirtschaftliche
Gesamtrechnung und die Ziele der Stabilitätspolitik.
\end{merksatz}



\subsection{Wettbewerb und Wettbewerbspolitik}

\subsubsection{Was ist Wettbewerb?}

Unter \textbf{Wettbewerb} versteht man allgemein die \emph{Rivalität der Marktteilnehmer
unter- und gegeneinander}, um ihre Ziele zulasten der Konkurrenten durchzusetzen. Der
Wettbewerb sichert u.\,a. die persönliche und wirtschaftliche Freiheit, die
Güterverteilung und Bedarfsdeckung über den Marktmechanismus und den ökonomischen
Einsatz der Produktionsfaktoren.

\subsubsection{Funktionen des Wettbewerbs}

Idealtypisch erfüllt der Wettbewerb folgende Funktionen:

\begin{itemize}
  \item \textbf{Freiheitsfunktion:} Jeder kann seine Fähigkeiten und Mittel für frei
    gewählte Zwecke einsetzen (Freiheit der Unternehmertätigkeit, Freiheit der
    Konsumwahl) --- freier Wettbewerb schafft mehr Handlungsalternativen.
  \item \textbf{Kontrollfunktion:} Anbieter und Nachfrager kontrollieren sich gegenseitig;
    jeder muss sich an der Leistung der Konkurrenz messen lassen. Die wirtschaftliche
    Macht des Einzelnen wird begrenzt.
  \item \textbf{Steuerungsfunktion (Lenkungsfunktion):} Mit zwei Unterfunktionen ---
    \emph{Koordinations- und Anpassungsfunktion} (Wettbewerb beschleunigt die Koordination
    von Angebot und Nachfrage; langfristig wird nur angeboten, was der Nachfrage
    entspricht) und \emph{Allokationsfunktion} (die Produktionsfaktoren Arbeit, Boden,
    Kapital, Know-how werden so kombiniert, dass ein Maximum an Produktivität entsteht).
  \item \textbf{Innovationsfunktion:} Wettbewerb schafft Anreize, Produkte zu verbessern,
    die Produktpalette zu erneuern, kostengünstiger zu produzieren und den technischen
    Fortschritt rasch zu nutzen.
  \item \textbf{Auslesefunktion:} Auslese der Marktteilnehmer nach wirtschaftlicher
    Leistungsfähigkeit --- stets kommen die leistungsfähigsten Anbieter und Nachfrager
    zum Zug.
  \item \textbf{Verteilungsfunktion:} Die Entlohnung der Produktionsfaktoren entspricht
    deren Leistung --- leistungsorientierte Einkommensverteilung.
\end{itemize}

\begin{eselsbruecke}
\textbf{F-K-S-I-A-V} --- „\textbf{F}reier \textbf{K}onkurrenzkampf \textbf{s}ichert
\textbf{i}mmer \textbf{a}lle \textbf{V}orteile``: \textbf{F}reiheit, \textbf{K}ontrolle,
\textbf{S}teuerung, \textbf{I}nnovation, \textbf{A}uslese, \textbf{V}erteilung ---
die sechs Funktionen des Wettbewerbs.
\end{eselsbruecke}

\subsubsection{Ziele und Instrumente der Wettbewerbspolitik}

Weil unvollkommene Märkte zu Wettbewerbsverzerrungen führen und Märkte die Tendenz haben,
Wettbewerb abzubauen, steuert der Staat im Rahmen der \textbf{Ordnungspolitik} gegen.
Schwerpunkte: freier Zugang zu Märkten, Abbau von Handelshemmnissen, Verhinderung
wettbewerbsbeschränkender Verhaltensweisen. Leitbild ist heute nicht mehr die
vollständige Konkurrenz, sondern der Erhalt von \emph{Qualität und Intensität} des
Wettbewerbs --- Wissen, Kreativität und Innovationen gelten als treibende Marktkräfte.

Die beiden gesetzlichen Grundlagen in Deutschland:

\begin{itemize}
  \item \textbf{UWG} --- Gesetz gegen den unlauteren Wettbewerb: garantiert fairen
    Wettbewerb und faire Marktpraktiken (Schutz der \emph{Art und Weise} des Wettbewerbs).
    Unzulässig sind z.\,B. irreführende und aggressive Werbung, Schleichwerbung oder die
    gezielte Herabsetzung von Mitbewerbern.
  \item \textbf{GWB} --- Gesetz gegen Wettbewerbsbeschränkungen („Kartellgesetz``):
    unterbindet Kartelle, Absprachen und Preisbindungen und wirkt der Konzentration der
    Wirtschaft entgegen (Schutz des \emph{Bestands} von Wettbewerb).
\end{itemize}

Über die Einhaltung des GWB wacht das \textbf{Bundeskartellamt}. In der EU liegt die
Wettbewerbspolitik bei der \textbf{EU-Kommission} (Kommissar für Wettbewerb);
Rechtsgrundlagen sind der \textbf{AEUV} (Vertrag über die Arbeitsweise der EU) und die
europäische \textbf{Fusionskontrollverordnung (FKVO)}. Durch die GWB-Novellen (7.~Novelle
2005: Angleichung an EU-Recht; 8.~Novelle 2013: Neustrukturierung der Missbrauchsaufsicht,
Angleichung der Fusionskontrolle; 9.~Novelle 2017: Umsetzung der
EU-Kartellschadenersatzrichtlinie) ähneln sich deutsches und europäisches Recht sehr.
Jüngere Novellen: Die 10.~Novelle (GWB-Digitalisierungsgesetz 2021) führte u.\,a.
\S\,19a GWB ein --- erweiterte Aufsicht über marktübergreifend bedeutende
Digitalkonzerne ---, die 11.~Novelle (2023) erweiterte die Befugnisse für
Sektoruntersuchungen.
\emph{EU-Recht steht über nationalem Recht} --- es greift, wenn Wettbewerbsbeschränkungen
den Handel zwischen Mitgliedstaaten beeinträchtigen. Ein weltweites Wettbewerbsrecht oder
Kartellamt existiert nicht.

\begin{merksatz}[title={\bfseries Die drei Säulen des GWB}]
\textbf{1. Kartellverbot} (\S\,1 GWB) \quad \textbf{2. Missbrauchsaufsicht}
(\S\S\,19\,ff. GWB) \quad \textbf{3. Fusionskontrolle} (\S\,35 GWB).
Merkhilfe: \textbf{Ka-Mi-Fu}.
\end{merksatz}

\begin{center}
\begin{tikzpicture}[font=\small,
    saeule/.style={rounded corners=4pt, thick, align=center, text width=3.6cm,
                   minimum height=2.9cm, inner sep=6pt}]
  \node[roughbox, minimum width=13.4cm, minimum height=1.05cm, align=center] (dach) at (0,0)
    {\bfseries GWB --- Gesetz gegen Wettbewerbsbeschränkungen („Kartellgesetz``)};
  \node[saeule, draw=bfwPrimaryDark, fill=bfwBgTeal] (s1) at (-4.7,-2.9)
    {{\color{bfwPrimaryDark}\bfseries 1. Kartellverbot}\\[1pt]
     {\footnotesize \S\,1 GWB}\\[4pt]
     {\footnotesize Preis-, Submissions-, Gebiets- und Quotenkartelle verboten;
      Ausnahmen: \S\S\,2--3 GWB, EU-Freistellungen}};
  \node[saeule, draw=bfwAccent!85!black, fill=bfwBgAmber] (s2) at (0,-2.9)
    {{\color{bfwAccent!85!black}\bfseries 2. Missbrauchsaufsicht}\\[1pt]
     {\footnotesize \S\S\,19\,ff. GWB}\\[4pt]
     {\footnotesize Ausbeutungs- und Behinderungsmissbrauch (z.\,B. Dumping)
      marktbeherrschender Unternehmen}};
  \node[saeule, draw=bfwWarn!85!black, fill=bfwBgRose] (s3) at (4.7,-2.9)
    {{\color{bfwWarn!85!black}\bfseries 3. Fusionskontrolle}\\[1pt]
     {\footnotesize \S\,35 GWB}\\[4pt]
     {\footnotesize Anmeldepflicht, Untersagung, Ministererlaubnis (\S\,42 GWB)
      als Ausnahme}};
  \node[rounded corners=4pt, thick, draw=bfwInkSoft, fill=bfwBgSoft, align=center,
        minimum width=13.4cm, minimum height=0.95cm] (amt) at (0,-5.6)
    {\bfseries Bundeskartellamt {\normalfont\footnotesize --- wacht über das GWB
     \textperiodcentered{} EU-Ebene: EU-Kommission (Art.~101 AEUV, FKVO)}};
  \draw[roughline] (s1.north) -- (s1.north |- dach.south);
  \draw[roughline] (s2.north) -- (s2.north |- dach.south);
  \draw[roughline] (s3.north) -- (s3.north |- dach.south);
  \draw[roughline] (amt.north -| s1.south) -- (s1.south);
  \draw[roughline] (amt.north -| s2.south) -- (s2.south);
  \draw[roughline] (amt.north -| s3.south) -- (s3.south);
\end{tikzpicture}
\end{center}

\subsubsection{Das Kartellverbot (\S\,1 GWB)}

\textbf{Kartelle} sind Kooperationen von Unternehmen auf Basis (schriftlicher oder
mündlicher) vertraglicher Vereinbarungen, um die Wettbewerbsposition zu verbessern. Die
Beteiligten bleiben \emph{rechtlich selbstständig}, schränken ihre \emph{wirtschaftliche}
Selbstständigkeit aber teilweise ein. Nach \S\,1 GWB sind Vereinbarungen zwischen
Unternehmen, Beschlüsse von Unternehmensvereinigungen und aufeinander abgestimmte
Verhaltensweisen \textbf{verboten}, wenn sie eine \emph{Verhinderung, Einschränkung oder
Verfälschung des Wettbewerbs} bewirken. Das europarechtliche Pendant ist
\textbf{Art.~101 AEUV}. Verbotene Kartellarten insbesondere:

\begin{center}
\begin{tabular}{p{4.2cm} p{9.8cm}}
\toprule
\textbf{Kartellart} & \textbf{Inhalt der Absprache}\\
\midrule
Preiskartell & Absprache einer einheitlichen Preispolitik\\
Submissionskartell & Vereinbarung von Preisuntergrenzen bzw. Festlegung, welches Kartellmitglied bei öffentlichen Ausschreibungen zum Zug kommt\\
Gebietskartell & Aufteilung des Absatzmarktes\\
Quotenkartell & künstliche Verknappung des Angebots durch Absprache von Produktionsquoten\\
\bottomrule
\end{tabular}
\end{center}

\begin{achtung}
Abgrenzung (beliebte Prüfungsfrage): Beim \textbf{Kartell} bleiben die Unternehmen
\emph{rechtlich} selbstständig und geben nur einen Teil ihrer \emph{wirtschaftlichen}
Selbstständigkeit auf. Beim \textbf{Konzern} bleiben sie rechtlich selbstständig, stehen
aber unter einheitlicher Leitung. Bei der \textbf{Fusion} verschmelzen sie auch
rechtlich zu einer Einheit --- dann greift die Fusionskontrolle (\S\,35 GWB), nicht das
Kartellverbot.
\end{achtung}

\subsubsection{Ausnahmen vom Kartellverbot}

Obwohl Kartelle grundsätzlich verboten sind, gibt es \textbf{Ausnahmen} in drei Gruppen:

\begin{enumerate}
  \item \textbf{Besondere Ausnahmeregelungen nach dem GWB:} z.\,B.
    \emph{Mittelstandskartelle} (\S\,3 GWB) --- Vereinbarungen zwischen kleinen und
    mittleren Unternehmen, die wegen ihres begrenzten Potenzials Wettbewerbsnachteile
    gegenüber Großunternehmen haben.
  \item \textbf{EU-Gruppenfreistellungsverordnungen:} Bestimmte Gruppen von Vereinbarungen
    sind freigestellt, weil sie sich trotz gewisser Wettbewerbsbeschränkung positiv auf
    den Markt auswirken; Einzelfallprüfung entfällt.
  \item \textbf{Einzelfreistellungen (Legalausnahmen, \S\,2 GWB / Art.~101 Abs.~3 AEUV):}
    Vereinbarungen sind freigestellt, wenn sie (a) den Verbraucher angemessen am Gewinn
    beteiligen, (b) der Verbesserung der Warenerzeugung oder -verteilung dienen oder
    (c) zur Förderung des technischen Fortschritts beitragen --- \emph{ohne} dass der
    Wettbewerb zwischen den Beteiligten ausgeschaltet wird.
\end{enumerate}

Nach dem \textbf{Selbstprüfungssystem} müssen die Unternehmen selbst einschätzen, ob ihr
Verhalten mit dem Kartellrecht vereinbar ist --- Anmelde- und Überprüfungsverfahren beim
Kartellamt entfallen (weniger Bürokratie, mehr Eigenverantwortung). Bei Verstößen drohen
empfindliche Strafen (Bußgelder, Gewinnabschöpfung).

\subsubsection{Missbrauchsaufsicht (\S\S\,19\,ff. GWB)}

Trotz Wettbewerbspolitik entstehen immer mehr Unternehmen mit erheblicher Marktmacht.
Über \emph{marktbeherrschende} Unternehmen übt das Bundeskartellamt die
\textbf{Missbrauchsaufsicht} aus. Zwei Missbrauchsformen:

\begin{itemize}
  \item \textbf{Ausbeutungsmissbrauch:} Forderung überhöhter Preise von Kunden oder
    extremer Preisnachlässe von abhängigen Zulieferern (viele Zulieferer sind so stark an
    einen Kunden gebunden, dass ihre Existenz bedroht ist, wenn sie ihn verlieren).
  \item \textbf{Behinderungsmissbrauch:} Marktmacht wird gezielt eingesetzt, um aktuelle
    oder potenzielle Konkurrenten zu behindern, zu verdrängen oder ihnen den Marktzugang
    zu versperren --- z.\,B. über \emph{Kampfpreise (Dumping)} oder die Verweigerung des
    Zugangs zu Infrastruktureinrichtungen.
\end{itemize}

\subsubsection{Fusionskontrolle (\S\,35 GWB)}

Während die Missbrauchsaufsicht \emph{bestehende} marktbeherrschende Unternehmen betrifft,
ist die \textbf{Fusionskontrolle} auch eine \emph{vorbeugende} Maßnahme: Geplante Fusionen
müssen beim Bundeskartellamt \emph{angemeldet} werden. Es kann den Zusammenschluss
\textbf{untersagen}, wenn die Gefahr besteht, dass eine marktbeherrschende Stellung
entsteht. Der Bundeswirtschaftsminister kann einen untersagten Zusammenschluss dennoch
erlauben (\textbf{Ministererlaubnis}, \S\,42 GWB), wenn er von gesamtwirtschaftlichem
Vorteil ist und die marktwirtschaftliche Ordnung nicht gefährdet wird.

\begin{beispiel}
Drei Bauunternehmen sprechen vor einer öffentlichen Ausschreibung ab, wer das günstigste
Angebot abgibt und den Auftrag erhält $\Rightarrow$ verbotenes
\textbf{Submissionskartell} (\S\,1 GWB). Ein marktbeherrschender Onlinehändler verkauft
gezielt unter Einstandspreis, um einen neuen Konkurrenten zu verdrängen $\Rightarrow$
\textbf{Behinderungsmissbrauch} (Dumping). Zwei große Lebensmittelketten wollen
fusionieren und hätten danach regional über 50~\% Marktanteil $\Rightarrow$ Anmeldepflicht,
das Bundeskartellamt kann die Fusion \textbf{untersagen}.
\end{beispiel}

\subsection{Eingriffe des Staates in die Preisbildung}

In einer Marktwirtschaft bildet sich der Preis durch Angebot und Nachfrage. In der
\emph{sozialen} Marktwirtschaft übernimmt der Staat regulierende Aufgaben, um
Benachteiligungen von Anbietern oder Nachfragern auszugleichen. Grundsätzlich gilt:

\begin{merksatz}
\textbf{Marktkonforme (indirekte) Eingriffe} beeinflussen Angebot und Nachfrage, ohne den
Marktpreismechanismus außer Kraft zu setzen (z.\,B. Steuern, Zölle, Subventionen).
\textbf{Marktkonträre (direkte) Eingriffe} greifen unmittelbar in die Preisbildung ein
und hebeln den Preismechanismus aus (z.\,B. Höchst- und Mindestpreise).
\end{merksatz}

\begin{center}
\begin{tabular}{p{3.2cm} p{5.2cm} p{5.2cm}}
\toprule
\rowcolor{bfwTabHead} & \textbf{marktkonform (indirekt)} & \textbf{marktkonträr (direkt)}\\
\midrule
Preismechanismus & bleibt in Kraft & wird außer Kraft gesetzt\\
Ansatzpunkt & Angebot bzw. Nachfrage werden verschoben & Preis wird unmittelbar staatlich festgesetzt\\
Beispiele & Steuern, Zölle, Subventionen, Transferleistungen, Mengenmaßnahmen & Höchstpreis, Mindestpreis\\
Typische Folgen & Preissignale bleiben erhalten (aber Verzerrung bei Dauersubventionen) & dauerhafte Überschüsse bzw. Knappheiten, graue und schwarze Märkte, Folgekosten\\
\bottomrule
\end{tabular}
\end{center}

\subsubsection{Indirekte (marktkonforme) Maßnahmen}

\begin{itemize}
  \item \textbf{Verbrauchssteuern:} Höhere Verbrauchssteuern (z.\,B. Tabaksteuer)
    verteuern das Produkt $\Rightarrow$ höherer Gleichgewichtspreis, geringere
    Gleichgewichtsmenge. So kann Verbraucherverhalten aus umwelt- und
    gesundheitspolitischen Gründen gelenkt werden (weniger Raucher). Je nach
    Nachfragereaktion können die Steuereinnahmen trotz höheren Steuersatzes sinken ---
    bei unelastischer Nachfrage steigen sie sogar. Steuersenkungen wirken umgekehrt:
    niedrigerer Gleichgewichtspreis, höhere Menge.
  \item \textbf{Zölle:} Senkung von Einfuhrzöllen erhöht die Importe (wirkt gegen
    Preissteigerungen im Inland); Erhöhung von Einfuhrzöllen verteuert Importwaren und
    schützt die einheimische Industrie. Die Gesetzgebungs- und Ertragskompetenz für Zölle
    liegt bei der \emph{EU}.
  \item \textbf{Subventionen und Steuervergünstigungen (Anbieterförderung):} Staatliche
    Finanzhilfen ermöglichen es Anbietern, zu einem Preis unterhalb der
    Produktionskosten zu verkaufen --- z.\,B. zur Markteinführung neuer Produkte
    (Windkraftanlagen). \emph{Problematisch bei Dauersubvention:} Sie behindert sämtliche
    Preisfunktionen; der künstlich niedrige Preis sendet falsche Signale, das Ausscheiden
    kostenungünstiger Anbieter wird verhindert, Produktionsfaktoren bleiben in nicht
    (mehr) gefragten Produkten gebunden.
  \item \textbf{Transferleistungen (Subvention der Nachfrager):} Der Staat gibt Haushalten
    zusätzliche Mittel, damit sie ein bestimmtes Gut erwerben können --- z.\,B.
    \emph{Wohngeld} oder \emph{Kindergeld}. Die Finanzierung von Subventionen und
    Transfers tragen die Steuerzahler insgesamt.
  \item \textbf{Mengenmaßnahmen:} Regulierung der Angebots- und Nachfragemenge, z.\,B.
    Einfuhrkontingente, Einfuhrverbote, staatliche Vorratshaltung.
\end{itemize}

\subsubsection{Direkte (marktkonträre) Maßnahmen: Mindest- und Höchstpreise}

Beide Preisgrenzen lassen sich im Preis-Mengen-Diagramm darstellen --- der staatlich
festgesetzte Preis „klemmt`` den Markt außerhalb des Gleichgewichts fest:

\begin{center}
\begin{tikzpicture}[font=\footnotesize, scale=0.94]
  % ---------- linkes Diagramm: Mindestpreis ----------
  \begin{scope}
    \node[font=\small\bfseries, color=bfwPrimaryDark, align=center] at (2.7,5.3)
      {Mindestpreis\\[-1pt]{\footnotesize\color{bfwInkSoft}\mdseries Schutz der Anbieter}};
    \draw[-{Stealth[length=2.5mm]}, semithick, color=bfwInk] (0,0) -- (0,4.55) node[above] {$p$};
    \draw[-{Stealth[length=2.5mm]}, semithick, color=bfwInk] (0,0) -- (5.3,0) node[right] {$x$};
    \draw[bfwPrimaryDark, very thick] (0.4,0.5) -- (4.4,3.9) node[above] {$A$};
    \draw[bfwWarn!90!black, very thick] (0.4,3.9) -- (4.4,0.5) node[right] {$N$};
    \draw[densely dashed, color=bfwInkSoft] (0,2.2) -- (2.4,2.2) -- (2.4,0);
    \node[left, color=bfwInkSoft] at (0,2.2) {$p_0$};
    \node[below, color=bfwInkSoft] at (2.4,0) {$x_0$};
    \fill[bfwInk] (2.4,2.2) circle (1.7pt);
    \draw[bfwAccent, line width=1.5pt] (0,3.05) -- (5.05,3.05);
    \node[left, color=bfwAccent!80!black] at (0,3.05) {$p_{\mathrm{min}}$};
    \draw[densely dashed, color=bfwInkSoft] (1.4,3.05) -- (1.4,0);
    \draw[densely dashed, color=bfwInkSoft] (3.4,3.05) -- (3.4,0);
    \node[below, color=bfwInkSoft] at (1.4,0) {$x_N$};
    \node[below, color=bfwInkSoft] at (3.4,0) {$x_A$};
    \fill[bfwAccent!80!black] (1.4,3.05) circle (1.5pt) (3.4,3.05) circle (1.5pt);
    \draw[{Stealth[length=2mm]}-{Stealth[length=2mm]}, color=bfwAccent!80!black, semithick]
      (1.4,3.5) -- (3.4,3.5)
      node[midway, above, align=center] {\bfseries Angebots-\\[-2pt]\bfseries überschuss};
  \end{scope}
  % ---------- rechtes Diagramm: Höchstpreis ----------
  \begin{scope}[xshift=7.5cm]
    \node[font=\small\bfseries, color=bfwWarn!85!black, align=center] at (2.7,5.3)
      {Höchstpreis\\[-1pt]{\footnotesize\color{bfwInkSoft}\mdseries Schutz der Nachfrager}};
    \draw[-{Stealth[length=2.5mm]}, semithick, color=bfwInk] (0,0) -- (0,4.55) node[above] {$p$};
    \draw[-{Stealth[length=2.5mm]}, semithick, color=bfwInk] (0,0) -- (5.3,0) node[right] {$x$};
    \draw[bfwPrimaryDark, very thick] (0.4,0.5) -- (4.4,3.9) node[above] {$A$};
    \draw[bfwWarn!90!black, very thick] (0.4,3.9) -- (4.4,0.5) node[right] {$N$};
    \draw[densely dashed, color=bfwInkSoft] (0,2.2) -- (2.4,2.2) -- (2.4,0);
    \node[left, color=bfwInkSoft] at (0,2.2) {$p_0$};
    \node[below, color=bfwInkSoft] at (2.4,0) {$x_0$};
    \fill[bfwInk] (2.4,2.2) circle (1.7pt);
    \draw[bfwAccent, line width=1.5pt] (0,1.35) -- (5.05,1.35);
    \node[left, color=bfwAccent!80!black] at (0,1.35) {$p_{\mathrm{max}}$};
    \draw[densely dashed, color=bfwInkSoft] (1.4,1.35) -- (1.4,0);
    \draw[densely dashed, color=bfwInkSoft] (3.4,1.35) -- (3.4,0);
    \node[below, color=bfwInkSoft] at (1.4,0) {$x_A$};
    \node[below, color=bfwInkSoft] at (3.4,0) {$x_N$};
    \fill[bfwAccent!80!black] (1.4,1.35) circle (1.5pt) (3.4,1.35) circle (1.5pt);
    \draw[{Stealth[length=2mm]}-{Stealth[length=2mm]}, color=bfwAccent!80!black, semithick]
      (1.4,0.75) -- (3.4,0.75)
      node[midway, fill=white, inner sep=1.5pt, align=center]
      {\bfseries Nachfrage-\\[-2pt]\bfseries überschuss};
  \end{scope}
\end{tikzpicture}
\end{center}

\begin{beispiel}
\textbf{Durchgerechnet:} Auf einem Markt gelte die Nachfragefunktion $x_N = 100 - 2p$
und die Angebotsfunktion $x_A = -20 + 2p$ ($p$ in €, $x$ in Stück).\\[0.3em]
\textbf{Schritt 1 --- Gleichgewicht:} $100 - 2p = -20 + 2p \Rightarrow 4p = 120
\Rightarrow p_0 = 30$~€; eingesetzt: $x_0 = 100 - 60 = 40$ Stück.\\
\textbf{Schritt 2 --- Mindestpreis 35~€ (über $p_0$):} $x_N = 100 - 70 = 30$,
$x_A = -20 + 70 = 50$ $\Rightarrow$ \textbf{Angebotsüberschuss von 20 Stück}
($x_N < x_0 < x_A$).\\
\textbf{Schritt 3 --- Höchstpreis 25~€ (unter $p_0$):} $x_N = 100 - 50 = 50$,
$x_A = -20 + 50 = 30$ $\Rightarrow$ \textbf{Nachfrageüberschuss von 20 Stück}
($x_A < x_0 < x_N$).
\end{beispiel}

\subsubsection{Mindestpreis}

Ein \textbf{Mindestpreis} ist eine staatlich verordnete \emph{Preisuntergrenze}, die
nicht unterschritten werden darf; er liegt \emph{über} dem Gleichgewichtspreis.
\textbf{Ziel: Schutz der Anbieter.} Folgen: Das Angebot wird ausgeweitet, die Nachfrage
sinkt $\Rightarrow$ \textbf{Angebotsüberschuss} ($x_N < x_0 < x_A$). Die Überschüsse
stellen den Staat vor neue Probleme --- er muss mengenregulierend eingreifen
(Kontingente, Prämien für Angebotsreduzierung, Flächenstilllegung). Bei
Mengenkontingenten kann ein \textbf{grauer Markt} entstehen, auf dem die Mehrproduktion
unter dem festgesetzten Mindestpreis angeboten wird.

\begin{beispiel}
\textbf{EU-Agrarmarkt:} Landwirten wurden jahrzehntelang Mindestpreise mit
Abnahmegarantien zugestanden --- Folge: massive Überproduktion („Butterberge``,
„Milchseen``). Der Staat musste die Überschüsse durch Einlagerung, Verarbeitung,
Veräußerung oder Vernichtung beseitigen (erhebliche Kosten) und das Angebot durch
Produktionsquoten und Stilllegungsprämien eindämmen.
\textbf{Gesetzlicher Mindestlohn:} Er hebelt die Findung des Gleichgewichtslohns aus und
kann zu einem dauerhaften Angebotsüberhang an Arbeitskräften (Arbeitslosigkeit) führen,
wenn Arbeitgeber Arbeitskräfte durch Technik substituieren oder Arbeitsumfänge ins
Ausland verlagern.
\end{beispiel}

\subsubsection{Höchstpreis}

Ein \textbf{Höchstpreis} ist eine staatlich verordnete \emph{Preisobergrenze}; er liegt
\emph{unter} dem Gleichgewichtspreis. \textbf{Ziel: Schutz der Nachfrager.} Beispiele:
Mietpreisbindung im sozialen Wohnungsbau, \emph{Mietpreisbremse} auf angespannten
Wohnungsmärkten. Folgen: Die Nachfrage steigt, das Angebot sinkt $\Rightarrow$
\textbf{Nachfrageüberschuss} ($x_A < x_0 < x_N$). Damit der Höchstpreis wirkt, sind
weitere Maßnahmen nötig: \emph{Subventionen}, damit überhaupt ein Angebot vorhanden ist,
sowie \emph{Kontingentierung und Rationierung} der knappen Menge. Häufig bilden sich
\textbf{Schwarzmärkte}, auf denen die Güter zu höheren Preisen verkauft werden. Selten
kann das verminderte Angebot über \emph{Importe} gelöst werden, wenn ausländische
Anbieter aufgrund von Kostenvorteilen unter dem Höchstpreis anbieten können.

\begin{eselsbruecke}
\textbf{Mi}ndestpreis liegt \textbf{über}, \textbf{Hö}chstpreis liegt \textbf{unter} dem
Gleichgewichtspreis --- sonst wären sie wirkungslos! Mindestpreis schützt die
\textbf{A}nbieter $\Rightarrow$ \textbf{A}ngebotsüberschuss; Höchstpreis schützt die
\textbf{N}achfrager $\Rightarrow$ \textbf{N}achfrageüberschuss. Merkhilfe: \emph{Die
geschützte Marktseite weitet ihre Menge aus --- auf ihrer Seite entsteht der Überschuss.}
\end{eselsbruecke}

\begin{merksatz}[title={\bfseries Beurteilung staatlicher Preiseingriffe (prüfungsrelevant!)}]
Staatliche Preisfestsetzungen setzen die Funktionen des Marktpreises (Signal, Ausgleich,
Allokation, Selektion) außer Kraft: Es entstehen dauerhafte Überschüsse bzw. Knappheiten,
Folgekosten für den Staat, Fehlallokation von Produktionsfaktoren sowie graue und schwarze
Märkte. Marktkonforme Eingriffe (Steuern, Subventionen, Transfers) sind daher
ordnungspolitisch vorzuziehen --- aber auch Dauersubventionen verzerren die Preissignale.
\end{merksatz}

\begin{merksatz}[title={\bfseries Wichtig für die Prüfungsvorbereitung (Teil 1)}]
Sie sollten sicher können: die \emph{Funktionen des Wettbewerbs} erläutern, die
\emph{drei Säulen des GWB} mit Kartellarten, Ausnahmen und Beispielen erklären,
\emph{marktkonforme und marktkonträre} Eingriffe unterscheiden sowie Höchst- und
Mindestpreise \emph{im Preis-Mengen-Diagramm} einzeichnen und ihre Folgen
(Überschüsse, graue und schwarze Märkte, Folgekosten) beurteilen. Üben Sie mit dem
Aufgabenheft (Aufgabe~4: Mindestpreis) und dem Quiz zu Teil~1!
\end{merksatz}

\vspace{1em}
\hrule
\vspace{0.8em}
\noindent\textsf{\small Quellen: Rahmenplan Wirtschaftsbezogene Qualifikationen (DIHK), Abschnitte 1.1.1.2, 1.1.1.3;
Textband Volkswirtschaft (DIHK-Bildungs-GmbH), Kap.~1.2--1.3 (S.~12--18).
Weiter geht es in Session~2 (Teil~2): Volkswirtschaftliche Gesamtrechnung \& Ziele der Stabilitätspolitik.}
\par\medskip{\bfwlabelfont\color{bfwPrimary}SESSION 2\par}\smallskip
\begin{merksatz}[title={\bfseries Worum geht es in Teil 2?}]
Ob Wettbewerbs- und Wirtschaftspolitik erfolgreich sind, misst die
\emph{Volkswirtschaftliche Gesamtrechnung} (BIP, BNE, Volkseinkommen) --- und die Ziele,
an denen sich die Wirtschaftspolitik ausrichten muss, formuliert das Stabilitätsgesetz
im \emph{magischen Viereck}. Dazu kommen die Primär- und Sekundärverteilung des
Volkseinkommens mit den prüfungsrelevanten Rechnungen (Lohnquote, verfügbares
Einkommen). Alle Themen sind Standardstoff der IHK-Prüfung und bauen auf Teil~1
(Wettbewerbspolitik \& Staatseingriffe, Session~1) auf.
\end{merksatz}



\subsection{Volkswirtschaftliche Gesamtrechnung (VGR)}

\subsubsection{Aufgaben der VGR}

Die \textbf{VGR} ist das statistische Instrument zur zahlenmäßigen Erfassung
gesamtwirtschaftlicher Vorgänge. Sie gibt \emph{rückschauend} für eine abgeschlossene
Periode Auskunft über \textbf{Entstehung, Verwendung und Verteilung} des
Bruttoinlandsprodukts (BIP) und des Bruttonationaleinkommens (BNE). Grundlage ist die
\textbf{Kreislauftheorie} (Wirtschaftskreislauf) mit den Sektoren private Haushalte,
Unternehmen, Staat, Kapitalsammelstellen und Ausland. Die Daten dienen u.\,a. zur
Ermittlung von Wirtschaftswachstum und Strukturwandel, als Konjunkturbarometer, als
Grundlage wirtschaftspolitischer Entscheidungen, als Anhaltspunkt für Lohnverhandlungen
der Tarifpartner und für den internationalen Vergleich. Für EU-Staaten gilt das
\textbf{Europäische System volkswirtschaftlicher Gesamtrechnungen (ESVG 2010)} --- es
sichert die Vergleichbarkeit der Daten zwischen den EU-Ländern.

\subsubsection{BIP und BNE: Inlandskonzept vs. Inländerkonzept}

\begin{itemize}
  \item Das \textbf{Bruttoinlandsprodukt (BIP)} misst den Wert der Güter und
    Dienstleistungen, die innerhalb eines Zeitraums \emph{innerhalb der geografischen
    Grenzen} einer Volkswirtschaft produziert wurden (\textbf{Inlandskonzept}) ---
    bewertet zu Marktpreisen. Es spielt keine Rolle, ob die Produzenten ihren Wohnsitz
    im Inland haben.
  \item Das \textbf{Bruttonationaleinkommen (BNE)} ist das Einkommen der \emph{Inländer},
    das sie durch Bereitstellung von Produktionsfaktoren im In- \emph{und} Ausland
    erzielt haben (\textbf{Inländerkonzept}) --- z.\,B. Pendler, Monteure, Künstler,
    Auslandsproduktion deutscher Unternehmen. „Inländer`` sind alle natürlichen Personen
    mit ständigem Wohnsitz im Inland (unabhängig von der Staatsbürgerschaft) und alle im
    Inland wirtschaftlich tätigen Unternehmen (auch in ausländischem Besitz).
\end{itemize}

\begin{merksatz}[title={\bfseries Formeln: Bruttowertschöpfung und BNE}]
\textbf{Bruttowertschöpfung} = (Brutto-)Produktionswerte $-$ Vorleistungen
\emph{(vermeidet Doppelzählungen!)}\\[0.4em]
\textbf{BIP} $-$ Inlandseinkommen von Nichtinländern $+$ Auslandseinkommen von Inländern
= \textbf{BNE}\\
(die beiden Korrekturposten bilden den \emph{Saldo der Primäreinkommen aus der übrigen
Welt})
\end{merksatz}

\begin{beispiel}
Ein Küchengerätehersteller bezieht Elektromotoren von einem Zulieferer. Würde man beim
BIP sowohl die Leistung des Motorenherstellers als auch die des Küchengeräteherstellers
voll zählen, wäre der Motor \emph{doppelt} erfasst --- denn im Preis der Rührmaschine
steckt der Motorpreis bereits. Deshalb werden von den Produktionswerten die
\textbf{Vorleistungen} abgezogen.
\end{beispiel}

Inlandsprodukt und Nationaleinkommen können \emph{brutto} (inkl. Abschreibungen) oder
\emph{netto} (nach Abzug der Abschreibungen) dargestellt werden.

\subsubsection{Nominales und reales BIP}

Das \textbf{nominale BIP} wird in laufenden/aktuellen Preisen (Marktpreisen) ausgewiesen.
Steigt es um 5~\%, kann das an gestiegenen \emph{Mengen}, an gestiegenen \emph{Preisen}
(Inflation) oder an beidem liegen. Um den \emph{echten} Leistungszuwachs zu ermitteln,
errechnet das Statistische Bundesamt das \textbf{reale BIP}: Es legt die Preise eines
\emph{festgelegten Basisjahres} zugrunde und eliminiert so die Preissteigerungen. Die
Wachstumsrate des realen BIP zeigt an, um wie viel Prozent die Produktion im Vergleich
zum Vorjahr gestiegen ist. Im internationalen Vergleich werden stets die Werte des
\emph{realen} BIP bzw. BNE zugrunde gelegt.

\begin{achtung}
Typische Prüfungsfalle: Wirtschaftswachstum wird \emph{immer} an der Zuwachsrate des
\textbf{realen} BIP gemessen. Ein nominaler BIP-Anstieg von 5~\% bei 5~\% Inflation
bedeutet: reales Wachstum \textbf{null} --- die Wirtschaft hat nicht mehr produziert,
nur die Preise sind gestiegen (Rechenweg in Abschnitt~4.4).
\end{achtung}

\subsubsection{Entstehung, Verwendung und Verteilung des BIP}

Das BIP kann über drei Rechnungen ermittelt werden --- alle führen zum gleichen Ergebnis:

\begin{center}
\begin{tabular}{p{3.4cm} p{4.6cm} p{6.4cm}}
\toprule
\textbf{Rechnung} & \textbf{Leitfrage} & \textbf{Schema (vereinfacht)}\\
\midrule
Entstehungs\-rechnung & \emph{Wo} ist die Wirtschaftsleistung entstanden? & Bruttowertschöpfung der Wirtschaftsbereiche $+$ Gütersteuern abzgl. Gütersubventionen $=$ BIP\\
Verwendungs\-rechnung & \emph{Wofür} wurde die Inlandsproduktion verwendet? & Private Konsumausgaben $+$ Konsumausgaben des Staates $+$ Bruttoinvestitionen $+$ Außenbeitrag (Exporte $-$ Importe) $=$ BIP\\
Verteilungs\-rechnung & \emph{An wen} fließen die Einkommen der Inländer? & Arbeitnehmerentgelt $+$ Unternehmens- und Vermögenseinkommen $=$ \textbf{Volkseinkommen}; $+$ Produktions-/Importabgaben abzgl. Subventionen $+$ Abschreibungen $-$ Saldo der Primäreinkommen aus der übrigen Welt $=$ BIP\\
\bottomrule
\end{tabular}
\end{center}

Hinweise zur Entstehungsrechnung: Leistungen der Unternehmen werden zu \emph{Marktpreisen}
bewertet; die (unverkaufte) Produktionsleistung des \emph{Staates} (Bildung, Polizei) zu
den \emph{Kosten der eingesetzten Produktionsfaktoren}. Bei privaten Haushalten zählen nur
Leistungen, für die Produktionskosten durch fremde Personen (z.\,B. Haushaltshilfen)
entstanden sind. Bei der Verteilungsrechnung wechselt die Bezugsgrundlage vom
\emph{Inland} zu den \emph{Inländern}.

\begin{center}
\begin{tikzpicture}[font=\small,
    rechnung/.style={rounded corners=4pt, thick, align=center, text width=3.9cm,
                     minimum height=3.1cm, inner sep=6pt}]
  \node[rechnung, draw=bfwPrimaryDark, fill=bfwBgTeal] (re) at (-4.9,0)
    {{\color{bfwPrimaryDark}\bfseries Entstehungsrechnung}\\[1pt]
     {\footnotesize\itshape Wo ist die Leistung entstanden?}\\[4pt]
     {\footnotesize Bruttowertschöpfung der Wirtschaftsbereiche
      $+$ Gütersteuern $-$ Gütersubventionen}};
  \node[rechnung, draw=bfwAccent!85!black, fill=bfwBgAmber] (rv) at (0,0)
    {{\color{bfwAccent!85!black}\bfseries Verwendungsrechnung}\\[1pt]
     {\footnotesize\itshape Wofür wurde die Produktion verwendet?}\\[4pt]
     {\footnotesize privater Konsum $+$ Staatskonsum $+$ Bruttoinvestitionen
      $+$ Außenbeitrag (Ex $-$ Im)}};
  \node[rechnung, draw=bfwWarn!85!black, fill=bfwBgRose] (rt) at (4.9,0)
    {{\color{bfwWarn!85!black}\bfseries Verteilungsrechnung}\\[1pt]
     {\footnotesize\itshape An wen fließen die Einkommen?}\\[4pt]
     {\footnotesize Arbeitnehmerentgelt $+$ Unternehmens- und Vermögenseinkommen
      $=$ Volkseinkommen}};
  \node[roughbox, minimum width=7.4cm, minimum height=1.05cm, align=center] (bip) at (0,-3.3)
    {\bfseries BIP --- ein und dasselbe Ergebnis};
  \draw[roughline] (re.south) to[out=-90, in=180] (bip.west);
  \draw[roughline] (rv.south) -- (bip.north);
  \draw[roughline] (rt.south) to[out=-90, in=0] (bip.east);
  \node[below=3pt of bip, font=\footnotesize, color=bfwInkSoft, align=center]
    {Alle drei Rechnungen führen zum gleichen Bruttoinlandsprodukt (ESVG 2010)};
\end{tikzpicture}
\end{center}

\begin{beispiel}
\textbf{Verwendungsrechnung durchgerechnet:} Private Konsumausgaben 1.900~Mrd.~€,
Konsumausgaben des Staates 800~Mrd.~€, Bruttoinvestitionen 500~Mrd.~€, Exporte
1.600~Mrd.~€, Importe 1.300~Mrd.~€.\\[0.3em]
\textbf{Schritt 1:} Außenbeitrag $=$ Exporte $-$ Importe $= 1.600 - 1.300 = 300$~Mrd.~€.\\
\textbf{Schritt 2:} BIP $= 1.900 + 800 + 500 + 300 = \textbf{3.500~Mrd.~€}$.\\[0.3em]
Dieselben 3.500~Mrd.~€ ergeben sich über die Entstehungs- und die Verteilungsrechnung ---
und dienen in Abschnitt~4.1 als Ausgangswert für die BNE-Berechnung.
\end{beispiel}

\subsubsection{Grenzen der Aussagekraft des BIP}

\begin{itemize}
  \item \textbf{Eigenleistungen} privater Haushalte gelten nicht als Produktion;
    \textbf{ehrenamtliche} und andere unentgeltliche Tätigkeiten (z.\,B. Hausarbeit)
    bleiben unberücksichtigt.
  \item Die \textbf{Schwarzarbeit} (Teil der \emph{Schattenwirtschaft}) wird vom
    Statistischen Bundesamt nur \emph{geschätzt}.
  \item Erfasst werden auch Transaktionen, die den \emph{Wohlstand nicht erhöhen} (z.\,B.
    Reparatur von Unfall- und Umweltschäden zählt positiv zum BIP); umgekehrt bleiben
    wohlstandserhöhende, aber unbezahlte Leistungen unerfasst.
  \item Das BIP sagt nichts über die \emph{Verteilung} des Wohlstands, Umweltbelastungen
    oder Lebensqualität aus --- ein hohes Pro-Kopf-Einkommen schließt Armut von Teilen
    der Bevölkerung nicht aus.
\end{itemize}

\subsection{Primär- und Sekundärverteilung des Volkseinkommens}

\subsubsection{Primärverteilung: Lohn- und Gewinnquote}

Das über die Verteilungsrechnung errechnete \textbf{Volkseinkommen} ist die Summe aller
Arbeitnehmerentgelte sowie Unternehmens- und Vermögenseinkommen, die Inländern aus dem
In- und Ausland zugeflossen sind. Diese Verteilung heißt \textbf{Primärverteilung}, die
Einkommen heißen \emph{Primäreinkommen}. Sie basiert auf dem \textbf{Leistungsprinzip}:
Grundlage ist die wirtschaftliche Betätigung und die Bereitstellung von
Produktionsfaktoren. Man betrachtet die Verteilung \emph{funktional} (nach
Produktionsfaktoren: Lohnquote, Gewinnquote) oder \emph{personell} (nach Haushalten:
Pro-Kopf-Einkommen, Lorenzkurve).

\begin{merksatz}[title={\bfseries Formeln: Lohnquote und Gewinnquote}]
\[
\text{Lohnquote in \%} = \frac{\text{Arbeitnehmerentgelt}}{\text{Volkseinkommen}} \cdot 100
\qquad
\text{Gewinnquote in \%} = \frac{\text{Unternehmens- und Vermögenseinkommen}}{\text{Volkseinkommen}} \cdot 100
\]
2018 betrug die Lohnquote rund 69~\%, die Gewinnquote rund 31~\%. Beide Quoten ergänzen
sich stets zu 100~\%.
\end{merksatz}

Eine \emph{steigende Lohnquote} deutet an, dass der Anteil des Arbeitnehmereinkommens
zulasten der Unternehmens- und Vermögenseinkommen gestiegen ist. Vorsicht bei der
Interpretation: Auch Arbeitnehmer haben Zinseinkünfte (zählen zur Gewinnquote), und
Vorstandsgehälter zählen zum Arbeitnehmerentgelt (Lohnquote).

Die \textbf{personelle Einkommensverteilung} fragt, wie gleichmäßig das Volkseinkommen
auf die privaten Haushalte verteilt ist --- unabhängig von der Einkommensquelle.
Vereinfachter Verteilungsschlüssel ist das \textbf{Pro-Kopf-Einkommen} (Volkseinkommen
$\div$ Einwohnerzahl), das häufig als Wohlstandsindikator im Ländervergleich dient. Die
Verteilung wird mit der \textbf{Lorenzkurve} dargestellt: Sie zeigt, wie viel Prozent der
Einkommensempfänger wie viel Prozent des Volkseinkommens verdienen. Je weiter die Kurve
von der Diagonalen (Gleichverteilung) entfernt ist, desto ungleicher die Verteilung.

\begin{center}
\begin{tikzpicture}[font=\footnotesize]
  % Fläche A zwischen Diagonale und Lorenzkurve
  \fill[bfwAccent!30] (0,0) -- (4.5,4.5) .. controls (4.1,2.2) and (2.8,0.5) .. (0,0);
  % Hilfsrahmen (100 %-Quadrat)
  \draw[densely dashed, color=bfwInkSoft!60] (4.5,0) -- (4.5,4.5) -- (0,4.5);
  % Achsen
  \draw[-{Stealth[length=2.5mm]}, semithick, color=bfwInk] (0,0) -- (0,5.15);
  \draw[-{Stealth[length=2.5mm]}, semithick, color=bfwInk] (0,0) -- (5.5,0);
  \node[right, align=left, color=bfwInk] at (0.15,5.05)
    {Anteil am Einkommen (kumuliert)};
  \node[below, color=bfwInk] at (3.0,-0.55) {Anteil der Haushalte (kumuliert)};
  % Gleichverteilungsgerade
  \draw[color=bfwInkSoft, thick, densely dashed] (0,0) -- (4.5,4.5)
    node[pos=0.62, sloped, above, color=bfwInkSoft] {Gleichverteilungsgerade (45°)};
  % Lorenzkurve
  \draw[bfwPrimaryDark, line width=1.5pt]
    (0,0) .. controls (2.8,0.5) and (4.1,2.2) .. (4.5,4.5);
  \node[color=bfwPrimaryDark, font=\footnotesize\bfseries] at (3.55,1.15) {Lorenzkurve};
  % Flächen-Labels
  \node[color=bfwAccent!70!black, font=\footnotesize\bfseries] at (2.45,1.7) {$A$};
  \node[color=bfwInkSoft, font=\footnotesize\bfseries] at (4.05,2.2) {$B$};
  % Ablesebeispiel: 50 % der Haushalte -> ca. 18 % des Einkommens
  \draw[densely dashed, color=bfwWarn!80!black] (2.25,0) -- (2.25,0.83) -- (0,0.83);
  \fill[bfwWarn!80!black] (2.25,0.83) circle (1.5pt);
  \node[below, color=bfwInkSoft] at (2.25,0) {50\,\%};
  \node[left, color=bfwWarn!80!black] at (0,0.83) {$\approx$\,18\,\%};
  \node[below, color=bfwInkSoft] at (4.5,0) {100\,\%};
  \node[left, color=bfwInkSoft] at (0,4.5) {100\,\%};
\end{tikzpicture}
\end{center}

Aus der Lorenzkurve wird der \textbf{Gini-Koeffizient} abgeleitet: das Verhältnis der
Fläche~$A$ (zwischen Gleichverteilungsgerade und Lorenzkurve) zur gesamten Fläche unter
der Diagonalen ($A + B$). Er liegt zwischen \textbf{0} (völlige Gleichverteilung ---
Lorenzkurve und Diagonale fallen zusammen) und \textbf{1} (maximale Ungleichverteilung
--- ein Haushalt erhält alles). Je größer der Gini-Koeffizient, desto ungleicher die
Verteilung. Im Ablesebeispiel der Grafik erhalten die „unteren`` 50~\% der Haushalte
zusammen nur rund 18~\% des Einkommens.

\subsubsection{Sekundärverteilung: Einkommensumverteilung}

Die Primärverteilung ist ungleich und wird häufig als ungerecht empfunden. Der Staat
korrigiert sie durch die \textbf{Sekundärverteilung} nach dem \textbf{Bedarfsprinzip}:
staatliche Umverteilung von Teilen des Primäreinkommens mit dem Ziel einer gerechteren
Einkommensverteilung. Ansatzpunkte:

\begin{itemize}
  \item \textbf{Progressive Einkommensteuer:} Wer über die Primärverteilung ein hohes
    Einkommen erzielt, wird zu einem höheren Prozentsatz besteuert.
  \item \textbf{Sozialversicherungsbeiträge} und \textbf{Sozialleistungen des Staates:}
    Wer keine Faktorleistung erbringen kann, erhält \emph{Transfereinkommen} --- z.\,B.
    Arbeitslosengeld, Renten, Wohngeld, Kindergeld.
\end{itemize}

\begin{merksatz}[title={\bfseries Formel: Verfügbares Einkommen der privaten Haushalte}]
\begin{tabular}{ll}
  & Faktoreinkommen der privaten Haushalte (Primäreinkommen)\\
$-$ & direkte Steuern\\
$-$ & Sozialversicherungsbeiträge\\
$+$ & Sozialleistungen des Staates (Transferzahlungen)\\
\midrule
$=$ & \textbf{verfügbares Einkommen} der privaten Haushalte
\end{tabular}\\[0.4em]
Das verfügbare Einkommen steht für \emph{Konsum} und \emph{Sparen} zur Verfügung ---
es ist das Ergebnis der Sekundärverteilung.
\end{merksatz}

Auch die \textbf{Gewerkschaften} beeinflussen die Einkommensverteilung über die
Tarifpolitik (Tariflöhne, Lohnsteigerungen über dem Produktivitätsfortschritt,
überproportionale Steigerungen für untere Lohngruppen). Zu hohe Forderungen können jedoch
die Stabilität der Unternehmen, ihre internationale Wettbewerbsfähigkeit und damit
Arbeitsplätze gefährden.

\subsection{Ziele der Stabilitätspolitik}

\subsubsection{Das Stabilitätsgesetz und das magische Viereck}

Das \textbf{Gesetz zur Förderung der Stabilität und des Wachstums der Wirtschaft (StabG)}
von 1967 gilt als \emph{Grundgesetz der Wirtschaftspolitik}. Nach \S\,1 StabG haben Bund
und Länder bei ihren wirtschafts- und finanzpolitischen Maßnahmen die Erfordernisse des
\textbf{gesamtwirtschaftlichen Gleichgewichts} zu beachten. Die Maßnahmen sollen im Rahmen
der marktwirtschaftlichen Ordnung \emph{gleichzeitig} beitragen zu:

\begin{enumerate}
  \item \textbf{Stabilität des Preisniveaus},
  \item \textbf{hohem Beschäftigungsstand},
  \item \textbf{außenwirtschaftlichem Gleichgewicht},
  \item \textbf{stetigem und angemessenem Wirtschaftswachstum}.
\end{enumerate}

Da es nahezu unmöglich ist, alle vier Ziele gleichzeitig vollständig zu erreichen, spricht
man vom \textbf{magischen Viereck}. Seine Ziele sind \emph{quantitativ}, d.\,h. der Grad
der Zielerreichung ist über Kennzahlen \emph{messbar}. Ergänzt um die \emph{qualitativen}
Ziele \textbf{Erhaltung einer lebenswerten Umwelt (Umweltschutz)} und \textbf{gerechte
Einkommens- und Vermögensverteilung} spricht man vom \textbf{magischen Sechseck} --- für
die qualitativen Ziele gibt es keinen klaren Maßstab der Zielerreichung.

\begin{center}
\begin{tikzpicture}[font=\footnotesize,
    ziel/.style={rounded corners=4pt, thick, align=center, text width=3.6cm, inner sep=5pt}]
  \node[ziel, draw=bfwPrimaryDark, fill=bfwBgTeal] (z1) at (0,2.9)
    {{\bfseries Preisniveaustabilität}\\[1pt]
     {\scriptsize Inflationsrate 2\,\% (symmetrisches EZB-Ziel)}};
  \node[ziel, draw=bfwAccent!85!black, fill=bfwBgAmber] (z2) at (5.3,0)
    {{\bfseries Hoher Beschäftigungsstand}\\[1pt]
     {\scriptsize Arbeitslosenquote ca. 3--4\,\%}};
  \node[ziel, draw=bfwWarn!85!black, fill=bfwBgRose] (z3) at (0,-2.9)
    {{\bfseries Stetiges u. angemessenes Wachstum}\\[1pt]
     {\scriptsize reales BIP im Schnitt $+2$\,\% p.\,a.}};
  \node[ziel, draw=bfwPrimaryDark, fill=bfwBgTeal] (z4) at (-5.3,0)
    {{\bfseries Außenwirtschaftliches Gleichgewicht}\\[1pt]
     {\scriptsize Außenbeitrag $+1$ bis $2$\,\% des BIP}};
  % Rautenkanten (Spannungsboegen)
  \draw[densely dashed, thick, color=bfwInkSoft!65] (z1) -- (z2) -- (z3) -- (z4) -- (z1);
  \node[align=center, font=\small\bfseries, color=bfwPrimaryDark] at (0,0.4)
    {Magisches Viereck\\ (\S\,1 StabG)};
  \node[align=center, font=\scriptsize, color=bfwInkSoft] at (0,-0.55)
    {alle vier Ziele sind gleichzeitig\\ kaum vollständig erreichbar};
  % Zielkonflikt: Preisstabilitaet <-> Wachstum
  \draw[bfwWarn, thick, {Stealth[length=2mm]}-{Stealth[length=2mm]},
        decorate, decoration={zigzag, segment length=7pt, amplitude=1.3pt,
                              pre length=6pt, post length=6pt}]
    (-1.7,2.0) -- (-1.7,-2.0);
  \node[rotate=90, font=\scriptsize\bfseries, color=bfwWarn] at (-2.05,0) {Zielkonflikt};
  % Zielharmonie: Wachstum <-> Beschaeftigung
  \draw[bfwSuccess!65!black, thick, {Stealth[length=2mm]}-{Stealth[length=2mm]}]
    (1.7,-1.62) -- (3.15,-0.86)
    node[midway, sloped, above, font=\scriptsize\bfseries] {Zielharmonie};
  % Legende
  \node[font=\scriptsize, color=bfwInkSoft, align=center] at (0,-4.15)
    {gestrichelte Kanten: Spannungsverhältnis („magisch``) \textperiodcentered{}
     \textcolor{bfwWarn}{Zickzack: konkurrierende Ziele} \textperiodcentered{}
     \textcolor{bfwSuccess!65!black}{Doppelpfeil: komplementäre Ziele}};
\end{tikzpicture}
\end{center}

Der \textbf{Jahreswirtschaftsbericht} (nach \S\,2 StabG zu Jahresbeginn dem Bundestag
vorzulegen) informiert über die wirtschaftlichen Zukunftserwartungen der Bundesregierung,
formuliert die wirtschafts- und finanzpolitischen Ziele für das laufende Jahr und nimmt
Stellung zum Jahresgutachten des \emph{Sachverständigenrates}.

\subsubsection{Die vier Ziele und ihre Messgrößen}

\begin{center}
\begin{tabular}{p{4.2cm} p{4.6cm} p{5.4cm}}
\toprule
\textbf{Ziel} & \textbf{Messgröße} & \textbf{Zielwert (Faustregel)}\\
\midrule
Hoher Beschäftigungsstand (Vollbeschäftigung) & Arbeitslosenquote (Bundesagentur für Arbeit) & erreicht bei ca. 3--4~\% (Sockelarbeitslosigkeit bleibt)\\
Preisniveaustabilität & Inflationsrate auf Basis des Verbraucherpreisindex (Warenkorb, Statistisches Bundesamt) & Inflationsrate von 2~\% (symmetrisches Ziel der EZB seit 2021, HVPI)\\
Außenwirtschaftliches Gleichgewicht & Außenbeitrag (Saldo aus Handels- und Dienstleistungsbilanz) & positiver Außenbeitrag von 1--2~\% des nominalen BIP\\
Stetiges und angemessenes Wirtschaftswachstum & jährliche Zuwachsrate des \emph{realen} BIP & Durchschnitt rund 2~\% pro Jahr, ohne heftige Konjunkturausschläge\\
\bottomrule
\end{tabular}
\end{center}

Erläuterungen: Die \emph{Arbeitslosenquote} setzt registrierte Arbeitslose ins Verhältnis
zu den zivilen Erwerbspersonen. \emph{Preisniveaustabilität} bedeutet nicht, dass kein
Preis steigen darf --- einzelne Preissteigerungen sollen durch Preissenkungen bei anderen
Gütern ausgeglichen werden. Der \emph{Außenbeitrag} (Exporte $-$ Importe von Waren und
Dienstleistungen) zeigt bei positivem Wert, dass die Volkswirtschaft mehr produziert als
verbraucht hat (Netto-Güterexport). \emph{Wirtschaftswachstum} wird stets an der
Zuwachsrate des realen BIP gemessen.

\begin{eselsbruecke}
Magisches Viereck: „\textbf{P}reise, \textbf{A}rbeit, \textbf{W}achstum,
\textbf{A}ußenhandel`` --- \textbf{PAWA}. Für das Sechseck kommen \textbf{U}mwelt und
\textbf{V}erteilung hinzu.
\end{eselsbruecke}

\subsubsection{Zielkonflikte und Zielharmonien}

Die Ziele stehen in unterschiedlichen Beziehungen zueinander --- deshalb „magisch``:

\begin{itemize}
  \item \textbf{Konkurrierende Zielbeziehungen (Zielkonflikte):}
    \emph{Wirtschaftswachstum vs. Preisniveaustabilität} (starkes Wachstum zieht i.\,d.\,R.
    Preiserhöhungen nach sich); \emph{Wirtschaftswachstum vs. Ökologie} (Wachstum um jeden
    Preis belastet die Umwelt); \emph{gerechte Verteilung vs. Wachstum} (zu starke
    Nivellierung der Einkommensunterschiede mindert Leistungsanreize und kann das
    Wachstum bremsen).
  \item \textbf{Komplementäre Zielbeziehungen (Zielharmonien):}
    \emph{Wirtschaftswachstum und Vollbeschäftigung} (steigt das BIP, werden mehr
    Produktionsfaktoren benötigt --- die Arbeitslosigkeit geht zurück);
    \emph{Wirtschaftswachstum und Umweltschutz} können harmonieren, wenn boomende
    Unternehmen in umweltfreundliche Produktionsverfahren investieren.
  \item \textbf{Indifferente Zielbeziehungen:} Ziele beeinflussen sich nicht ---
    z.\,B. außenwirtschaftliches Gleichgewicht und Erhaltung der Umwelt.
\end{itemize}

\subsection{Formeln \& Rechenbeispiele}

Alle prüfungsrelevanten Berechnungen dieses Teils auf einen Blick --- jeweils mit
Formel, einfachem Zahlenbeispiel und Rechenweg. Die gleichen Zahlen finden Sie im Quiz
und auf den Folien wieder.

\subsubsection{Vom BIP zum BNE}

\begin{merksatz}[title={\bfseries Formel: BIP $\rightarrow$ BNE}]
\[
\text{BNE} = \text{BIP} - \text{Inlandseinkommen der Nichtinländer}
+ \text{Auslandseinkommen der Inländer}
\]
\end{merksatz}

\begin{beispiel}
BIP = 3.500~Mrd.~€; Nichtinländer verdienten im Inland 120~Mrd.~€; Inländer verdienten
im Ausland 180~Mrd.~€.\\[0.3em]
\textbf{Schritt 1:} Einkommen der Nichtinländer abziehen (zählt nicht zu den Inländern):
$3.500 - 120 = 3.380$~Mrd.~€.\\
\textbf{Schritt 2:} Auslandseinkommen der Inländer hinzurechnen:
$3.380 + 180 = \textbf{3.560~Mrd.~€} = \text{BNE}$.\\[0.3em]
Merke: Verdienen Inländer im Ausland \emph{mehr} als Nichtinländer im Inland, ist das
BNE größer als das BIP.
\end{beispiel}

\subsubsection{Volkseinkommen}

\begin{merksatz}[title={\bfseries Formel: Volkseinkommen}]
\[
\text{Volkseinkommen} = \text{Arbeitnehmerentgelt}
+ \text{Unternehmens- und Vermögenseinkommen}
\]
\end{merksatz}

\begin{beispiel}
Arbeitnehmerentgelt = 1.800~Mrd.~€; Unternehmens- und Vermögenseinkommen = 700~Mrd.~€.\\
\textbf{Rechnung:} $1.800 + 700 = \textbf{2.500~Mrd.~€}$ Volkseinkommen ---
die Summe aller Primäreinkommen der Inländer.
\end{beispiel}

\subsubsection{Lohnquote und Gewinnquote}

\begin{merksatz}[title={\bfseries Formeln: Lohnquote und Gewinnquote}]
\[
\text{Lohnquote in \%} = \frac{\text{Arbeitnehmerentgelt}}{\text{Volkseinkommen}} \cdot 100
\qquad
\text{Gewinnquote in \%} = 100 - \text{Lohnquote}
\]
\end{merksatz}

\begin{beispiel}
Mit den Zahlen von oben (Arbeitnehmerentgelt 1.800~Mrd.~€, Volkseinkommen
2.500~Mrd.~€):\\[0.3em]
\textbf{Schritt 1:} Lohnquote $= 1.800 \div 2.500 \cdot 100 = \textbf{72~\%}$.\\
\textbf{Schritt 2:} Gewinnquote $= 700 \div 2.500 \cdot 100 = \textbf{28~\%}$
(oder einfach $100 - 72$).\\
\textbf{Kontrolle:} $72~\% + 28~\% = 100~\%$ --- beide Quoten ergänzen sich stets zu 100~\%.
\end{beispiel}

\subsubsection{Reales vs. nominales BIP (Preisbereinigung)}

\begin{merksatz}[title={\bfseries Formel: reales BIP}]
\[
\text{reales BIP} = \frac{\text{nominales BIP}}{\text{Preisindex}} \cdot 100
\qquad (\text{Preisindex des Basisjahres} = 100)
\]
\end{merksatz}

\begin{beispiel}
Nominales BIP im Basisjahr: 2.000~Mrd.~€. Ein Jahr später: nominales BIP =
2.100~Mrd.~€, Preisindex = 105 (d.\,h. die Preise sind um 5~\% gestiegen).\\[0.3em]
\textbf{Schritt 1:} Reales BIP $= 2.100 \div 105 \cdot 100 = \textbf{2.000~Mrd.~€}$.\\
\textbf{Schritt 2:} Vergleich mit dem Vorjahr: real $2.000$ vs. $2.000$~Mrd.~€
$\Rightarrow$ \textbf{reales Wachstum = 0~\%}.\\[0.3em]
Interpretation: Der nominale Zuwachs von 5~\% bestand \emph{vollständig} aus
Preissteigerungen --- die Wirtschaft hat real nicht mehr produziert.
\end{beispiel}

\subsubsection{Verfügbares Einkommen der privaten Haushalte}

\begin{merksatz}[title={\bfseries Schema: verfügbares Einkommen}]
\begin{tabular}{ll}
  & Primäreinkommen (Faktoreinkommen)\\
$-$ & direkte Steuern\\
$-$ & Sozialversicherungsbeiträge\\
$+$ & Sozialleistungen des Staates (Transferzahlungen)\\
\midrule
$=$ & \textbf{verfügbares Einkommen} (für Konsum und Sparen)
\end{tabular}
\end{merksatz}

\begin{beispiel}
Ein Haushalt: Primäreinkommen 3.000~€, direkte Steuern 400~€,
Sozialversicherungsbeiträge 500~€, Transferzahlungen (z.\,B. Kindergeld) 600~€.\\[0.3em]
\textbf{Schritt 1:} Abzüge: $3.000 - 400 - 500 = 2.100$~€.\\
\textbf{Schritt 2:} Transfers hinzurechnen: $2.100 + 600 = \textbf{2.700~€}$
verfügbares Einkommen.
\end{beispiel}

\begin{merksatz}[title={\bfseries Wichtig für die Prüfungsvorbereitung (Teil 2)}]
Sie sollten sicher können: \emph{BIP und BNE} abgrenzen (Inlands-/Inländerkonzept), die
\emph{Entstehungs-, Verwendungs- und Verteilungsrechnung} des BIP schematisch darstellen,
\emph{nominales und reales BIP} unterscheiden, Kritikpunkte an der BIP-Berechnung
aufzählen, die \emph{Lohnquote} berechnen und das \emph{verfügbare Einkommen} ermitteln
(alle Rechenwege in Abschnitt~4), die \emph{Lorenzkurve} skizzieren und mit dem
Gini-Koeffizienten interpretieren sowie die Ziele des magischen Vierecks/Sechsecks
nennen, ihre Messbarkeit beschreiben und Zielkonflikte bzw. Zielharmonien mit Beispielen
erläutern. Üben Sie mit dem Aufgabenheft (Aufgabe~7: VGR-Rechnung) und dem Quiz zu Teil~2!
\end{merksatz}

\vspace{1em}
\hrule
\vspace{0.8em}
\noindent\textsf{\small Quellen: Rahmenplan Wirtschaftsbezogene Qualifikationen (DIHK), Abschnitte 1.1.2, 1.1.3.1;
Textband Volkswirtschaft (DIHK-Bildungs-GmbH), Kap.~2 (S.~19--27), Kap.~3.1 (S.~29--38).
Nächste Lektion (Sa, 26.09.2026): Konjunktur, Wirtschaftspolitik \& Außenwirtschaft.}
\clearpage
\section{Lektion 3: Konjunktur, Wirtschaftspolitik \& Außenwirtschaft}
\par\medskip{\bfwlabelfont\color{bfwPrimary}SESSION 1\par}\smallskip
\begin{merksatz}[title={\bfseries Worum geht es in diesem Teil?}]
Die wirtschaftliche Entwicklung verläuft in Wellen (Konjunktur). Staat und Zentralbank
versuchen, diese Schwankungen zu glätten und die Ziele des Stabilitätsgesetzes zu erreichen.
In Teil~1 der Lektion lernen Sie den \textbf{Konjunkturzyklus} mit seinen Indikatoren, die
\textbf{Geldpolitik der EZB} (Ziel, Instrumente, Transmissionskette) und die
\textbf{Finanzpolitik} (antizyklische Fiskalpolitik, automatische Stabilisatoren,
Staatsverschuldung) kennen. In Teil~2 (Session~2, ab 10:55~Uhr) folgen die weiteren
Politikfelder, die wirtschaftspolitischen Konzeptionen sowie Außenwirtschaft und EU ---
alle Themen sind regelmäßig Gegenstand der IHK-Prüfung.
\end{merksatz}



\subsection{Konjunktur: Zyklus und Indikatoren (Kurzüberblick)}

Jede volkswirtschaftliche Entwicklung ist Schwankungen unterworfen. Man unterscheidet
\textbf{Konjunkturschwankungen} (mehrjährige, wellenförmige Schwankungen der
gesamtwirtschaftlichen Aktivität um den langfristigen Wachstumstrend) von
\textbf{Saisonschwankungen} (regelmäßige, jahreszeitlich bedingte Schwankungen, z.\,B.
Bau im Winter, Tourismus im Sommer) und zufallsbedingten Schwankungen (z.\,B.
Naturkatastrophen, Krisen).

\subsubsection{Die vier Phasen des Konjunkturzyklus}

\begin{enumerate}
  \item \textbf{Aufschwung (Expansion):} steigende Nachfrage, Produktion und
    Investitionen, sinkende Arbeitslosigkeit, wachsender Optimismus.
  \item \textbf{Hochkonjunktur (Boom):} volle Kapazitätsauslastung, hohe Gewinne,
    Vollbeschäftigung, aber steigende Preise und Zinsen (Überhitzungsgefahr).
  \item \textbf{Abschwung (Rezession):} sinkende Nachfrage und Produktion,
    Auftragsrückgänge, zunehmende Arbeitslosigkeit, Pessimismus.
  \item \textbf{Tiefphase (Depression):} geringe Auslastung, hohe Arbeitslosigkeit,
    niedrige Preise/Zinsen --- Ausgangspunkt für einen neuen Aufschwung.
\end{enumerate}

\begin{center}
\begin{tikzpicture}[font=\small, line cap=round]
  % Achsen
  \draw[-{Stealth[length=2.5mm]}, thick, bfwInk] (0,0) -- (0,5.3)
    node[anchor=west, font=\footnotesize\bfseries] {BIP / wirtschaftliche Aktivität};
  \draw[-{Stealth[length=2.5mm]}, thick, bfwInk] (0,0) -- (13.4,0)
    node[anchor=north east, font=\footnotesize\bfseries, yshift=-2pt] {Zeit};
  % Hilfslinien: Boom-zu-Boom-Zyklus
  \draw[dotted, thick, bfwInkSoft!70] (3.0,3.35) -- (3.0,-0.55);
  \draw[dotted, thick, bfwInkSoft!70] (9.2,4.25) -- (9.2,-0.55);
  \draw[{Stealth[length=2mm]}-{Stealth[length=2mm]}, thick, bfwPrimary] (3.0,-0.55) -- (9.2,-0.55)
    node[midway, below=2pt, font=\scriptsize, text=bfwInkSoft]
    {ein vollständiger Konjunkturzyklus (z.\,B. von Boom zu Boom gemessen)};
  % Wachstumstrend (gestrichelt, leicht steigend)
  \draw[dashed, thick, bfwInkSoft] (0,1.7) -- (12.9,3.5)
    node[anchor=south east, font=\footnotesize\bfseries, text=bfwInkSoft] {Trend};
  % Konjunkturwelle um den Trend
  \draw[line width=1.4pt, bfwPrimaryDark]
    plot[smooth, tension=0.75] coordinates
    {(0,1.0) (1.5,1.5) (3.0,3.35) (4.5,2.3) (6.0,1.1) (7.6,2.5) (9.2,4.25) (10.8,3.1) (12.3,2.1)};
  % Markierung Hoch- und Tiefpunkt
  \draw[thick, bfwAccent, fill=bfwBgAmber] (3.0,3.35) circle (1.5mm);
  \draw[thick, bfwRot, fill=bfwBgRose] (6.0,1.1) circle (1.5mm);
  % Phasenbeschriftung
  \node[text=bfwSuccess!75!black, font=\small\bfseries, align=center] at (1.3,2.8)
    {Aufschwung\\{\scriptsize\color{bfwInkSoft}(Expansion)}};
  \node[text=bfwAccent!85!black, font=\small\bfseries, align=center] at (3.0,4.25)
    {Boom\\{\scriptsize\color{bfwInkSoft}(Hochkonjunktur)}};
  \node[text=bfwWarn, font=\small\bfseries, align=center] at (5.05,3.1)
    {Abschwung\\{\scriptsize\color{bfwInkSoft}(Rezession)}};
  \node[text=bfwRot, font=\small\bfseries, align=center] at (6.0,0.45)
    {Tief\\{\scriptsize\color{bfwInkSoft}(Depression)}};
  \node[text=bfwSuccess!75!black, font=\footnotesize\bfseries, align=center, anchor=east] at (8.2,3.6)
    {neuer\\Aufschwung};
\end{tikzpicture}
\end{center}

\emph{Faustregel:} Sinkt das reale BIP in \textbf{zwei aufeinanderfolgenden Quartalen},
spricht man von einer \emph{(technischen) Rezession}.

\subsubsection{Zyklusdauern: kurze, mittlere und lange Wellen}

Nach der Dauer der Schwankungen unterscheidet die Konjunkturforschung drei Wellenarten,
die sich überlagern:

\begin{itemize}
  \item \textbf{Kitchin-Zyklen} (kurze Wellen, ca. 3--4 Jahre): Lagerzyklen ---
    Unternehmen bauen Lagerbestände auf und ab;
  \item \textbf{Juglar-Zyklen} (mittlere Wellen, ca. 7--11 Jahre): der „klassische``
    Konjunkturzyklus dieses Kapitels, getragen von den Investitionen der Unternehmen;
  \item \textbf{Kondratjew-Zyklen} (lange Wellen, ca. 40--60 Jahre): ausgelöst durch
    \emph{Basisinnovationen}, die ganze Investitionswellen tragen (Dampfmaschine,
    Eisenbahn, Elektrotechnik/Chemie, Automobil, Informationstechnik).
\end{itemize}

\subsubsection{Konjunkturindikatoren}

\begin{itemize}
  \item \textbf{Frühindikatoren} (laufen der Entwicklung voraus): Auftragseingänge,
    Geschäftsklimaindex (ifo), Baugenehmigungen, Aktienkurse.
  \item \textbf{Präsenzindikatoren} (zeigen die aktuelle Lage): reales BIP,
    Kapazitätsauslastung, Industrieproduktion, Kurzarbeit, offene Stellen.
  \item \textbf{Spätindikatoren} (laufen hinterher): Arbeitslosenquote, Preisniveau
    (Inflationsrate), Insolvenzen, Lohnentwicklung.
\end{itemize}

\begin{eselsbruecke}
\textbf{F-P-S wie „Fahrplan der Konjunktur``:} \textbf{F}rüh kündigt an,
\textbf{P}räsenz zeigt an, \textbf{S}pät hinkt nach. Die Arbeitslosenquote ist ein
\emph{Spät}indikator --- Unternehmen entlassen erst, wenn der Abschwung schon läuft.
\end{eselsbruecke}

\subsection{Geldpolitik der Europäischen Zentralbank (EZB)}

Die Geldpolitik der Eurozone liegt seit 1999 in den Händen der \textbf{Europäischen
Zentralbank (EZB)}. Sie steuert die Geldmengenentwicklung in der Eurozone und unterstützt
zudem die allgemeine Wirtschaftspolitik in der Gemeinschaft. \textbf{Vorrangiges Ziel}
ist nach dem Vertrag von Maastricht die \textbf{Preisniveaustabilität}: Der Anstieg des
\emph{harmonisierten Verbraucherpreisindex (HVPI)} soll mittelfristig \textbf{zwei
Prozent} betragen (seit der Strategieüberprüfung 2021 ein \emph{symmetrisches}
Inflationsziel; zuvor: „unter, aber nahe zwei Prozent``) --- damit sind sowohl Inflation
als auch Deflation ausgeschlossen. Die EZB
berücksichtigt bei ihrer Geldmengensteuerung sowohl die Preisentwicklung als auch die
Geldmengenentwicklung (\textbf{Zwei-Säulen-Strategie}; seit der Strategieüberprüfung 2021
weiterentwickelt zum integrierten Analyserahmen: wirtschaftliche Analyse + monetäre und
finanzielle Analyse).

\begin{merksatz}
Geschäftsbanken benötigen Zentralbankgeld, um Kredite vergeben zu können. Die Beschaffung
dieses Geldes heißt \textbf{Refinanzierung} --- wichtigste Refinanzierungsquelle ist die
Zentralbank. Genau hier setzen die geldpolitischen Instrumente an: Sie steuern
\emph{Liquidität} (Kreditmenge) und \emph{Zinsen} (Kreditkosten).
\end{merksatz}

\subsubsection{Instrument 1: Offenmarktgeschäfte}

Die Offenmarktgeschäfte stehen im Zentrum des Instrumentariums. Die EZB tritt als Käufer
oder Verkäufer von Wertpapieren am offenen Markt auf: \emph{Kauft} sie Wertpapiere, fließt
Zentralbankgeld in den Bankensektor (Geldschöpfung, Nachfrage wird belebt); \emph{verkauft}
sie, fließt Geld zurück (Geldvernichtung, Zinsen steigen, Nachfrage sinkt). Vier Kategorien:

\begin{itemize}
  \item \textbf{Hauptrefinanzierungsgeschäfte (Haupttender):} befristete Transaktionen im
    \emph{wöchentlichen} Rhythmus mit einer Laufzeit von \emph{einer Woche}
    (Pensionsgeschäfte bzw. Kredite gegen Verpfändung von Wertpapieren). Der hierfür
    festgelegte Zinssatz ist der \textbf{wichtigste Leitzins} im Euroraum.
  \item \textbf{Längerfristige Refinanzierungsgeschäfte (Basistender):} gleiche Technik,
    ursprünglich monatlich mit \emph{drei Monaten} Laufzeit (heute auch länger).
  \item \textbf{Feinsteuerungsoperationen:} gleichen kurzfristige Schwankungen der
    Bankenliquidität aus.
  \item \textbf{Strukturelle Operationen:} beeinflussen den Bedarf an Zentralbankgeld
    langfristig (z.\,B. Emission von Schuldverschreibungen).
\end{itemize}

Abgewickelt werden die Geschäfte in \textbf{Tenderverfahren} (standardisierte
Ausschreibungen): nach der Durchführung \emph{Standardtender} (24 Stunden, für Haupt- und
Basistender) oder \emph{Schnelltender} (wenige Stunden, für Feinsteuerung); nach der
Zuteilung \emph{Mengentender} (die EZB gibt den Zinssatz vor, die Banken bieten nur die
gewünschten Beträge; bei Überzeichnung wird anteilig zugeteilt --- \emph{Repartierung})
oder \emph{Zinstender} (die Banken bieten Zinssatz \emph{und} Betrag; die Zuteilung
richtet sich nach der Höhe der Zinsgebote).

\begin{beispiel}
\textbf{Mengentender mit Repartierung (durchgerechnet):} Die EZB schreibt einen
Haupttender zum Festzinssatz aus und will 50~Mrd.~Euro zuteilen. Die Geschäftsbanken
bieten zusammen 200~Mrd.~Euro (Überzeichnung).\\[0.3em]
\textbf{Schritt 1:} Repartierungsquote = Zuteilungsbetrag $\div$ Gebotssumme
$= 50 \div 200 = 0{,}25 = \textbf{25~\%}$.\\
\textbf{Schritt 2:} Eine Bank, die 8~Mrd.~Euro geboten hat, erhält
$8 \times 0{,}25 = \textbf{2~Mrd.~Euro}$ Zentralbankgeld.\\[0.3em]
\emph{Merke:} Beim Mengentender wird bei Überzeichnung \emph{anteilig} zugeteilt ---
beim Zinstender entscheidet dagegen die Höhe des gebotenen Zinssatzes.
\end{beispiel}

\subsubsection{Instrument 2: Ständige Fazilitäten}

Die Fazilitäten stehen den Geschäftsbanken \emph{ständig} zur Verfügung (keine
Ausschreibung, Laufzeit: ein Geschäftstag):

\begin{itemize}
  \item \textbf{Spitzenrefinanzierungsfazilität} (Kreditfazilität): Banken können sich
    über Nacht unbegrenzt Zentralbankgeld leihen --- ihr Zinssatz bildet die
    \emph{Obergrenze} des Geldmarktzinses.
  \item \textbf{Einlagefazilität:} Banken legen überschüssige Liquidität bis zum nächsten
    Tag bei der EZB an --- ihr Zinssatz bildet die \emph{Untergrenze}.
\end{itemize}

Beide Sätze bilden zusammen den \textbf{Zinskorridor}, in dessen Mitte sich der
Hauptrefinanzierungssatz bewegt. Diese \textbf{drei Leitzinsen} (Einlagefazilität,
Hauptrefinanzierungssatz, Spitzenrefinanzierungsfazilität) müssen Sie für die Prüfung
nennen können.

\begin{center}
\begin{tikzpicture}[font=\footnotesize, line cap=round]
  % Korridor-Band
  \fill[bfwBgTeal] (0,0) rectangle (11.8,2.6);
  % Obergrenze
  \draw[line width=1.2pt, bfwWarn] (0,2.6) -- (11.8,2.6);
  \node[anchor=south west, text=bfwWarn, font=\footnotesize\bfseries] at (0,2.68)
    {Spitzenrefinanzierungsfazilität = Obergrenze (Banken leihen über Nacht)};
  % Mitte
  \draw[line width=1.2pt, bfwPrimaryDark] (0,1.3) -- (11.8,1.3);
  \node[anchor=south west, text=bfwPrimaryDark, font=\footnotesize\bfseries] at (0,1.38)
    {Hauptrefinanzierungssatz = wichtigster Leitzins (Haupttender, 1 Woche)};
  % Untergrenze
  \draw[line width=1.2pt, bfwAccent!85!black] (0,0) -- (11.8,0);
  \node[anchor=north west, text=bfwAccent!85!black, font=\footnotesize\bfseries] at (0,-0.1)
    {Einlagefazilität = Untergrenze (Banken legen über Nacht an)};
  % Geldmarktzins als Schlangenlinie im Korridor
  \draw[thick, bfwInkSoft, decorate, decoration={snake, amplitude=1mm, segment length=10mm}]
    (0.4,0.72) -- (8.4,0.72);
  \node[anchor=west, text=bfwInkSoft, font=\scriptsize] at (8.55,0.72) {Geldmarktzins};
  % Doppelpfeil Zinskorridor
  \draw[{Stealth[length=2mm]}-{Stealth[length=2mm]}, thick, bfwPrimaryDark]
    (11.35,0.08) -- (11.35,2.52);
  \node[anchor=south, rotate=90, text=bfwPrimaryDark, font=\scriptsize\bfseries] at (11.2,1.3)
    {Zinskorridor};
\end{tikzpicture}
\end{center}

\subsubsection{Instrument 3: Mindestreservepolitik}

Die EZB verpflichtet die Geschäftsbanken, bestimmte prozentuale Anteile der bei ihnen in
Sicht-, Termin- und Spareinlagen angelegten Gelder als \textbf{Pflichtguthaben} auf Konten
bei der Zentralbank zu unterhalten. Eine \emph{Erhöhung} des Mindestreservesatzes entzieht
den Banken Liquidität (weniger Kreditvergabe möglich), eine \emph{Senkung} wirkt umgekehrt.
Über Änderungen entscheidet der EZB-Rat.

\subsubsection{Expansive und restriktive Geldpolitik --- die Transmissionskette}

\begin{itemize}
  \item \textbf{Expansive Geldpolitik} („Politik des billigen Geldes``): Leitzinsen senken,
    Liquidität zuführen $\Rightarrow$ Kreditzinsen sinken $\Rightarrow$ Kreditnachfrage und
    Geldmenge steigen $\Rightarrow$ Investitionen und Konsum steigen $\Rightarrow$ BIP
    steigt, Arbeitslosigkeit sinkt --- Mittel gegen Rezession und Deflation.
  \item \textbf{Restriktive Geldpolitik}: Leitzinsen erhöhen, Liquidität verknappen
    (z.\,B. Mindestreserve erhöhen, Zuteilungsvolumen verringern) $\Rightarrow$ Zinsniveau
    steigt, Sparneigung nimmt zu $\Rightarrow$ kreditfinanzierte Investitionen und
    Konsumentenkredite gehen zurück $\Rightarrow$ gesamtwirtschaftliche Nachfrage sinkt
    $\Rightarrow$ Preisanstieg wird verringert, Konjunktur gedämpft --- Mittel gegen
    Inflation.
\end{itemize}

Die folgende Wirkungskette (\textbf{Transmissionsmechanismus}) zeigt die
\emph{restriktive} Variante --- bei expansiver Geldpolitik drehen sich alle Pfeile um
(Zinssenkung $\Rightarrow$ billigere Kredite $\Rightarrow$ mehr Investitionen und Konsum
$\Rightarrow$ Nachfrage und Beschäftigung steigen):

\begin{center}
\begin{tikzpicture}[font=\small,
  kette/.style 2 args={draw=#1, fill=#2, thick, rounded corners=4pt, align=center,
    text width=3.9cm, minimum height=1.25cm, inner sep=5pt},
  pfeil/.style={-{Stealth[length=3mm]}, line width=1.2pt, bfwPrimaryDark}]
  % Obere Reihe
  \node[kette={bfwRot}{bfwBgRose}] (s1) at (0,0)
    {\textbf{EZB erhöht}\\{\footnotesize die Leitzinsen}};
  \node[kette={bfwAccent!85!black}{bfwBgAmber}] (s2) at (5.7,0)
    {\textbf{Refinanzierung}\\{\footnotesize der Banken wird teurer}};
  \node[kette={bfwAccent!85!black}{bfwBgAmber}] (s3) at (11.4,0)
    {\textbf{Kreditzinsen steigen}\\{\footnotesize Sparneigung nimmt zu}};
  % Untere Reihe (von rechts nach links)
  \node[kette={bfwPrimaryDark}{bfwBgTeal}] (s4) at (11.4,-2.2)
    {\textbf{Kreditnachfrage sinkt}\\{\footnotesize Unternehmen und Haushalte}};
  \node[kette={bfwPrimaryDark}{bfwBgTeal}] (s5) at (5.7,-2.2)
    {\textbf{Investitionen und Konsum}\\{\footnotesize gehen zurück}};
  \node[kette={bfwPrimaryDark}{bfwBgTeal}] (s6) at (0,-2.2)
    {\textbf{Gesamtnachfrage sinkt}\\{\footnotesize Konjunktur wird gedämpft}};
  % Zielbox
  \node[draw=bfwSuccess!65!black, fill=bfwBgGruen, thick, rounded corners=4pt, align=center,
        text width=14.6cm, inner sep=6pt] (ziel) at (5.7,-4.2)
    {\textbf{\color{bfwSuccess!55!black}Ziel erreicht: Der Preisauftrieb lässt nach ---
      Inflation wird bekämpft}\\
     {\footnotesize\color{bfwInkSoft}Nebenwirkung: Wachstum und Beschäftigung können
      leiden (restriktive Geldpolitik)}};
  % Pfeile
  \draw[pfeil] (s1) -- (s2);
  \draw[pfeil] (s2) -- (s3);
  \draw[pfeil] (s3) -- (s4);
  \draw[pfeil] (s4) -- (s5);
  \draw[pfeil] (s5) -- (s6);
  \draw[pfeil] (s6.south) -- (s6.south |- ziel.north);
\end{tikzpicture}
\end{center}

\begin{beispiel}
Die Inflationsrate liegt deutlich über zwei Prozent. Die EZB \emph{erhöht die Leitzinsen}:
Die Refinanzierung wird für die Geschäftsbanken teurer, sie geben die höheren Zinsen an
Unternehmen und Haushalte weiter. Kredite für Maschinen, Bauvorhaben und Konsum lohnen sich
weniger, die gesamtwirtschaftliche Nachfrage sinkt --- der Preisauftrieb lässt nach.
Kehrseite: Wachstum und Beschäftigung können leiden.
\end{beispiel}

\begin{eselsbruecke}
\textbf{„Zins runter --- Wirtschaft munter. Zins rauf --- Preisbremse drauf.``}
Expansiv = Gas geben (gegen Rezession), restriktiv = bremsen (gegen Inflation).
\end{eselsbruecke}

\subsubsection{Grenzen der Geldpolitik}

Die Instrumente wirken nur bedingt: Internationale Finanzströme und hohe Lohnabschlüsse
(über dem Produktivitätsfortschritt) können Inflation erzeugen; ob und wann Banken
Zinsänderungen weitergeben, ist unklar; bei pessimistischer Grundhaltung fördern selbst
niedrige Zinsen die Investitions- und Konsumbereitschaft nicht; und die EZB ist auf eine
einvernehmliche Fiskalpolitik der Regierungen angewiesen (erhöht der Staat in der
Inflationsbekämpfung gleichzeitig schuldenfinanziert seine Ausgaben, verpufft die
restriktive Wirkung).

\begin{achtung}
Klassische Verwechslung in der Prüfung: \textbf{Geldpolitik} betreibt die \emph{EZB}
(Leitzinsen, Offenmarktgeschäfte, Mindestreserve), \textbf{Fiskalpolitik} betreibt der
\emph{Staat} (Steuern, Staatsausgaben). Formulierungen wie „der Staat senkt die
Leitzinsen`` sind immer falsch --- die EZB ist von Weisungen der Regierungen
\emph{unabhängig}.
\end{achtung}

\subsection{Finanzpolitik (Fiskalpolitik)}

Die Finanzpolitik umfasst alle Maßnahmen, die den \textbf{Staatshaushalt} (Einnahmen und
Ausgaben von Bund, Ländern, Gemeinden und Sozialversicherung) betreffen. Aufgaben:

\begin{itemize}
  \item \textbf{Beschaffung von Einnahmen} (v.\,a. Steuern, daneben Beiträge, Gebühren) zur
    Finanzierung öffentlicher Ausgaben; reichen sie nicht, steigt die Staatsverschuldung;
  \item \textbf{Umverteilung von Einkommen (Distribution)}, z.\,B. progressive
    Einkommensteuer, Transferleistungen, Sozialversicherung;
  \item \textbf{Stabilitätsaufgaben:} Steuerung der Konjunktur im Sinne des
    \textbf{Stabilitätsgesetzes (StabG, 1967)}.
\end{itemize}

\subsubsection{Antizyklische Fiskalpolitik nach Keynes}

Nach der Lehre von \textbf{John Maynard Keynes} sollen staatliche Einnahmen- und
Ausgabenpolitik dem Konjunkturverlauf \emph{entgegenwirken}:

\begin{center}
\begin{tabular}{p{3.2cm} p{5.2cm} p{5.2cm}}
\toprule
 & \textbf{Boom (restriktiv)} & \textbf{Rezession (expansiv)}\\
\midrule
Ausgaben & senken (Investitionen, Subventionen, Transfers zurückfahren) & erhöhen (öffentliche Aufträge, Konjunkturprogramme)\\
Einnahmen & erhöhen (Steuern) & senken (Steuererleichterungen)\\
Budgetwirkung & Überschüsse bilden (Rücklagen) & Defizite zulassen (\emph{Deficit Spending})\\
\bottomrule
\end{tabular}
\end{center}

Nimmt der Staat in der Rezession bewusst neue Schulden auf, um zusätzliche Nachfrage zu
erzeugen, spricht man von \textbf{Deficit Spending}.

\begin{beispiel}
In der Finanz- und Wirtschaftskrise wurden die \emph{Konjunkturpakete I und II} (2008/2009)
verabschiedet: Investitionen in Infrastruktur und öffentliche Gebäude,
Abschreibungserleichterungen (wirkten wie eine Steuererleichterung) und die
\emph{Umweltprämie (Abwrackprämie)}, die die Nachfrage nach Kraftfahrzeugen fördern sollte.
\end{beispiel}

\begin{achtung}
Grenzen der antizyklischen Fiskalpolitik: \textbf{Timelags} (bis Konjunkturprogramme
wirken, ist die Rezession oft vorbei --- die Maßnahme wirkt dann prozyklisch),
\textbf{Crowding-out} (kreditfinanzierte Staatsnachfrage verdrängt private Investitionen
über steigende Zinsen) und die \emph{politische Asymmetrie}: Schuldenmachen in der
Rezession ist populär, Sparen und Tilgen im Boom nicht --- die Staatsverschuldung steigt
deshalb tendenziell dauerhaft.
\end{achtung}

\subsubsection{Automatische Stabilisatoren}

Neben gezielten Maßnahmen wirken \textbf{automatische (eingebaute) Stabilisatoren} ohne
politische Entscheidung antizyklisch: Die \emph{progressive Einkommensteuer} entzieht im
Boom überproportional Kaufkraft und entlastet in der Rezession; \emph{Arbeitslosengeld und
Sozialleistungen} stützen in der Rezession die Nachfrage, während im Aufschwung die
Beiträge die Kaufkraft dämpfen. Der Haushalt „atmet`` also mit der Konjunktur.

\subsubsection{Staatsverschuldung und ihre Grenzen}

Zu unterscheiden sind die jährliche \textbf{Neuverschuldung} (Kredite zum Ausgleich eines
Haushaltsdefizits) und die \textbf{Gesamtverschuldung} des Staates. Hohe Schulden und
Zinsverpflichtungen beschränken die Handlungsfähigkeit künftiger Politik. Juristische
Grenzen: die \textbf{Schuldenbremse} im Grundgesetz (Art.\,109 GG, seit 2009) und der
\emph{europäische Fiskalpakt} (2013), der die Mitgliedsländer verpflichtet, die
Schuldenbremse national zu regeln. Die Schuldenbremse begrenzt die \emph{strukturelle}
(konjunkturunabhängige) Neuverschuldung des Bundes auf \textbf{0,35~\%} des BIP; seit der
Grundgesetzänderung 2025 dürfen auch die Länder \emph{zusammen} strukturell Kredite von
bis zu 0,35~\% des BIP aufnehmen (vorher: gar keine). Ausnahmen gelten für
Naturkatastrophen und außergewöhnliche Notsituationen (z.\,B. Corona-Pandemie 2020);
seit 2025 bleiben zudem Verteidigungsausgaben oberhalb von 1~\% des BIP außen vor, und für
Infrastruktur wurde ein kreditfinanziertes Sondervermögen eingerichtet.

\begin{beispiel}
Das BIP beträgt 4.000~Mrd.~Euro. Zulässige strukturelle Neuverschuldung des Bundes:
$4.000 \times 0{,}0035 = \textbf{14~Mrd.~Euro}$. Plant der Bund 20~Mrd.~Euro neue
Schulden, liegt er 6~Mrd.~Euro über der Grenze --- zulässig nur bei einer vom Bundestag
festgestellten Notsituation.
\end{beispiel}

\subsection{Formeln \& Rechenbeispiele}

Die folgenden Berechnungen sind Prüfungsklassiker --- die ersten beiden folgen demselben
Muster: \emph{(neu $-$ alt) durch alt mal 100}. Prägen Sie sich Formel und Rechenweg
anhand der einfachen Zahlenbeispiele ein.

\begin{merksatz}[title={\bfseries Formel 1: Wirtschaftswachstum (Veränderungsrate des realen BIP)}]
\[
\text{Wachstumsrate in \%} =
\frac{\text{BIP}_{\text{neu}} - \text{BIP}_{\text{alt}}}{\text{BIP}_{\text{alt}}} \times 100
\]
Immer mit dem \emph{realen} BIP rechnen (Preise eines Basisjahres) --- sonst misst man
Inflation statt echtem Wachstum.
\end{merksatz}

\begin{beispiel}
Das reale BIP steigt von 3.000~Mrd.~Euro (Vorjahr) auf 3.060~Mrd.~Euro.\\[0.3em]
\textbf{Schritt 1:} Veränderung berechnen: $3.060 - 3.000 = 60$~Mrd.~Euro.\\
\textbf{Schritt 2:} Durch den \emph{alten} Wert teilen: $60 \div 3.000 = 0{,}02$.\\
\textbf{Schritt 3:} Mal 100: $0{,}02 \times 100 = \textbf{2{,}0~\%}$ Wirtschaftswachstum.
\end{beispiel}

\begin{merksatz}[title={\bfseries Formel 2: Inflationsrate (Veränderungsrate des Preisindex)}]
\[
\text{Inflationsrate in \%} =
\frac{\text{Preisindex}_{\text{neu}} - \text{Preisindex}_{\text{alt}}}{\text{Preisindex}_{\text{alt}}} \times 100
\]
Grundlage ist der Verbraucherpreisindex (in der Eurozone: HVPI), der die Preisentwicklung
eines repräsentativen Warenkorbs misst.
\end{merksatz}

\begin{beispiel}
Der Verbraucherpreisindex steigt von 110 auf 115,5 Punkte.\\[0.3em]
\textbf{Schritt 1:} Veränderung: $115{,}5 - 110 = 5{,}5$ Punkte.\\
\textbf{Schritt 2:} Durch den \emph{alten} Indexstand teilen: $5{,}5 \div 110 = 0{,}05$.\\
\textbf{Schritt 3:} Mal 100: $0{,}05 \times 100 = \textbf{5{,}0~\%}$ Inflationsrate.\\[0.3em]
\emph{Typischer Fehler:} durch 100 statt durch den Basiswert 110 teilen --- das ergäbe
fälschlich 5,5~\%.
\end{beispiel}

\begin{merksatz}[title={\bfseries Formel 3: Deflationierung --- vom nominalen zum realen BIP}]
\[
\text{reales BIP} =
\frac{\text{nominales BIP}}{\text{Preisindex}} \times 100
\]
Das nominale BIP wird mit dem Preisindex (Basisjahr $=100$) „preisbereinigt``. Erst das
reale BIP zeigt, ob tatsächlich mehr Güter produziert wurden.
\end{merksatz}

\begin{beispiel}
Vorjahr: BIP 3.000~Mrd.~Euro (Preisindex 100). Aktuelles Jahr: nominales BIP
3.150~Mrd.~Euro (${}+5~\%$), Preisindex 105.\\[0.3em]
\textbf{Schritt 1:} Reales BIP $= 3.150 \div 105 \times 100 = 3.000$~Mrd.~Euro.\\
\textbf{Schritt 2:} Reale Wachstumsrate $= (3.000 - 3.000) \div 3.000 \times 100
= \textbf{0{,}0~\%}$.\\[0.3em]
\emph{Interpretation:} Das gesamte nominale „Wachstum`` von 5~\% ist reine Inflation ---
real wurde nicht mehr produziert. Deshalb immer mit dem \emph{realen} BIP rechnen.
\end{beispiel}

\begin{merksatz}[title={\bfseries Wichtig für die Prüfungsvorbereitung (Teil 1)}]
Sie sollten sicher können: den \emph{Konjunkturzyklus} mit Indikatoren (F-P-S), das Ziel
und die \emph{drei Instrumente der EZB} samt der drei \emph{Leitzinsen} (Zinskorridor!)
und der Wirkungskette einer Leitzinsänderung (Transmission), die Unterscheidung
\emph{expansiv vs. restriktiv}, die \emph{antizyklische Fiskalpolitik} (Boom vs.
Rezession) mit automatischen Stabilisatoren und Schuldenbremse (0,35~\%) sowie die
\emph{Formeln} für Wachstumsrate, Inflationsrate und Deflationierung (Abschnitt~4). In Teil~2
folgen die weiteren Politikfelder, die Konzeptionen sowie Außenwirtschaft und EU ---
inklusive Wiederholung des gesamten Blocks~1.1.
\end{merksatz}

\vspace{1em}
\hrule
\vspace{0.8em}
\noindent\textsf{\small Quellen: Rahmenplan Wirtschaftsbezogene Qualifikationen (DIHK), Abschnitt 1.1.3.2;
Textband Volkswirtschaft (DIHK-Bildungs-GmbH), Kap.~3.2, S.~39--52. Weiter geht es in Session~2 (10:55--13:15~Uhr):
Teil~2 --- Politikfelder, Außenwirtschaft \& EU + Wiederholung Block~1.1.}
\par\medskip{\bfwlabelfont\color{bfwPrimary}SESSION 2\par}\smallskip
\begin{merksatz}[title={\bfseries Worum geht es in diesem Teil?}]
Nach Geld- und Finanzpolitik (Teil~1) lernen Sie in Teil~2 die weiteren
\textbf{wirtschaftspolitischen Maßnahmen} kennen: Wachstums-, Tarif-, Arbeitsmarkt- und
Umweltpolitik --- sowie die beiden großen \textbf{Konzeptionen} (nachfrage- vs.
angebotsorientierte Wirtschaftspolitik). Anschließend geht es um die
\textbf{Außenwirtschaft}: Freihandel und Protektionismus, Zahlungsbilanz und
Wechselkurse sowie die Besonderheiten der EU (Binnenmarkt, Wirtschafts- und
Währungsunion). Damit schließen wir den Themenblock~1.1
des Rahmenplans ab und wiederholen den gesamten Prüfungsstoff der Lektionen~1--3 ---
alle Themen sind regelmäßig Gegenstand der IHK-Prüfung.
\end{merksatz}



\subsection{Wachstumspolitik}

Wachstumspolitik umfasst alle staatlichen Maßnahmen für ein \emph{angemessenes und
stetiges} Wirtschaftswachstum. Zu unterscheiden sind \textbf{quantitatives Wachstum}
(reine Zunahme des realen BIP bzw. Pro-Kopf-Einkommens) und \textbf{qualitatives Wachstum}
(zusätzlich Schonung der Umwelt und Verbesserung der Lebensqualität; Messkonzepte: soziale
Indikatoren der OECD, Net Economic Welfare). Ansatzpunkte der Wachstumspolitik:

\begin{itemize}
  \item Förderung von \textbf{Investitionen} (besseres Investitionsklima, Gewinnerwartung),
  \item Beschleunigung des technischen Fortschritts durch Förderung von
    \textbf{Forschung und Entwicklung},
  \item Förderung von \textbf{Bildung}, um das Angebot an qualifizierten Arbeitskräften zu
    erhöhen,
  \item im weiteren Sinne auch die \textbf{Wettbewerbspolitik} (funktionierender
    Leistungswettbewerb) und die wirtschaftliche Infrastruktur.
\end{itemize}

\begin{merksatz}
\textbf{Quantitativ} = „mehr`` (Zunahme des realen BIP), \textbf{qualitativ} = „besser``
(Wachstum unter Schonung von Umwelt und Ressourcen, mehr Lebensqualität). Maßstab für den
Wohlstandsvergleich ist das \emph{BIP pro Kopf}: Wächst die Bevölkerung schneller als das
BIP, sinkt der materielle Wohlstand \emph{trotz} Wachstums.
\end{merksatz}

\begin{beispiel}
Das reale BIP eines Landes steigt von 500 auf 515~Mrd.~Euro (${}+3{,}0~\%$), die
Bevölkerung wächst von 10,0 auf 10,4~Mio. (${}+4{,}0~\%$).\\[0.3em]
\textbf{Schritt 1:} BIP pro Kopf alt: $500\,\text{Mrd.} \div 10{,}0\,\text{Mio.}
= 50.000$~Euro.\\
\textbf{Schritt 2:} BIP pro Kopf neu: $515\,\text{Mrd.} \div 10{,}4\,\text{Mio.}
\approx 49.519$~Euro.\\
\textbf{Schritt 3:} Veränderung: $(49.519 - 50.000) \div 50.000 \times 100 \approx
\textbf{$-$1{,}0~\%}$.\\[0.3em]
\emph{Interpretation:} Trotz 3~\% Wachstum sinkt der Wohlstand je Einwohner --- das BIP
pro Kopf ist der bessere Wohlstandsindikator als das absolute BIP.
\end{beispiel}

\subsection{Tarifpolitik}

Die Tarifpolitik umfasst alle Fragen rund um Verhandlung und Abschluss von
\textbf{Tarifverträgen} zwischen den Tarifparteien --- in Deutschland in erster Linie
\emph{Arbeitgeberverbände} und \emph{Gewerkschaften}; auch ein einzelner Arbeitgeber kann
Tarifvertragspartei sein (\emph{Haustarifvertrag}).

\begin{merksatz}[title={\bfseries Tarifautonomie (Art. 9 Abs. 3 GG)}]
Die Tarifpartner sind von Weisungen des Staates \textbf{unabhängig}. Der Abschluss von
Tarifverträgen liegt ausschließlich in den Händen der Tarifparteien; der Staat ist zur
Neutralität verpflichtet. Staatliche Vorgaben sind nur als \emph{Gesetze} zulässig
(z.\,B. Bundesurlaubsgesetz, Entgeltfortzahlungsgesetz, Arbeitszeitgesetz,
Mindestlohngesetz) --- sie setzen Mindestansprüche, die nicht unterschritten werden dürfen.
\end{merksatz}

\subsubsection{Tarifvertragsarten}

\begin{itemize}
  \item \textbf{Lohn- und Gehalts(rahmen)tarifverträge:} regeln die Entgeltfindung ---
    Höhe des Entgelts, Eingruppierung, Zuschläge, Zulagen;
  \item \textbf{Manteltarifverträge:} regeln die arbeitsrechtlichen Rahmenbedingungen ---
    Verteilung der Wochenarbeitszeit, Pausen, Urlaub, Schichtarbeit, Freistellungen,
    Altersversorgung (meist mit langer Laufzeit).
\end{itemize}

\subsubsection{Lohnpolitik: Kaufkrafttheorie vs. produktivitätsorientierte Lohnpolitik}

\begin{itemize}
  \item Die \textbf{Gewerkschaften} argumentieren \emph{nachfrageorientiert}
    (\textbf{Kaufkrafttheorie}): Höhere Löhne steigern die Nachfrage, regen Produktion und
    Investitionen an und verbessern die Beschäftigung.
  \item Die \textbf{Arbeitgeber} argumentieren \emph{angebotsorientiert}: Löhne sind ein
    Kostenfaktor, der die (internationale) Wettbewerbsfähigkeit gefährdet. Sie fordern eine
    \textbf{produktivitätsorientierte Lohnpolitik}: Lohnsteigerungen nicht höher als der
    Produktivitätsfortschritt --- das verhindert Lohnkosteninflation und lohnkostenbedingte
    Arbeitslosigkeit.
\end{itemize}

Lohnerhöhungen wirken doppelt: Sie erhöhen die Produktionskosten
(\emph{Kostenerhöhungseffekt}), können aber durch höhere Produktivität der eingesetzten
Arbeit kompensiert werden (\emph{Kosteneinsparungseffekt}). Volkswirtschaftlich schädlich
sind überhöhte Abschlüsse dann, wenn Unternehmen mit Preiserhöhungen, Rationalisierung oder
Produktionsverlagerung ins Ausland reagieren.

\begin{merksatz}[title={\bfseries Kostenniveauneutraler Verteilungsspielraum}]
\[
\text{Verteilungsspielraum} = \text{Produktivitätsfortschritt} + \text{Zielinflationsrate}
\]
Bleiben die Lohnabschlüsse innerhalb dieses Spielraums, steigen die
\emph{Lohnstückkosten} nicht stärker als das Preisziel der EZB --- es droht keine
Lohnkosteninflation.
\end{merksatz}

\begin{beispiel}
Die Arbeitsproduktivität steigt um 1,5~\%, die Zielinflation der EZB beträgt 2,0~\%.\\[0.3em]
\textbf{Schritt 1:} Verteilungsspielraum $= 1{,}5 + 2{,}0 = \textbf{3{,}5~\%}$.\\
\textbf{Schritt 2:} Die Gewerkschaft setzt 5,0~\% durch: $5{,}0 - 3{,}5 =
1{,}5$~Prozentpunkte über dem Spielraum.\\[0.3em]
\emph{Folge:} Die Lohnstückkosten steigen --- die Unternehmen reagieren mit
Preiserhöhungen (Inflationsdruck), Rationalisierung oder Verlagerung
(lohnkostenbedingte Arbeitslosigkeit).
\end{beispiel}

\subsection{Arbeitsmarktpolitik}

Der Arbeitsmarkt ist kein „freier Markt`` im Sinne der Preisbildung: Wegen der hohen
sozialen Bedeutung des Lohns besteht Konsens, dass der Arbeitsmarkt nicht vollständig dem
freien Spiel der Marktkräfte überlassen wird. Zuständig für die Arbeitsmarktpolitik im
engeren Sinne ist v.\,a. die \textbf{Bundesagentur für Arbeit} (Nürnberg) mit ihren
Regionalagenturen, Arbeitsagenturen und Jobcentern.

\begin{itemize}
  \item \textbf{Passive Arbeitsmarktpolitik:} mildert die wirtschaftlichen Folgen
    eingetretener oder drohender Arbeitslosigkeit --- \emph{Lohnersatzleistungen} wie
    Arbeitslosengeld I und Bürgergeld (bis 2022: ALG II), Insolvenzgeld, Kurzarbeitergeld.
  \item \textbf{Aktive Arbeitsmarktpolitik:} versetzt Erwerbspersonen in die Lage, im
    Arbeitsprozess zu bleiben oder schnell wieder einzutreten --- Beratung und
    \emph{Vermittlung}, Aktivierung und berufliche Eingliederung (Trainings- und
    Qualifizierungsmaßnahmen), Förderung der \emph{Aus- und Weiterbildung}, Umschulung und
    berufliche Rehabilitation, \emph{Existenzgründungsförderung}, Mobilitätshilfen (Reise-
    und Umzugskosten).
\end{itemize}

Leitgedanke ist heute \textbf{„Fördern und Fordern``}: Arbeitslose müssen sich aktiv
bemühen und zumutbare Jobs annehmen, sonst drohen Kürzungen der Lohnersatzleistungen.
Kritikpunkte an der Arbeitsmarktpolitik: geschönte Arbeitslosenquoten (Geförderte fallen
aus der Statistik), teils ineffiziente Qualifizierungsmaßnahmen, Verschlechterung der
Qualität der Arbeitsverhältnisse (Befristung, Leiharbeit, Scheinselbstständigkeit).

\subsubsection{Formen der Arbeitslosigkeit}

Die richtige Gegenmaßnahme hängt von der \emph{Ursache} der Arbeitslosigkeit ab:

\begin{center}
\begin{tabular}{p{3.0cm} p{6.0cm} p{4.6cm}}
\toprule
\rowcolor{bfwTabHead}\textbf{Form} & \textbf{Ursache} & \textbf{Typische Gegenmaßnahme}\\
\midrule
friktionell & kurze Suchphase beim Arbeitsplatzwechsel („Sucharbeitslosigkeit``) & schnellere Vermittlung, Jobbörsen\\
saisonal & jahreszeitliche Schwankungen (Bau, Tourismus, Landwirtschaft) & Saison-Kurzarbeitergeld, Winterbauförderung\\
konjunkturell & Nachfragerückgang in der Rezession & expansive Geld- und Fiskalpolitik, Kurzarbeitergeld\\
strukturell & dauerhafte Diskrepanz zwischen Angebot und Nachfrage: Qualifikation, Region, Branche, technischer Wandel & Qualifizierung, Umschulung, Mobilitätshilfen, Strukturpolitik\\
\bottomrule
\end{tabular}
\end{center}

Der „harte Kern``, der auch in der Hochkonjunktur nicht abgebaut wird, heißt
\textbf{Sockelarbeitslosigkeit} (überwiegend strukturell bedingt).

\begin{eselsbruecke}
Die vier Formen als \textbf{„vier W``}: friktionell = \textbf{W}echsel, saisonal =
\textbf{W}etter, konjunkturell = \textbf{W}elle, strukturell = \textbf{W}andel.
Und: \textbf{Passiv = Portemonnaie} (Geld ersetzt Lohn), \textbf{Aktiv = Arbeit}
(zurück in den Job).
\end{eselsbruecke}

\subsection{Umweltpolitik}

\subsubsection{Prinzipien der Umweltpolitik}

\begin{itemize}
  \item \textbf{Verursacherprinzip:} Die Kosten für Vermeidung, Beseitigung oder Ausgleich
    von Umweltbelastungen trägt der Verursacher --- externe Kosten werden
    \emph{internalisiert} (gehen in Kostenrechnung und Preise ein); Beispiele:
    Rücknahmepflicht für Umverpackungen, Abwasserabgabe.
  \item \textbf{Gemeinlastprinzip:} Ist kein Verursacher feststellbar, übernimmt der Staat
    die Maßnahmen --- finanziert über Steuern (z.\,B. kommunale Kläranlagen, Mülldeponien);
    Nachteil: keine Anreize für umweltbewusstes Handeln.
  \item \textbf{Vorsorgeprinzip:} Umweltgefahren vorausschauend vermeiden --- „besser
    Vorsorge statt Nachsorge`` (viele Schäden, z.\,B. das Ozonloch, sind kaum reparabel).
  \item \textbf{Kooperationsprinzip:} Zusammenarbeit von Bürgern, Unternehmen und Staat.
  \item \textbf{Nachhaltigkeitsprinzip} (sustainable development, ursprünglich aus der
    Forstwirtschaft): künftige Generationen nicht schlechterstellen --- erneuerbare
    Ressourcen nur im Umfang ihrer Regenerationsfähigkeit nutzen, nicht erneuerbare sparsam
    nutzen, Emissionen an der Aufnahmefähigkeit der Ökosysteme ausrichten (vgl. Agenda~21).
\end{itemize}

\subsubsection{Instrumente der staatlichen Umweltpolitik}

\begin{itemize}
  \item \textbf{Nicht fiskalische Instrumente:} \emph{Aufklärung und Information};
    \emph{Auflagen} (Ordnungsrecht): umweltbezogene Ge- und Verbote, z.\,B. Grenzwerte für
    Schadstoffausstoß, FCKW-Verbot, Umweltplaketten/Umweltzonen.
  \item \textbf{Fiskalische Instrumente:}
    \emph{Umweltabgaben} --- Steuern (Ökosteuer), Sonderabgaben (Abwasserabgabe), Gebühren
    (Abfall-/Abwassergebühren) verteuern umweltschädliches Verhalten;
    \emph{Umweltzertifikate} --- Umweltbelastung nur gegen Berechtigung: Ein Zertifikat
    gestattet eine bestimmte Menge Schadstoffausstoß pro Jahr; die Emissionsrechte werden
    börsenmäßig gehandelt (EU-Emissionshandel für CO\textsubscript{2});
    \emph{Umweltsubventionen} --- Zuschüsse oder Steuererleichterungen für
    umweltschützende Maßnahmen (z.\,B. Wärmedämmung).
\end{itemize}

\begin{eselsbruecke}
Die drei großen Instrumentenfamilien: \textbf{A-A-Z} --- \textbf{A}uflagen (Ordnungsrecht:
verbieten), \textbf{A}bgaben (verteuern), \textbf{Z}ertifikate (begrenzen und handeln).
\end{eselsbruecke}

\subsection{Nachfrage- vs. angebotsorientierte Wirtschaftspolitik}

Nachfrage- und Angebotstheorie liefern unterschiedliche idealtypische Konzepte zur
Bewältigung von Wachstums- und Beschäftigungskrisen. Sie unterscheiden sich darin, welcher
\emph{Marktseite} die größere Bedeutung beigemessen wird und welche \emph{Rolle der Staat}
spielen soll.

\subsubsection{Nachfragetheorie nach Keynes (Fiskalismus)}

Die Nachfragetheorie geht auf \textbf{John Maynard Keynes} (1883--1946) zurück und sieht
Wachstumsschwächen und Arbeitslosigkeit in einer \emph{unzureichenden gesamtwirtschaftlichen
Nachfrage} begründet. Da der Markt allein das Gleichgewicht nicht herstellt, soll der Staat
eine \textbf{aktive Rolle} übernehmen und für zusätzliche Nachfrage sorgen
(\emph{Symptombekämpfung}, kurzfristig): Stärkung der Massenkaufkraft durch Lohnerhöhungen
und staatliche Zuschüsse, Erhöhung des Staatskonsums durch öffentliche Ausgabenprogramme,
Ausweitung des öffentlichen Sektors --- notfalls über Verschuldung (\emph{Deficit
Spending}). Geldpolitisch soll die Zentralbank die Finanzpolitik unterstützen (Liquidität
bereitstellen, Zinsen niedrig halten); Fiskal- und Geldpolitik wirken \emph{antizyklisch}.
Anhänger heißen \textbf{Fiskalisten}. Grundlage der Lohnargumentation ist die
\emph{Kaufkrafttheorie}.

\textbf{Schwächen:} zeitliche Verzögerungen (\emph{Timelags}) zwischen Entscheidung und
Wirkung, Verdrängung privater Investitionen durch staatliche (\emph{Crowding-out}),
steigende Staatsquote sowie steigende Abgaben-, Steuer- und Schuldenlast.

\subsubsection{Angebotstheorie (Neoklassik, Monetarismus)}

Die Angebotstheorie geht auf die klassische Nationalökonomie zurück
(\textbf{Neoklassizismus}; prominenter Vertreter: \textbf{Milton Friedman}, 1912--2006).
Beschäftigung und Wachstum hängen demnach von den Bedingungen auf der
\emph{Angebotsseite} ab; der Staat soll sich \textbf{passiv} verhalten und nur die
\emph{Rahmenbedingungen} verbessern (\emph{Ursachenbekämpfung}, mittel- bis langfristig):
Kostendämpfung zur Erhöhung der Unternehmensrentabilität, Steuersenkungen, Abbau von
Subventionen und Regulierungen (Deregulierung), Verringerung des effizienzschwachen
Staatskonsums, Investitionsförderung, moderate Lohnabschlüsse. Motor der Entwicklung sind
die \emph{Investitionen der Unternehmen}: Bessere Renditeerwartungen führen zu
Investitionen und neuen Arbeitsplätzen. Die Forderung nach \emph{stetigem
Geldmengenwachstum} basiert auf dem \textbf{Monetarismus}: Zwischen Geldmenge und BIP
besteht eine enge Beziehung.

\textbf{Schwächen:} einseitige Betonung der Kosten --- Löhne sind auch ein
\emph{Nachfragefaktor}; eine bessere Gewinnsituation erhöht zwar die
Investitions\emph{fähigkeit}, nicht automatisch die Investitions\emph{bereitschaft} (bei
geringer Absatzerwartung bleiben Erweiterungsinvestitionen aus).

\begin{eselsbruecke}
\textbf{Keynes = Nachfrage = Staat aktiv = kurzfristig Symptome behandeln.}
\textbf{Angebot = Rahmenbedingungen = Staat passiv = langfristig Ursachen beheben.}
Merkhilfe: „\textbf{K}eynes \textbf{k}auft`` (Staat schafft Nachfrage) ---
„\textbf{F}riedman \textbf{f}reit die Märkte`` (Deregulierung).
\end{eselsbruecke}

\subsubsection{Gemeinsamer Zielrahmen: das magische Viereck (StabG)}

Beide Konzeptionen streben dieselben gesamtwirtschaftlichen Ziele des
\textbf{Stabilitätsgesetzes (1967)} an --- sie unterscheiden sich nur im \emph{Weg}
dorthin. Die vier Ziele bilden das \textbf{magische Viereck} (Wiederholung aus
Lektion~2): „magisch``, weil nie alle vier Ziele gleichzeitig vollständig erreichbar
sind.

\begin{center}
\begin{tikzpicture}[font=\small,
  ziel/.style={draw=bfwPrimaryDark, fill=bfwBgTeal, thick, rounded corners=4pt,
    align=center, text width=3.6cm, inner sep=5pt}]
  % Diagonalen (Spannungsbögen zwischen allen Zielen)
  \draw[dashed, thick, bfwInkSoft!70] (0,3.4) -- (0,-3.4);
  \draw[dashed, thick, bfwInkSoft!70] (-5.2,0) -- (5.2,0);
  % Die vier Ziele
  \node[ziel] (p) at (0,3.4)
    {\textbf{Preisniveaustabilität}\\{\footnotesize Inflationsrate $\approx$ 2~\% (HVPI)}};
  \node[ziel] (b) at (5.2,0)
    {\textbf{hoher Beschäftigungsstand}\\{\footnotesize niedrige Arbeitslosenquote}};
  \node[ziel] (w) at (0,-3.4)
    {\textbf{stetiges und angemessenes Wirtschaftswachstum}\\{\footnotesize reales BIP-Wachstum}};
  \node[ziel] (a) at (-5.2,0)
    {\textbf{außenwirtschaftliches Gleichgewicht}\\{\footnotesize ausgeglichene Leistungsbilanz}};
  % Seiten des Vierecks
  \draw[thick, bfwPrimary] (p) -- (a);
  \draw[thick, bfwPrimary] (a) -- (w);
  \draw[thick, bfwPrimary] (w) -- (b);
  % Klassischer Zielkonflikt als Spannungsbogen
  \draw[line width=1.2pt, bfwAccent!85!black, decorate,
        decoration={zigzag, amplitude=0.9mm, segment length=4mm}] (p) -- (b);
  \node[text=bfwAccent!85!black, font=\scriptsize\bfseries, align=center, anchor=west]
    at (3.6,2.1) {klassischer\\Zielkonflikt};
  % Zentrum
  \node[fill=white, draw=bfwInkSoft!50, rounded corners=3pt, inner sep=4pt, align=center,
        font=\scriptsize, text=bfwInkSoft] at (0,0)
    {„magisch``:\\nie alle vier Ziele\\gleichzeitig erreichbar};
\end{tikzpicture}
\end{center}

\emph{Klassischer Zielkonflikt:} Preisniveaustabilität $\leftrightarrow$ hoher
Beschäftigungsstand --- eine expansive Politik schafft Beschäftigung, treibt aber die
Preise (und umgekehrt). \emph{Zielharmonie:} Wachstum und Beschäftigung stützen sich
gegenseitig. Erweitert um \emph{gerechte Einkommensverteilung} und \emph{Umweltschutz}
(bzw. lebenswerte Umwelt) spricht man vom \textbf{magischen Sechseck}.

\subsection{Außenwirtschaft: Freihandel oder Protektionismus}

Unter \textbf{Außenwirtschaft} versteht man die Gesamtheit aller wirtschaftlichen
Aktivitäten zwischen Wirtschaftssubjekten im Inland und im Ausland (Handel mit Gütern und
Dienstleistungen, Kapital- und Devisenverkehr; ausgewiesen in der \emph{Zahlungsbilanz}).
Gründe für internationalen Handel: unterschiedliche klimatische Bedingungen, ungleich
verteilte Rohstoff- und Energievorräte, unterschiedliche Spezialisierungen und
Technologieschwerpunkte --- Grundgedanke der \textbf{internationalen Arbeitsteilung}.

\subsubsection{Freihandel}

\textbf{Freihandel} ist der völlig unbehinderte internationale Güteraustausch --- der Staat
verzichtet auf marktbeeinflussende Maßnahmen. Vorteile: große Güterauswahl in verschiedenen
Qualitäten und Preislagen für die Verbraucher, \emph{Spezialisierungsvorteile} und
Kostensenkungen in der Produktion, Absatzmärkte im In- und Ausland, Innovationsförderung
durch internationalen Wettbewerb. Gewinner sind allerdings v.\,a. die starken,
international wettbewerbsfähigen Unternehmen.

\subsubsection{Protektionismus: tarifäre und nichttarifäre Handelshemmnisse}

Greift der Staat zum Schutz der einheimischen Wirtschaft in den freien Handel ein, spricht
man von \textbf{Protektionismus} (Motive: Schutz vor Billigimporten, Vermeidung von
Monostrukturen und Auslandsabhängigkeit).

\begin{itemize}
  \item \textbf{Tarifäre Handelshemmnisse} (preispolitisch): \emph{Importzölle} (größte
    Bedeutung --- verbessern die Wettbewerbsposition der einheimischen Wirtschaft und
    bringen Staatseinnahmen), Steuertarife, \emph{Exportsubventionen} (inländische Produkte
    können im Ausland günstiger angeboten werden). Innerhalb der EU werden keine Zölle
    erhoben; gegenüber Drittländern gilt ein einheitliches EU-Vorgehen (keine autonome
    deutsche Zollpolitik mehr).
  \item \textbf{Nichttarifäre Handelshemmnisse} (mengenpolitisch/administrativ):
    \emph{Importkontingente} (mengenmäßige Beschränkungen), \emph{Ein- und Ausfuhrverbote}
    (schärfste Form, z.\,B. Rüstungsgüter, Drogen), \emph{Embargos} (außen- oder
    sicherheitspolitisch angeordnete Beschränkungen, i.\,d.\,R. Umsetzung internationaler
    Sanktionen von UNO oder EU), administrative Erschwernisse (Abfertigungs- und
    Genehmigungsverfahren), technische Normen, Gesundheits-, Sicherheits- und
    Verbraucherschutzstandards --- letztere mit stark zunehmender Bedeutung.
\end{itemize}

\begin{merksatz}
Rechtsgrundlagen in Deutschland: \textbf{Außenwirtschaftsgesetz (AWG)} und
Außenwirtschaftsverordnung (AWV); das \textbf{BAFA} (Bundesamt für Wirtschaft und
Ausfuhrkontrolle) informiert über Verbote, Genehmigungen und Beschränkungen.
\textbf{Globalisierung} = weltweiter Austausch von Gütern, Dienstleistungen, Informationen
und Kapital sowie weltweite Vernetzung von Unternehmen. Die EU ist Mitglied der
\textbf{WTO} und schließt Freihandelsabkommen (z.\,B. CETA mit Kanada, JEFTA mit Japan).
\end{merksatz}

\subsubsection{Zahlungsbilanz}

Die \textbf{Zahlungsbilanz} erfasst systematisch alle wirtschaftlichen Transaktionen
zwischen In- und Ausland innerhalb einer Periode. Ihre wichtigsten Teilbilanzen:

\begin{itemize}
  \item \textbf{Leistungsbilanz:} \emph{Handelsbilanz} (Warenexporte und -importe),
    \emph{Dienstleistungsbilanz} (z.\,B. Reiseverkehr, Transport),
    \emph{Primäreinkommen} (Erwerbs- und Vermögenseinkommen, z.\,B. Zinsen und
    Dividenden) und \emph{Sekundäreinkommen} (laufende Übertragungen ohne Gegenleistung,
    z.\,B. Überweisungen ausländischer Arbeitskräfte, EU-Beiträge);
  \item \textbf{Kapitalbilanz:} Kapitalexporte und -importe (Direktinvestitionen,
    Wertpapieranlagen, Kredite);
  \item \textbf{Devisenbilanz:} Veränderung der Währungsreserven der Zentralbank.
\end{itemize}

\begin{center}
\begin{tikzpicture}[font=\footnotesize, line cap=round,
  kasten/.style={draw=bfwPrimaryDark, fill=bfwBgTeal, thick, rounded corners=4pt,
    align=center, inner sep=5pt, text width=3.6cm},
  blatt/.style={draw=bfwInkSoft, fill=bfwBgSoft, thick, rounded corners=3pt,
    align=left, inner sep=4pt, anchor=west, text width=4.9cm}]
  % Wurzel
  \node[draw=bfwPrimaryDark, fill=bfwPrimaryDark, text=white, thick, rounded corners=4pt,
        align=center, inner sep=6pt, font=\footnotesize\bfseries] (z) at (7.0,3.4)
    {Zahlungsbilanz\\{\scriptsize formal stets ausgeglichen}};
  % Teilbilanzen
  \node[kasten] (l) at (2.4,1.6)
    {\textbf{Leistungsbilanz}\\{\scriptsize realwirtschaftliche Transaktionen}};
  \node[kasten] (k) at (7.0,1.6)
    {\textbf{Kapitalbilanz}\\{\scriptsize Direktinvestitionen, Wertpapiere, Kredite}};
  \node[kasten] (d) at (11.6,1.6)
    {\textbf{Devisenbilanz}\\{\scriptsize Veränderung der Währungsreserven}};
  \draw[thick, bfwPrimaryDark] (z.south) -- ++(0,-0.25) -| (l.north);
  \draw[thick, bfwPrimaryDark] (z.south) -- (k.north);
  \draw[thick, bfwPrimaryDark] (z.south) -- ++(0,-0.25) -| (d.north);
  % Unterbilanzen der Leistungsbilanz
  \node[blatt] (t1) at (2.9,0.3)
    {\textbf{Handelsbilanz:} Warenexporte und -importe};
  \node[blatt] (t2) at (2.9,-0.55)
    {\textbf{Dienstleistungsbilanz:} z.\,B. Reiseverkehr, Transport};
  \node[blatt] (t3) at (2.9,-1.4)
    {\textbf{Primäreinkommen:} Löhne, Zinsen, Dividenden};
  \node[blatt] (t4) at (2.9,-2.25)
    {\textbf{Sekundäreinkommen:} Übertragungen, z.\,B. EU-Beiträge};
  \draw[thick, bfwPrimaryDark] (l.south) -- (2.4,-2.25);
  \draw[thick, bfwPrimaryDark] (2.4,0.3) -- (t1.west);
  \draw[thick, bfwPrimaryDark] (2.4,-0.55) -- (t2.west);
  \draw[thick, bfwPrimaryDark] (2.4,-1.4) -- (t3.west);
  \draw[thick, bfwPrimaryDark] (2.4,-2.25) -- (t4.west);
\end{tikzpicture}
\end{center}

\begin{merksatz}
\textbf{Die Zahlungsbilanz ist formal immer ausgeglichen:} Einem Leistungsbilanzüberschuss
steht ein entsprechender Kapitalexport gegenüber (doppelte Verbuchung). Ungleichgewichte
zeigen sich daher in den \emph{Teilbilanzen} --- Deutschland erzielt traditionell hohe
\textbf{Leistungsbilanzüberschüsse} (Exportstärke), was international auch kritisiert wird.
\end{merksatz}

\begin{beispiel}
Für ein Land liegen folgende Werte vor (Mrd.~Euro): Warenexporte 1.560, Warenimporte
1.320, Saldo der Dienstleistungsbilanz $-20$, Primäreinkommen $+130$, Sekundäreinkommen
$-60$.\\[0.3em]
\textbf{Schritt 1:} Handelsbilanzsaldo $= 1.560 - 1.320 = +240$ (Überschuss).\\
\textbf{Schritt 2:} Leistungsbilanzsaldo $= 240 - 20 + 130 - 60 =
\textbf{+290~Mrd.~Euro}$ (Überschuss).\\[0.3em]
\emph{Interpretation:} Dem Überschuss steht spiegelbildlich ein Nettokapitalexport in
gleicher Höhe gegenüber --- das Land baut Auslandsvermögen auf.
\end{beispiel}

\begin{achtung}
\textbf{Handelsbilanz} $\neq$ \textbf{Leistungsbilanz}: Die Handelsbilanz erfasst nur den
\emph{Warenhandel}; die Leistungsbilanz umfasst zusätzlich Dienstleistungen sowie Primär-
und Sekundäreinkommen. Wenn Prüfungsaufgaben vom „Exportüberschuss`` sprechen, ist meist
die Handelsbilanz gemeint --- genau lesen.
\end{achtung}

\subsubsection{Wechselkurse}

Der \textbf{Wechselkurs} ist der Preis einer Währung, ausgedrückt in einer anderen
Währung. \textbf{Preisnotierung:} Preis einer ausländischen Währungseinheit in Inlandswährung
(z.\,B. 1~USD = 0,90~€); \textbf{Mengennotierung} (in der Eurozone üblich): Menge
ausländischer Währung je Euro (z.\,B. 1~€ = 1,10~USD).

\begin{itemize}
  \item \textbf{Aufwertung} des Euro: Importe werden billiger, Exporte verteuern sich ---
    die preisliche Wettbewerbsfähigkeit der Exporteure sinkt.
  \item \textbf{Abwertung} des Euro: Exporte werden für ausländische Käufer billiger
    (Exportförderung), Importe verteuern sich (Gefahr importierter Inflation).
  \item \textbf{Flexible Wechselkurse} bilden sich frei durch Angebot und Nachfrage am
    Devisenmarkt (z.\,B. Euro/US-Dollar); bei \textbf{festen Wechselkursen} garantiert
    die Zentralbank durch Devisenmarktinterventionen eine feste Parität (z.\,B. WKM~II
    vor dem Euro-Beitritt).
  \item \textbf{Kaufkraftparität:} Langfristig tendieren Wechselkurse zu dem Kurs, bei
    dem ein Warenkorb in beiden Währungsräumen gleich viel kostet.
\end{itemize}

\begin{center}
\begin{tikzpicture}[font=\small,
  wk/.style 2 args={draw=#1, fill=#2, thick, rounded corners=4pt, align=center,
    text width=4.1cm, minimum height=1.35cm, inner sep=5pt},
  pfeil/.style={-{Stealth[length=3mm]}, line width=1.2pt, bfwInkSoft}]
  % Aufwertung (obere Kette)
  \node[wk={bfwAccent!85!black}{bfwBgAmber}] (a1) at (0,0)
    {\textbf{Aufwertung des Euro}\\{\footnotesize z.\,B. 1~€ = 1,10 $\rightarrow$ 1,25~USD}};
  \node[wk={bfwAccent!85!black}{bfwBgAmber}] (a2) at (5.85,0)
    {\textbf{Exporte verteuern sich,}\\\textbf{Importe werden billiger}};
  \node[wk={bfwAccent!85!black}{bfwBgAmber}] (a3) at (11.7,0)
    {\textbf{Exportnachfrage sinkt}\\{\footnotesize Wettbewerbsfähigkeit $\downarrow$,
     dämpft aber die Preise im Inland}};
  % Abwertung (untere Kette)
  \node[wk={bfwPrimaryDark}{bfwBgTeal}] (b1) at (0,-2.5)
    {\textbf{Abwertung des Euro}\\{\footnotesize z.\,B. 1~€ = 1,10 $\rightarrow$ 1,00~USD}};
  \node[wk={bfwPrimaryDark}{bfwBgTeal}] (b2) at (5.85,-2.5)
    {\textbf{Exporte werden billiger,}\\\textbf{Importe verteuern sich}};
  \node[wk={bfwPrimaryDark}{bfwBgTeal}] (b3) at (11.7,-2.5)
    {\textbf{Exportnachfrage steigt}\\{\footnotesize Exportförderung, aber Gefahr
     importierter Inflation}};
  \draw[pfeil] (a1) -- (a2);
  \draw[pfeil] (a2) -- (a3);
  \draw[pfeil] (b1) -- (b2);
  \draw[pfeil] (b2) -- (b3);
\end{tikzpicture}
\end{center}

\begin{achtung}
Rechenregel bei \textbf{Mengennotierung} (1~€ = x~USD): Auslandspreis \textbf{geteilt
durch} den Kurs = Europreis; Europreis \textbf{mal} Kurs = Auslandspreis. Und:
\emph{Steigt} die Mengennotierung (mehr Dollar je Euro), wertet der Euro \emph{auf} ---
bei der Preisnotierung ist es genau umgekehrt. Diese Verwechslung ist ein
Prüfungsklassiker.
\end{achtung}

\begin{beispiel}
Kurs: 1~€ = 1,10~USD (Mengennotierung). Ein deutscher Händler importiert eine
US-Maschine für 220.000~USD; ein Maschinenbauer exportiert eine Anlage für 200.000~€
in die USA.\\[0.3em]
\textbf{Schritt 1 (Import):} $220.000 \div 1{,}10 = 200.000$~€ Einkaufspreis.\\
\textbf{Schritt 2 (Export):} $200.000 \times 1{,}10 = 220.000$~USD Preis für den
US-Kunden.\\
\textbf{Schritt 3 (Aufwertung auf 1~€ = 1,25~USD):} Import: $220.000 \div 1{,}25 =
176.000$~€ (billiger); Export: $200.000 \times 1{,}25 = 250.000$~USD (teurer).\\[0.3em]
\emph{Interpretation:} Die Aufwertung entlastet Importeure und Verbraucher, verschlechtert
aber die preisliche Wettbewerbsfähigkeit der Exporteure --- genau die Wirkungskette aus
der Grafik oben.
\end{beispiel}

\subsection{Besonderheiten der EU}

Die Integration begann 1957 mit der \emph{Europäischen Wirtschaftsgemeinschaft (EWG)}
(Römische Verträge, sechs Gründungsstaaten); mit dem Vertrag von \textbf{Maastricht} (1993)
entstand die \emph{Europäische Union (EU)}.

\subsubsection{Stufen der wirtschaftlichen Integration}

Die europäische Integration verlief über mehrere Stufen --- jede Stufe schließt die
vorhergehende mit ein:

\begin{center}
\begin{tikzpicture}[font=\footnotesize,
  stufe/.style 2 args={draw=#1, fill=#2, thick, rounded corners=4pt, align=center,
    text width=3.15cm, minimum height=2.0cm, inner sep=4pt, anchor=south west},
  pfeil/.style={-{Stealth[length=2.5mm]}, thick, bfwInkSoft}]
  \node[stufe={bfwPrimaryDark}{bfwBgTeal}] (s1) at (0,0)
    {\textbf{1. Freihandelszone}\\{\scriptsize keine Binnenzölle\\(z.\,B. EFTA)}};
  \node[stufe={bfwPrimaryDark}{bfwPrimary!25}] (s2) at (4.3,0.7)
    {\textbf{2. Zollunion}\\{\scriptsize + gemeinsamer Außenzoll\\(EWG ab 1968)}};
  \node[stufe={bfwPrimaryDark}{bfwPrimary!45}] (s3) at (8.6,1.4)
    {\textbf{3. Binnenmarkt}\\{\scriptsize + vier Freiheiten\\(EU seit 1993)}};
  \node[stufe={bfwPrimaryDark}{bfwPrimaryDark}, text=white] (s4) at (12.9,2.1)
    {\textbf{4. Wirtschafts- und Währungsunion}\\{\scriptsize + Euro und gemeinsame
     Geldpolitik (seit 1999/2002)}};
  \draw[pfeil] (s1.east) -- (s2.west);
  \draw[pfeil] (s2.east) -- (s3.west);
  \draw[pfeil] (s3.east) -- (s4.west);
\end{tikzpicture}
\end{center}

Denkbare Endstufe wäre die \emph{politische Union}. Die Eurozonen-Staaten stehen heute
auf Stufe~4, alle übrigen EU-Staaten mindestens auf Stufe~3 (Binnenmarkt).

\subsubsection{Der Europäische Binnenmarkt: die vier Freiheiten}

Ein \textbf{Binnenmarkt} ist ein Wirtschaftsraum ohne innere Grenzen (gegründet 1993). Er
ist durch \textbf{vier Freiheiten} gekennzeichnet:

\begin{center}
\begin{tabular}{p{4.2cm} p{9.6cm}}
\toprule
\textbf{Freiheit} & \textbf{Kennzeichen}\\
\midrule
Freier Personenverkehr & Abbau der Grenzkontrollen (Schengen), Aufenthalts- und Niederlassungsfreiheit, freie Arbeitsplatzwahl, Anerkennung von Berufsabschlüssen\\
Freier Warenverkehr & keine Zölle oder mengenmäßigen Beschränkungen, Harmonisierung von Umsatz-/Verbrauchssteuern und technischen Normen\\
Freier Dienstleistungsverkehr & Selbstständige und freie Berufe dürfen EU-weit anbieten und sich niederlassen; Liberalisierung bei Banken und Versicherungen\\
Freier Geld- und Kapitalverkehr & Öffnung der Finanzmärkte, erleichterter Zahlungsverkehr (SEPA), Aufhebung von Kapitalverkehrsbeschränkungen\\
\bottomrule
\end{tabular}
\end{center}

\textbf{Chancen:} größerer Absatzmarkt, sinkende Stückkosten durch höhere Stückzahlen und
Rationalisierung, vielfältigeres und günstigeres Güterangebot, stabilisierende Wirkung auf
das Preisniveau. \textbf{Risiken/Nachteile:} Verlagerung von Produktionsstätten ins
kostengünstigere EU-Ausland (Arbeitsplatzverluste), Bürokratiekosten, aus Sicht der
Kritiker „Entmachtung`` der Einzelstaaten.

\begin{eselsbruecke}
Die vier Freiheiten: \textbf{P-W-D-K} --- \textbf{P}ersonen, \textbf{W}aren,
\textbf{D}ienstleistungen, \textbf{K}apital. Merkspruch: „\emph{Praktisch wandern durch
Kontinente}``.
\end{eselsbruecke}

\subsubsection{Europäische Wirtschafts- und Währungsunion (EWWU)}

Die Europäische Währungsunion ist Teil der EWWU; die Teilnehmerstaaten bilden die
\textbf{Eurozone}. Nach dem Maastrichter Vertrag dürfen nur EU-Staaten den Euro einführen,
die ihre Stabilität durch die \textbf{Konvergenzkriterien} bewiesen haben:

\begin{enumerate}
  \item \textbf{Inflationsrate:} im Jahr vor dem Beitritt höchstens 1,5~Prozentpunkte über
    dem Mittelwert der drei preisstabilsten EU-Staaten;
  \item \textbf{Nettoneuverschuldung:} nicht dauerhaft mehr als \textbf{3~\%} des BIP;
  \item \textbf{Gesamtschuldenstand:} i.\,d.\,R. nicht mehr als \textbf{60~\%} des BIP;
  \item \textbf{Zinssatz für langfristige Kredite:} höchstens 2~Prozentpunkte über dem
    Mittelwert der drei preisstabilsten EU-Staaten;
  \item \textbf{Wechselkursstabilität:} zwei Jahre stabiler Außenwert ohne Abwertung
    (Teilnahme am Wechselkursmechanismus WKM~II).
\end{enumerate}

\begin{eselsbruecke}
Konvergenzkriterien als Zahlenreihe \textbf{„1,5 --- 3 --- 60 --- 2 --- 2``}: höchstens
\textbf{1,5}~Punkte mehr Inflation, \textbf{3}~\% Neuverschuldung, \textbf{60}~\%
Schuldenstand, \textbf{2}~Punkte mehr Zinsen, \textbf{2}~Jahre stabiler Wechselkurs.
\end{eselsbruecke}

Die Einhaltung wurde in der Vergangenheit nicht konsequent überprüft; mit der Euro- und
Staatsschuldenkrise reformierten EU-Parlament und Finanzminister 2011 den europäischen
\emph{Stabilitäts- und Wachstumspakt} (härteres Vorgehen gegen „Haushaltssünder``; die
Referenzwerte 3~\% und 60~\% blieben erhalten). 2024 wurde der Pakt grundlegend
reformiert (länderspezifische Ausgabenpfade); die Referenzwerte 3~\% (Defizit) und
60~\% (Schuldenstand) gelten weiter.

\subsubsection{Vor- und Nachteile der gemeinsamen Währung}

\begin{itemize}
  \item \textbf{Vorteile:} keine Wechselkursrisiken und Umtauschkosten im Euroraum,
    Preistransparenz (direkter Preisvergleich), Förderung von Handel, Investitionen und
    Reiseverkehr, starke Währung im Weltmaßstab.
  \item \textbf{Nachteile:} Verzicht auf eigenständige nationale Geld- und
    Wechselkurspolitik (keine Abwertung zur Stärkung der Wettbewerbsfähigkeit möglich),
    einheitlicher Leitzins für unterschiedlich entwickelte Volkswirtschaften,
    Haftungs- und Transferrisiken bei Schuldenkrisen einzelner Mitglieder.
\end{itemize}

\subsection{Formeln \& Rechenbeispiele}

Die folgenden Berechnungen sind Prüfungsklassiker --- alle folgen demselben Muster:
\emph{Teil durch Ganzes mal 100}. Prägen Sie sich Formel und Rechenweg anhand der
einfachen Zahlenbeispiele ein.

\begin{merksatz}[title={\bfseries Formel: Arbeitslosenquote}]
\[
\text{Arbeitslosenquote in \%} =
\frac{\text{registrierte Arbeitslose}}{\text{alle zivilen Erwerbspersonen}} \times 100
\]
Zivile Erwerbspersonen = Erwerbstätige \emph{plus} registrierte Arbeitslose (also alle,
die arbeiten oder arbeiten wollen).
\end{merksatz}

\begin{beispiel}
In einem Land sind 2,3~Mio. Personen arbeitslos gemeldet; es gibt 43,7~Mio. Erwerbstätige.\\[0.3em]
\textbf{Schritt 1:} Erwerbspersonen bestimmen: $43{,}7 + 2{,}3 = 46$~Mio.\\
\textbf{Schritt 2:} Arbeitslose durch Erwerbspersonen teilen: $2{,}3 \div 46 = 0{,}05$.\\
\textbf{Schritt 3:} Mal 100: $0{,}05 \times 100 = \textbf{5{,}0~\%}$ Arbeitslosenquote.
\end{beispiel}

\begin{merksatz}[title={\bfseries Anwendung: Konvergenzkriterien nachrechnen (3~\% / 60~\%)}]
\[
\text{Neuverschuldungsquote} = \frac{\text{jährliche Neuverschuldung}}{\text{BIP}} \times 100
\qquad
\text{Schuldenstandsquote} = \frac{\text{Gesamtschulden}}{\text{BIP}} \times 100
\]
Grenzwerte: Neuverschuldung max. \textbf{3~\%}, Gesamtschuldenstand max. \textbf{60~\%}
des BIP.
\end{merksatz}

\begin{beispiel}
Ein Beitrittskandidat hat ein BIP von 500~Mrd.~Euro, eine Neuverschuldung von
20~Mrd.~Euro und Gesamtschulden von 280~Mrd.~Euro.\\[0.3em]
\textbf{Schritt 1:} Neuverschuldung: $20 \div 500 \times 100 = 4{,}0~\%$ ---
\emph{über} der 3-\%-Grenze~$\Rightarrow$ Kriterium \textbf{verletzt}.\\
\textbf{Schritt 2:} Schuldenstand: $280 \div 500 \times 100 = 56{,}0~\%$ ---
\emph{unter} der 60-\%-Grenze~$\Rightarrow$ Kriterium \textbf{erfüllt}.\\
\textbf{Ergebnis:} Das Land erfüllt die fiskalischen Konvergenzkriterien insgesamt
\emph{nicht} (beide Grenzwerte müssen eingehalten werden).
\end{beispiel}

\begin{merksatz}[title={\bfseries Wichtig für die Prüfungsvorbereitung (Block 1.1 komplett)}]
Sie sollten sicher können: die Gegenüberstellung \emph{nachfrage- vs. angebotsorientiert}
(Keynes vs. Friedman) samt dem \emph{magischen Viereck} als gemeinsamem Zielrahmen
(Abschnitt~5.3), \emph{tarifäre vs. nichttarifäre Handelshemmnisse}, die
\emph{Zahlungsbilanz} mit ihren Teilbilanzen, \emph{Wechselkurse} (Auf-/Abwertung,
Mengennotierung, Umrechnung), die \emph{vier Freiheiten} des Binnenmarkts mit den
\emph{Integrationsstufen}, die \emph{Konvergenzkriterien} (3~\%~/ 60~\%), die
Prinzipien und Instrumente der \emph{Umweltpolitik} (A-A-Z), \emph{aktive vs. passive
Arbeitsmarktpolitik} mit den \emph{Formen der Arbeitslosigkeit}, die
\emph{Tarifautonomie} mit Tarifvertragsarten und Verteilungsspielraum sowie die
\emph{Formeln} für Arbeitslosenquote und Konvergenzquoten (Abschnitt~8). Wiederholen Sie
außerdem Teil~1 dieser Lektion (Konjunktur, EZB, Fiskalpolitik),
Lektion~1 (Preisbildung, Marktformen) und Lektion~2 (Wettbewerb, Staatseingriffe, VGR,
magisches Viereck/Sechseck) --- die Quiz beider Teile und das Prüfungstraining helfen dabei.
\end{merksatz}

\vspace{1em}
\hrule
\vspace{0.8em}
\noindent\textsf{\small Quellen: Rahmenplan Wirtschaftsbezogene Qualifikationen (DIHK), Abschnitte 1.1.3.2 und 1.1.4;
Textband Volkswirtschaft (DIHK-Bildungs-GmbH), Kap.~3.2 und Kap.~4, S.~45--64. Nächster Termin (Sa, 24.10.2026):
Start des BWL-Teils --- Betriebliche Funktionen~I.}


\clearpage
\addcontentsline{toc}{part}{Teil B\quad Übungsaufgaben}
\vspace*{5cm}
{\bfwlabelfont\color{bfwPrimary}TEIL B\par}
\vspace{0.5em}
{\Huge\bfseries\color{bfwPrimaryDark}Übungsaufgaben\par}
\vspace{0.6em}
{\color{bfwPrimary}\rule{\linewidth}{2pt}}\par
\vspace{0.6em}
{\large\color{bfwInkSoft}Neue Aufgaben auf IHK-Niveau --- erst selbst lösen, dann in Teil C vergleichen\par}
\clearpage

\clearpage
\section*{Aufgaben zu Lektion 1: Grundbegriffe \& Preisbildung am Markt}
\addcontentsline{toc}{section}{Aufgaben zu Lektion 1: Grundbegriffe \& Preisbildung am Markt}
% ===========================================================================
% Gesamtskript VWL -- Übungsaufgaben Lektion 1: Grundbegriffe & Preisbildung
% (Fragment, wird per \input eingebunden; Lösungen in loesungen-l1.tex)
% Gesamtpunktzahl: 104 Punkte
% ===========================================================================

% ---------------------------------------------------------------------------
\subsection*{Aufgabe~1.1 --- Wirtschaften und ökonomisches Prinzip (11 Punkte)}
\addcontentsline{toc}{subsection}{Aufgabe~1.1 --- Wirtschaften und ökonomisches Prinzip}

Sie sind Mitarbeiterin bzw. Mitarbeiter der \emph{Nordlicht Fahrradmanufaktur GmbH} in Kiel,
die hochwertige E-Lastenräder herstellt. Eine Kundin, Frau Jansen, möchte ihre beiden Kinder
künftig ohne Auto zur Kita bringen. Sie hat dafür 4.500~€ angespart und bestellt schließlich
ein Lastenrad im Onlineshop der Nordlicht GmbH.

In der wöchentlichen Teamsitzung fallen folgende Aussagen:
\begin{itemize}
  \item[(1)] Einkaufsleiterin: „Für das nächste Quartal brauchen wir 2.400 Aluminiumrahmen.
    Ich hole Angebote ein und bestelle beim günstigsten Lieferanten.``
  \item[(2)] Vertriebsleiter: „Für die Fahrradmesse in Friedrichshafen stehen uns
    18.000~€ zur Verfügung. Damit wollen wir so viele Händlerkontakte wie möglich gewinnen.``
  \item[(3)] Geschäftsführer: „Unser Ziel ist klar: Wir produzieren mit den geringstmöglichen
    Kosten die größtmögliche Zahl an Rädern.``
\end{itemize}

\begin{enumerate}[label=\alph*)]
  \item \textbf{Erläutern} Sie am Beispiel von Frau Jansen die Begriffe Bedürfnis, Bedarf
    und Nachfrage. \hfill (3~P.)
  \item \textbf{Ordnen} Sie die Aussagen (1) bis (3) jeweils dem Maximalprinzip, dem
    Minimalprinzip oder keinem ökonomischen Prinzip zu und \textbf{begründen} Sie Ihre
    Zuordnung. \hfill (6~P.)
  \item \textbf{Ordnen} Sie die folgenden Gegebenheiten der Nordlicht GmbH den
    volkswirtschaftlichen Produktionsfaktoren zu: das zwölfköpfige Montageteam, das
    Gewerbegrundstück im Kieler Hafen, der neue Lackierroboter, das patentierte Verfahren
    zur Rahmenverklebung. \hfill (2~P.)
\end{enumerate}

% ---------------------------------------------------------------------------
\subsection*{Aufgabe~1.2 --- Marktarten und Marktformen (12 Punkte)}
\addcontentsline{toc}{subsection}{Aufgabe~1.2 --- Marktarten und Marktformen}

Sie sind Mitarbeiterin bzw. Mitarbeiter der \emph{Hanse Agrarhandel GmbH} in Lüneburg. Das
Unternehmen handelt mit Saatgut, Dünger und Pflanzenschutzmitteln und vermarktet die Ernte
seiner landwirtschaftlichen Kunden.

\begin{enumerate}[label=\alph*)]
  \item \textbf{Ordnen} Sie die folgenden Vorgänge jeweils der zutreffenden Marktart zu:
    \begin{enumerate}[label=(\arabic*)]
      \item Die Hanse Agrarhandel GmbH schreibt zwei Stellen für Lagerfachkräfte aus.
      \item Sie mietet für die Erntesaison eine zusätzliche Lagerhalle an.
      \item Sie nimmt bei ihrer Hausbank einen Investitionskredit über 250.000~€ auf.
      \item Sie verkauft 300~Tonnen Speisekartoffeln an eine Supermarktkette.
    \end{enumerate}
    \hfill (4~P.)
  \item \textbf{Bestimmen} Sie für die folgenden Situationen die Marktform und
    \textbf{begründen} Sie Ihre Antwort:
    \begin{enumerate}[label=(\arabic*)]
      \item Rund 1.800 Landwirte der Region liefern ihre Zuckerrüben an die beiden
        Zuckerfabriken in Norddeutschland.
      \item Eine einzige Reederei betreibt die Fährverbindung zu einer Nordseeinsel und
        befördert jährlich mehrere Hunderttausend Fahrgäste.
      \item Pflanzenschutzmittel werden von wenigen internationalen Chemiekonzernen
        hergestellt und an Tausende Händler und Landwirte verkauft.
    \end{enumerate}
    \hfill (6~P.)
  \item Ein Kollege behauptet: „Der Lüneburger Wochenmarkt ist ein vollkommener Markt, denn
    dort gibt es viele Anbieter und viele Nachfrager.`` \textbf{Beurteilen} Sie diese
    Aussage. \hfill (2~P.)
\end{enumerate}

% ---------------------------------------------------------------------------
\subsection*{Aufgabe~1.3 --- Preiselastizität und Umsatzwirkung (16 Punkte)}
\addcontentsline{toc}{subsection}{Aufgabe~1.3 --- Preiselastizität und Umsatzwirkung}

Sie sind im Controlling der \emph{Bohnenwerk Kaffeerösterei GmbH} in Bremen tätig. Im
letzten Quartal wurden in einem Testmarkt die Preise zweier Produkte erhöht. Alle übrigen
Einflussgrößen blieben unverändert. Es liegen folgende Daten vor (Absatz in Packungen je
Quartal):

\begin{center}
\begin{tabular}{l c c c c}
\toprule
\rowcolor{bfwTabHead}
\textbf{Produkt (je 500\,g)} & \textbf{Preis alt} & \textbf{Preis neu} & \textbf{Absatz alt} & \textbf{Absatz neu}\\
\midrule
Espresso „Crema Oro`` & 14,00~€ & 15,40~€ & 8.000 & 6.800\\
Filterkaffee „Hausmischung`` & 6,00~€ & 6,60~€ & 20.000 & 19.200\\
\bottomrule
\end{tabular}
\end{center}

\begin{enumerate}[label=\alph*)]
  \item \textbf{Berechnen} Sie für beide Produkte die direkte Preiselastizität der
    Nachfrage. Der Rechenweg ist darzustellen. \hfill (4~P.)
  \item \textbf{Interpretieren} Sie beide Ergebnisse und \textbf{begründen} Sie, warum
    die Kunden auf die Preiserhöhungen unterschiedlich stark reagieren. \hfill (4~P.)
  \item \textbf{Berechnen} Sie für beide Produkte den Umsatz vor und nach der
    Preiserhöhung sowie die Umsatzveränderung in Euro und in Prozent. \hfill (4~P.)
  \item Der Vertriebsleiter schlägt vor, die Preiserhöhung bei beiden Produkten dauerhaft
    und bundesweit umzusetzen. \textbf{Beurteilen} Sie diesen Vorschlag. \hfill (4~P.)
\end{enumerate}

% ---------------------------------------------------------------------------
\subsection*{Aufgabe~1.4 --- Marktgleichgewicht aus Angebots- und Nachfragefunktion (17 Punkte)}
\addcontentsline{toc}{subsection}{Aufgabe~1.4 --- Marktgleichgewicht aus Angebots- und Nachfragefunktion}

Sie sind Mitarbeiterin bzw. Mitarbeiter der \emph{Weserholz Energie GmbH}, die Holzpellets
in 15-kg-Säcken vertreibt. Für den regionalen Pelletmarkt, der annähernd die Bedingungen
der vollständigen Konkurrenz erfüllt, hat die Marktforschung folgende Funktionen ermittelt
($p$ = Preis in € je Sack, $x$ = Menge in Säcken pro Tag):

\begin{center}
Nachfrage: \quad $x_N = 1.200 - 40\,p$ \qquad\qquad Angebot: \quad $x_A = 60\,p - 300$
\end{center}

\begin{enumerate}[label=\alph*)]
  \item \textbf{Berechnen} Sie den Gleichgewichtspreis, die Gleichgewichtsmenge und den
    Umsatz im Gleichgewicht. \hfill (4~P.)
  \item \textbf{Berechnen} Sie für einen Preis von 12~€ und für einen Preis von 20~€
    jeweils die angebotene und die nachgefragte Menge. \textbf{Bestimmen} Sie Art und
    Höhe des jeweiligen Überhangs und \textbf{erläutern} Sie die Marktsituation
    (Käufer- oder Verkäufermarkt) sowie die zu erwartende Preisreaktion. \hfill (6~P.)
  \item \textbf{Stellen} Sie Angebots- und Nachfragekurve in einem Preis-Mengen-Diagramm
    \textbf{dar}. Kennzeichnen Sie das Marktgleichgewicht sowie den Überhang bei einem
    Preis von 20~€. \hfill (4~P.)
  \item \textbf{Berechnen} Sie die Konsumentenrente im Gleichgewicht und
    \textbf{erläutern} Sie, was diese Größe aussagt. \hfill (3~P.)
\end{enumerate}

% ---------------------------------------------------------------------------
\subsection*{Aufgabe~1.5 --- Kurvenverschiebungen und Kreuzpreiselastizität (16 Punkte)}
\addcontentsline{toc}{subsection}{Aufgabe~1.5 --- Kurvenverschiebungen und Kreuzpreiselastizität}

Sie arbeiten im Vertrieb der \emph{Alpenquell Getränke GmbH}, eines mittelständischen
Herstellers von Limonaden. Auf dem Limonadenmarkt ereignen sich folgende Entwicklungen:
\begin{itemize}
  \item[(1)] Der Deutsche Wetterdienst sagt einen außergewöhnlich heißen und langen Sommer
    voraus.
  \item[(2)] Der Weltmarktpreis für Zucker steigt innerhalb weniger Monate um 40~\%.
  \item[(3)] Der Preis für Limonade sinkt.
  \item[(4)] In der Branche setzt sich eine neue, deutlich schnellere Abfülltechnik durch.
\end{itemize}

\begin{enumerate}[label=\alph*)]
  \item \textbf{Erläutern} Sie für jede Entwicklung (1) bis (4), ob sie zu einer Bewegung
    auf einer Kurve oder zu einer Verschiebung einer Kurve führt. Geben Sie jeweils an,
    welche Kurve betroffen ist und in welche Richtung sie sich bewegt bzw. verschiebt.
    \hfill (8~P.)
  \item Die Entwicklungen (1) und (2) treten gleichzeitig ein. \textbf{Erläutern} Sie die
    Auswirkungen auf den Gleichgewichtspreis und die Gleichgewichtsmenge. \hfill (4~P.)
  \item Ein Wettbewerber senkt den Preis seines Eistees um 20~\%. Der Absatz der
    Alpenquell-Limonade geht daraufhin von 12.500 auf 11.750 Kisten pro Monat zurück.
    \textbf{Berechnen} Sie die Kreuzpreiselastizität und \textbf{bestimmen} Sie die
    Güterbeziehung zwischen Eistee und Limonade. \hfill (4~P.)
\end{enumerate}

% ---------------------------------------------------------------------------
\subsection*{Aufgabe~1.6 --- Funktionen des Marktpreises (10 Punkte)}
\addcontentsline{toc}{subsection}{Aufgabe~1.6 --- Funktionen des Marktpreises}

Sie sind im Einkauf der \emph{Holzbau Kramer GmbH} tätig, eines Zimmereibetriebs in
Thüringen. Durch einen Bauboom und eine stark gestiegene Exportnachfrage aus Übersee ist
die Nachfrage nach Konstruktionsvollholz sprunghaft gestiegen. Der Marktpreis ist binnen
eines Jahres von 420~€ auf 680~€ je Kubikmeter gestiegen. Die Sägewerke fahren
Sonderschichten, ein Investor plant ein neues Sägewerk in der Region. Mehrere kleine
Zimmereien, die die höheren Einkaufspreise nicht an ihre Kunden weitergeben können,
haben ihren Betrieb aufgegeben.

\begin{enumerate}[label=\alph*)]
  \item \textbf{Erläutern} Sie drei Funktionen des Marktpreises jeweils anhand der
    beschriebenen Situation. \hfill (6~P.)
  \item \textbf{Beschreiben} Sie die Marktsituation, die unmittelbar nach dem
    Nachfrageanstieg beim alten Preis von 420~€ je Kubikmeter bestand, und
    \textbf{erläutern} Sie, wie der Preismechanismus den Markt zu einem neuen
    Gleichgewicht führt. \hfill (4~P.)
\end{enumerate}

% ---------------------------------------------------------------------------
\subsection*{Aufgabe~1.7 --- Monopolistischer Absatzbereich im unvollkommenen Polypol (10 Punkte)}
\addcontentsline{toc}{subsection}{Aufgabe~1.7 --- Monopolistischer Absatzbereich im unvollkommenen Polypol}

Sie sind im Produktmanagement der \emph{Senfmanufaktur Brinkhaus GmbH} tätig. Auf dem
deutschen Senfmarkt konkurrieren viele Hersteller mit zahlreichen Marken um viele
Verbraucher. Die Brinkhaus GmbH verkauft ihren „Altmühltaler Grobkornsenf`` in einem
Steingut-Glas mit traditionellem Etikett für 2,60~€ je Glas. Eine Marktforschungsstudie
hat ergeben, dass die Stammkundschaft bei Preisen zwischen 2,20~€ und 3,10~€ je Glas
weitgehend treu bleibt.

\begin{enumerate}[label=\alph*)]
  \item \textbf{Bestimmen} Sie die vorliegende Marktform unter Berücksichtigung der
    Marktvollkommenheit und \textbf{erklären} Sie, wodurch die Brinkhaus GmbH einen
    preispolitischen Spielraum besitzt. \hfill (4~P.)
  \item \textbf{Beurteilen} Sie die Auswirkungen auf den Absatz, wenn die Brinkhaus GmbH
    den Preis (1) auf 2,95~€ bzw. (2) auf 3,40~€ je Glas anhebt. \hfill (4~P.)
  \item \textbf{Nennen} Sie zwei Maßnahmen, mit denen die Brinkhaus GmbH ihren
    preispolitischen Spielraum vergrößern kann. \hfill (2~P.)
\end{enumerate}

% ---------------------------------------------------------------------------
\subsection*{Aufgabe~1.8 --- Preisbildung im Oligopol und im Monopol (12 Punkte)}
\addcontentsline{toc}{subsection}{Aufgabe~1.8 --- Preisbildung im Oligopol und im Monopol}

Sie sind Mitarbeiterin bzw. Mitarbeiter der \emph{Baustoffhandel Lindner GmbH} in
Niederbayern. Zement wird in Ihrer Region von nur drei Zementwerken angeboten; Abnehmer
sind mehrere Hundert Bauunternehmen, Betonwerke und Baustoffhändler. Wegen der hohen
Transportkosten lohnt sich der Bezug aus weiter entfernten Werken kaum.

\begin{enumerate}[label=\alph*)]
  \item \textbf{Bestimmen} Sie die Marktform und \textbf{begründen} Sie Ihre Antwort.
    \hfill (2~P.)
  \item \textbf{Erläutern} Sie zwei für diese Marktform typische Verhaltensweisen der
    Anbieter in der Preispolitik. \hfill (4~P.)
  \item Eines der Werke erwägt, seinen Zementpreis um 10~\% zu senken, um Marktanteile
    zu gewinnen. \textbf{Beurteilen} Sie die Erfolgsaussichten dieser Strategie.
    \hfill (3~P.)
  \item Angenommen, eines der Werke übernimmt die beiden anderen. \textbf{Erläutern} Sie,
    wie sich dadurch seine Preissetzungsmöglichkeiten verändern, und \textbf{nennen} Sie
    zwei Grenzen seiner Preissetzungsmacht. \hfill (3~P.)
\end{enumerate}

\clearpage
\section*{Aufgaben zu Lektion 2: Wettbewerbspolitik, Staatseingriffe \& VGR}
\addcontentsline{toc}{section}{Aufgaben zu Lektion 2: Wettbewerbspolitik, Staatseingriffe \& VGR}
% ===========================================================================
% Gesamtskript VWL --- Übungsaufgaben Lektion 2
% Wettbewerb, Staatseingriffe & VGR (Rahmenplan DIHK 1.1)
% Alle Situationen, Unternehmen und Zahlen sind frei erfunden.
% ===========================================================================

\subsection*{Aufgabe~2.1 --- UWG oder GWB? (12 Punkte)}
\addcontentsline{toc}{subsection}{Aufgabe~2.1 --- UWG oder GWB?}

Sie sind Mitarbeiterin bzw. Mitarbeiter im Vertrieb der \emph{Rheinquell Getränke GmbH},
einem regionalen Hersteller von Mineralwasser und Erfrischungsgetränken. In der
wöchentlichen Vertriebsrunde werden drei Vorgänge besprochen:

\begin{enumerate}[label=(\arabic*)]
  \item Ein Wettbewerber wirbt auf seinen Etiketten mit dem Hinweis „Testsieger 2026``,
    obwohl es den genannten Test nie gegeben hat.
  \item Ein anderer Wettbewerber lässt Privathaushalte ohne deren vorherige Einwilligung
    anrufen, um ihnen ein Getränke-Abo zu verkaufen.
  \item Der Vertriebsleiter berichtet, dass auf einem Branchentreffen zwei Konkurrenten
    vorgeschlagen haben, die Listenpreise für die Gastronomie zum 1.~Januar gemeinsam
    um 8~\% anzuheben. Er findet die Idee „wirtschaftlich vernünftig``.
\end{enumerate}

\begin{enumerate}[label=\alph*)]
  \item \textbf{Unterscheiden} Sie das Gesetz gegen den unlauteren Wettbewerb (UWG) und
    das Gesetz gegen Wettbewerbsbeschränkungen (GWB) hinsichtlich ihres Schutzzwecks.
    \hfill (4~P.)
  \item \textbf{Ordnen} Sie die Vorgänge (1) bis (3) jeweils dem zutreffenden Gesetz zu und
    \textbf{begründen} Sie Ihre Zuordnung. \hfill (6~P.)
  \item \textbf{Nennen} Sie je zwei mögliche Rechtsfolgen, mit denen Verstöße gegen das UWG
    bzw. gegen das GWB geahndet werden können. \hfill (2~P.)
\end{enumerate}

\subsection*{Aufgabe~2.2 --- Kartellverbot und Ausnahmen (10 Punkte)}
\addcontentsline{toc}{subsection}{Aufgabe~2.2 --- Kartellverbot und Ausnahmen}

Sie sind Assistentin bzw. Assistent der Geschäftsleitung der \emph{Elbtal Tiefbau GmbH}
(85~Beschäftigte) in Sachsen. Die Geschäftsführerin bittet Sie, vier Vorhaben
kartellrechtlich einzuschätzen:

\begin{enumerate}[label=(\arabic*)]
  \item Vor einer Ausschreibung des Landkreises für die Sanierung einer Kreisstraße
    sollen Elbtal und zwei Mitbewerber vereinbaren, dass diesmal Elbtal das günstigste
    Angebot abgibt; bei der nächsten Ausschreibung ist ein anderer Mitbewerber „an der Reihe``.
  \item Elbtal will sich mit vier anderen kleinen Tiefbaubetrieben zu einer
    Einkaufsgemeinschaft für Baumaschinen und Baustoffe zusammenschließen, um gegenüber
    den großen Baukonzernen bessere Einkaufskonditionen zu erzielen. Ihre Preise und
    Angebote gestalten alle Betriebe weiterhin selbst.
  \item Elbtal und ein Wettbewerber wollen gemeinsam ein CO$_2$-reduziertes
    Asphaltverfahren entwickeln. Die Ergebnisse dürfen beide nutzen; ihre Preise legen
    sie getrennt fest. Ihr gemeinsamer Marktanteil liegt bei rund 12~\%.
  \item Ein Wettbewerber aus Brandenburg schlägt vor, dass Elbtal künftig keine Angebote
    mehr in Brandenburg abgibt; im Gegenzug verzichtet er auf Aufträge in Sachsen.
\end{enumerate}

\begin{enumerate}[label=\alph*)]
  \item \textbf{Grenzen} Sie das Kartell von der Fusion hinsichtlich der rechtlichen und
    wirtschaftlichen Selbstständigkeit der beteiligten Unternehmen \textbf{ab}. \hfill (2~P.)
  \item \textbf{Beurteilen} Sie die Vorhaben (1) bis (4) jeweils unter Angabe der Kartellart
    bzw. der einschlägigen Ausnahme, ob sie zulässig sind. \hfill (8~P.)
\end{enumerate}

\subsection*{Aufgabe~2.3 --- Fusionskontrolle und Missbrauchsaufsicht (12 Punkte)}
\addcontentsline{toc}{subsection}{Aufgabe~2.3 --- Fusionskontrolle und Missbrauchsaufsicht}

Sie arbeiten in der Strategieabteilung der \emph{Nordlicht Molkerei AG}. Der Vorstand plant,
sämtliche Anteile der \emph{Hansemilch GmbH} zu übernehmen. Für das letzte Geschäftsjahr
liegen folgende Umsatzerlöse vor:

\begin{center}
\begin{tabular}{l r r}
\toprule
\rowcolor{bfwTabHead} \textbf{Unternehmen} & \textbf{Umsatz weltweit} & \textbf{davon im Inland}\\
\midrule
Nordlicht Molkerei AG & 430~Mio.~€ & 390~Mio.~€\\
Hansemilch GmbH & 85~Mio.~€ & 22~Mio.~€\\
\bottomrule
\end{tabular}
\end{center}

\noindent Auf dem norddeutschen Markt für Frischmilchprodukte hätten beide Unternehmen
nach dem Zusammenschluss einen gemeinsamen Marktanteil von rund 43~\%.

\begin{enumerate}[label=\alph*)]
  \item \textbf{Prüfen} Sie anhand der Umsatzschwellen des \S\,35 GWB, ob der Zusammenschluss
    beim Bundeskartellamt anzumelden ist. \textbf{Prüfen} Sie außerdem, ob sich das Ergebnis
    ändert, wenn die Hansemilch GmbH im Inland nur 15~Mio.~€ umsetzen würde (übrige Werte
    unverändert, keine weiteren Besonderheiten). \hfill (5~P.)
  \item \textbf{Erläutern} Sie, wie das Bundeskartellamt nach der Anmeldung vorgeht und
    unter welcher Voraussetzung es den Zusammenschluss untersagen kann. Gehen Sie dabei auf
    den Marktanteil von 43~\% ein. \hfill (3~P.)
  \item Das Bundeskartellamt untersagt die Fusion. \textbf{Erklären} Sie, welche Möglichkeit
    Nordlicht dennoch hat, den Zusammenschluss zu verwirklichen. \hfill (2~P.)
  \item Angenommen, die Fusion kommt zustande und Nordlicht senkt anschließend den
    Milchabnahmepreis für die von ihr abhängigen Landwirte weit unter das übliche
    Marktniveau --- mit der Drohung, andernfalls keine Milch mehr abzunehmen.
    \textbf{Ordnen} Sie dieses Verhalten der zutreffenden Missbrauchsform \textbf{zu} und
    \textbf{begründen} Sie. \hfill (2~P.)
\end{enumerate}

\subsection*{Aufgabe~2.4 --- Höchstpreis auf dem Pelletmarkt (12 Punkte)}
\addcontentsline{toc}{subsection}{Aufgabe~2.4 --- Höchstpreis auf dem Pelletmarkt}

Sie sind Mitarbeiterin bzw. Mitarbeiter im Controlling der \emph{Pelletwerk Sauerland GmbH}.
Nach stark gestiegenen Energiepreisen diskutiert die Regierung einen gesetzlichen
Höchstpreis für Holzpellets. Für den (vereinfachten) Gesamtmarkt gelten folgende
Funktionen ($p$ = Preis in € je Tonne, $x$ = Menge in 1.000~Tonnen):
\[
\text{Nachfrage: } x_N = 900 - 1{,}5\,p \qquad\qquad \text{Angebot: } x_A = -300 + 2{,}5\,p
\]

\begin{enumerate}[label=\alph*)]
  \item \textbf{Berechnen} Sie Gleichgewichtspreis und Gleichgewichtsmenge. \hfill (3~P.)
  \item Der Höchstpreis wird auf 240~€/t festgelegt. \textbf{Berechnen} Sie die angebotene
    und die nachgefragte Menge und \textbf{bestimmen} Sie Art und Umfang der entstehenden
    Marktsituation. \textbf{Begründen} Sie außerdem, ob ein Höchstpreis von 320~€/t eine
    Wirkung hätte. \hfill (4~P.)
  \item Es bildet sich ein Schwarzmarkt. \textbf{Berechnen} Sie den Preis, den die
    Nachfrager für die tatsächlich angebotene Menge höchstens zu zahlen bereit wären, und
    \textbf{erläutern} Sie zwei Maßnahmen, die der Staat flankierend ergreifen müsste,
    damit der Höchstpreis seine Wirkung entfalten kann. \hfill (3~P.)
  \item \textbf{Berechnen} Sie, wie sich der Gesamtumsatz aller Anbieter durch den
    Höchstpreis absolut und prozentual verändert. \hfill (2~P.)
\end{enumerate}

\subsection*{Aufgabe~2.5 --- Marktkonforme und marktkonträre Eingriffe (12 Punkte)}
\addcontentsline{toc}{subsection}{Aufgabe~2.5 --- Marktkonforme und marktkonträre Eingriffe}

Sie arbeiten in der Stabsstelle Energiewirtschaft der \emph{Stadtwerke Lahntal GmbH}. Um
private Haushalte bei den Stromkosten zu entlasten, werden politisch vier Vorschläge
diskutiert:

\begin{enumerate}[label=(\arabic*)]
  \item eine gesetzliche Obergrenze von 30~Cent je kWh für Haushaltsstrom, die kein
    Versorger überschreiten darf --- ohne finanziellen Ausgleich für die Versorger,
  \item eine Senkung der Stromsteuer auf das EU-rechtlich zulässige Mindestmaß,
  \item ein einmaliger Energiekostenzuschuss für Haushalte mit geringem Einkommen,
  \item eine staatliche Investitionszulage für Hersteller von Wärmepumpen.
\end{enumerate}

\begin{enumerate}[label=\alph*)]
  \item \textbf{Ordnen} Sie die Vorschläge (1) bis (4) den marktkonformen bzw.
    marktkonträren Eingriffen \textbf{zu} und \textbf{begründen} Sie jeweils kurz.
    \hfill (4~P.)
  \item \textbf{Erläutern} Sie die Wirkung von Vorschlag (2) auf Angebot,
    Gleichgewichtspreis und Gleichgewichtsmenge am Strommarkt. \hfill (3~P.)
  \item Die Obergrenze aus Vorschlag (1) liegt unter dem Gleichgewichtspreis.
    \textbf{Erläutern} Sie drei mögliche Folgen für die Stadtwerke Lahntal bzw. den
    Strommarkt. \hfill (3~P.)
  \item \textbf{Beurteilen} Sie aus ordnungspolitischer Sicht, ob Vorschlag (3) oder
    Vorschlag (1) vorzuziehen ist. \hfill (2~P.)
\end{enumerate}

\subsection*{Aufgabe~2.6 --- VGR: Vom BIP zum Volkseinkommen (16 Punkte)}
\addcontentsline{toc}{subsection}{Aufgabe~2.6 --- VGR: Vom BIP zum Volkseinkommen}

Sie sind Mitarbeiterin bzw. Mitarbeiter der Research-Abteilung der \emph{Nordbank Consult
GmbH} und bereiten für einen Kundenvortrag die Daten der (fiktiven) Volkswirtschaft
\emph{Aldoria} auf. Für das abgelaufene Jahr liegen folgende Werte vor (in Mrd.~€):

\begin{center}
\begin{tabular}{l r}
\toprule
\rowcolor{bfwTabHead} \textbf{Größe} & \textbf{Mrd.~€}\\
\midrule
Private Konsumausgaben & 2.180\\
Konsumausgaben des Staates & 860\\
Bruttoinvestitionen & 740\\
Exporte von Waren und Dienstleistungen & 1.950\\
Importe von Waren und Dienstleistungen & 1.730\\
Primäreinkommen, die Inländer aus der übrigen Welt erhalten & 210\\
Primäreinkommen, die an die übrige Welt (Nichtinländer) fließen & 145\\
Abschreibungen & 690\\
Produktions- und Importabgaben & 420\\
Subventionen & 70\\
Arbeitnehmerentgelt & 1.936\\
\bottomrule
\end{tabular}
\end{center}

\begin{enumerate}[label=\alph*)]
  \item \textbf{Berechnen} Sie den Außenbeitrag und das Bruttoinlandsprodukt (BIP) über
    die Verwendungsrechnung. \hfill (4~P.)
  \item \textbf{Berechnen} Sie das Bruttonationaleinkommen (BNE) und \textbf{erläutern} Sie,
    warum es in Aldoria größer ist als das BIP. \hfill (3~P.)
  \item \textbf{Berechnen} Sie das Nettonationaleinkommen und das Volkseinkommen.
    \hfill (4~P.)
  \item \textbf{Berechnen} Sie die Unternehmens- und Vermögenseinkommen sowie die Lohn- und
    die Gewinnquote. \hfill (3~P.)
  \item Im Folgejahr steigt die Lohnquote auf 66~\%, obwohl die Arbeitnehmerentgelte
    unverändert geblieben sind. \textbf{Erläutern} Sie, wie dies zu erklären ist.
    \hfill (2~P.)
\end{enumerate}

\subsection*{Aufgabe~2.7 --- Nominales und reales Wachstum (12 Punkte)}
\addcontentsline{toc}{subsection}{Aufgabe~2.7 --- Nominales und reales Wachstum}

Sie sind im Controlling der \emph{Weserland Logistik GmbH} tätig. Die Geschäftsführung
wundert sich: „Das BIP steigt jedes Jahr --- warum sprechen die Medien trotzdem von einer
schrumpfenden Wirtschaft?`` Sie erhalten folgende Daten (Jahr~1 = Basisjahr):

\begin{center}
\begin{tabular}{l r r r}
\toprule
\rowcolor{bfwTabHead} & \textbf{Jahr 1} & \textbf{Jahr 2} & \textbf{Jahr 3}\\
\midrule
Nominales BIP (Mrd.~€) & 4.000,0 & 4.202,4 & 4.263,0\\
BIP-Deflator (Basisjahr = 100) & 100,0 & 103,0 & 105,0\\
\bottomrule
\end{tabular}
\end{center}

\begin{enumerate}[label=\alph*)]
  \item \textbf{Berechnen} Sie das reale BIP für die Jahre~2 und~3. \hfill (3~P.)
  \item \textbf{Berechnen} Sie die nominalen und die realen Wachstumsraten für die
    Jahre~2 und~3 (auf zwei Nachkommastellen). \hfill (4~P.)
  \item \textbf{Berechnen} Sie die Preissteigerung (gemessen am Deflator) von Jahr~2 auf
    Jahr~3 und \textbf{erläutern} Sie damit die Frage der Geschäftsführung. \hfill (3~P.)
  \item \textbf{Unterscheiden} Sie den BIP-Deflator vom Verbraucherpreisindex (VPI).
    \hfill (2~P.)
\end{enumerate}

\subsection*{Aufgabe~2.8 --- Magisches Viereck: Lage beurteilen (14 Punkte)}
\addcontentsline{toc}{subsection}{Aufgabe~2.8 --- Magisches Viereck: Lage beurteilen}

Sie sind in der Geschäftsleitungsassistenz der \emph{Allgäu Präzisionstechnik GmbH}, eines
stark exportorientierten Maschinenbauers, tätig. Für die Strategieklausur sollen Sie die
gesamtwirtschaftliche Lage eines Landes beurteilen, in dem das Unternehmen einen großen
Teil seines Umsatzes erzielt. Es liegen folgende Jahreswerte vor:

\begin{center}
\begin{tabular}{l r}
\toprule
\rowcolor{bfwTabHead} \textbf{Kennzahl} & \textbf{Wert}\\
\midrule
Inflationsrate (HVPI) & 4,1~\%\\
Arbeitslosenquote & 3,4~\%\\
Wachstumsrate des realen BIP & 0,2~\%\\
Außenbeitrag in \% des nominalen BIP & 5,5~\%\\
\bottomrule
\end{tabular}
\end{center}

\begin{enumerate}[label=\alph*)]
  \item \textbf{Beurteilen} Sie für jedes der vier Ziele des magischen Vierecks anhand der
    Kennzahlen, ob es erreicht ist. Nennen Sie dabei jeweils den Zielwert. \hfill (8~P.)
  \item Die Regierung will die Inflation durch deutliche Kürzungen der Staatsausgaben
    bekämpfen. \textbf{Erläutern} Sie, welcher Zielkonflikt dadurch entstehen kann.
    \hfill (3~P.)
  \item Der Geschäftsführer meint: „Ein Außenbeitrag von 5,5~\% zeigt doch nur, wie stark
    das Land ist --- je höher, desto besser.`` \textbf{Beurteilen} Sie diese Aussage.
    \hfill (3~P.)
\end{enumerate}

\clearpage
\section*{Aufgaben zu Lektion 3: Konjunktur, Wirtschaftspolitik \& Außenwirtschaft}
\addcontentsline{toc}{section}{Aufgaben zu Lektion 3: Konjunktur, Wirtschaftspolitik \& Außenwirtschaft}
% ============================================================
% Gesamtskript VWL --- Übungsaufgaben Lektion 3
% Konjunktur, Wirtschaftspolitik & Außenwirtschaft (Rahmenplan DIHK 1.1)
% 8 Aufgaben, 100 Punkte --- Lösungen in loesungen-l3.tex
% ============================================================

\subsection*{Aufgabe~3.1 --- Konjunkturlage richtig deuten (11 Punkte)}
\addcontentsline{toc}{subsection}{Aufgabe~3.1 --- Konjunkturlage richtig deuten}

Sie sind Mitarbeiter/-in in der Unternehmensplanung der \emph{Kranich Fördertechnik GmbH},
einem mittelständischen Hersteller von Förderbändern und Hebeanlagen. Die Geschäftsführung
will wissen, wie sie die Produktions- und Personalplanung für das kommende Jahr ausrichten
soll. Ihnen liegen folgende aktuelle Wirtschaftsdaten für Deutschland vor:

\begin{enumerate}[label=(\arabic*)]
  \item Der ifo-Geschäftsklimaindex für das Verarbeitende Gewerbe ist zum dritten Mal in
    Folge gestiegen.
  \item Die Auftragseingänge in der Industrie liegen 3,5~\% über dem Vorquartal.
  \item Die Kapazitätsauslastung der Industrie beträgt 76~\% (langjähriger Durchschnitt:
    83~\%).
  \item Die Arbeitslosenquote ist im letzten Monat erneut leicht gestiegen.
  \item Die Zahl der Unternehmensinsolvenzen liegt auf dem höchsten Stand seit mehreren
    Jahren.
\end{enumerate}

\begin{enumerate}[label=\alph*)]
  \item \textbf{Ordnen} Sie die fünf Daten jeweils den Früh-, Präsenz- oder
    Spätindikatoren \textbf{zu}. \hfill (5~P.)
  \item \textbf{Bestimmen} Sie anhand der Daten die aktuelle Konjunkturphase und
    \textbf{begründen} Sie Ihre Einschätzung. \hfill (3~P.)
  \item \textbf{Erläutern} Sie zwei Schlussfolgerungen, die die Kranich Fördertechnik GmbH
    aus Ihrer Einschätzung für ihre Planung ziehen sollte. \hfill (3~P.)
\end{enumerate}

\subsection*{Aufgabe~3.2 --- Inflation, Realzins und Lohnabschluss (14 Punkte)}
\addcontentsline{toc}{subsection}{Aufgabe~3.2 --- Inflation, Realzins und Lohnabschluss}

Sie sind Mitarbeiter/-in im Rechnungswesen der \emph{Holtkamp Logistik GmbH}. Für eine
Präsentation in der Geschäftsleitung sollen Sie die Auswirkungen der Preisentwicklung
aufbereiten. Der Verbraucherpreisindex (VPI, Basisjahr 2020~$=100$) lag im September des
Vorjahres bei 124,0 Punkten und im September dieses Jahres bei 127,1 Punkten
(angenommene Werte).

\begin{enumerate}[label=\alph*)]
  \item \textbf{Berechnen} Sie die Inflationsrate für diesen Zeitraum. \hfill (3~P.)
  \item Die Holtkamp Logistik GmbH hat vor einem Jahr 250.000~€ als Festgeld für ein Jahr
    zu einem Zinssatz von 3,4~\% p.\,a. angelegt. \textbf{Berechnen} Sie den Zinsertrag in
    Euro sowie (näherungsweise) den Realzins, und \textbf{beurteilen} Sie das Ergebnis
    kurz. Gehen Sie davon aus, dass die Inflationsrate aus a) für die gesamte
    Anlagedauer gilt; Steuern bleiben unberücksichtigt. \hfill (4~P.)
  \item Das monatliche Bruttogehalt eines Disponenten wurde im selben Zeitraum von
    3.200~€ auf 3.264~€ angehoben. \textbf{Berechnen} Sie die nominale Gehaltssteigerung
    in Prozent und \textbf{ermitteln} Sie näherungsweise die Veränderung des Reallohns.
    \textbf{Interpretieren} Sie das Ergebnis. \hfill (3~P.)
  \item In der anstehenden Tarifrunde fordert die Gewerkschaft eine Lohnerhöhung von
    5,5~\%. Die Arbeitsproduktivität ist im letzten Jahr um 1,3~\% gestiegen.
    \textbf{Berechnen} Sie den kostenniveauneutralen Verteilungsspielraum (Zielinflation
    der EZB) und \textbf{beurteilen} Sie die Forderung aus Sicht der Arbeitgeber und aus
    Sicht der Gewerkschaft. \hfill (4~P.)
\end{enumerate}

\subsection*{Aufgabe~3.3 --- Wechselkursänderung im Exportgeschäft (14 Punkte)}
\addcontentsline{toc}{subsection}{Aufgabe~3.3 --- Wechselkursänderung im Exportgeschäft}

Sie sind Mitarbeiter/-in im Vertrieb der \emph{Nordlicht Optik GmbH}, Kiel. Das
Unternehmen verkauft Laser-Messgeräte in die USA und nach Österreich. Ein Messgerät wird
zum Preis von 45.000~€ kalkuliert. Für die Fertigung bezieht Nordlicht Optik Sensoren von
einem Lieferanten in den USA; eine Jahreslieferung kostet 32.400~USD. Zu Jahresbeginn gilt der
Kurs 1~€ $=$ 1,08~USD (Mengennotierung).

\begin{enumerate}[label=\alph*)]
  \item \textbf{Berechnen} Sie den Preis eines Messgeräts in US-Dollar und den
    Euro-Betrag der Sensorlieferung zum Kurs von 1,08~USD. \hfill (2~P.)
  \item Zum Jahresende notiert der Kurs bei 1~€ $=$ 1,20~USD. \textbf{Stellen} Sie
    \textbf{fest}, ob der Euro auf- oder abgewertet hat. \textbf{Berechnen} Sie den neuen
    US-Dollar-Preis des Messgeräts (absolut und prozentuale Veränderung) sowie den neuen
    Euro-Betrag der Sensorlieferung. \hfill (5~P.)
  \item Ein wichtiger US-Kunde akzeptiert keine Preiserhöhung. Nordlicht Optik hält
    deshalb den bisherigen Dollarpreis aus a). \textbf{Berechnen} Sie den Euro-Erlös je
    Messgerät beim neuen Kurs und die Erlösminderung in Euro und Prozent. \hfill (2~P.)
  \item \textbf{Nennen} Sie zwei Möglichkeiten, wie sich Nordlicht Optik künftig gegen
    solche Wechselkursrisiken absichern kann. \hfill (2~P.)
  \item \textbf{Erläutern} Sie, warum bei den Lieferungen nach Österreich kein
    Wechselkursrisiko besteht, und \textbf{nennen} Sie zwei weitere Vorteile, die der
    Nordlicht Optik GmbH aus dem Europäischen Binnenmarkt bzw. der Währungsunion
    entstehen. \hfill (3~P.)
\end{enumerate}

\subsection*{Aufgabe~3.4 --- Regionaler Arbeitsmarkt (13 Punkte)}
\addcontentsline{toc}{subsection}{Aufgabe~3.4 --- Regionaler Arbeitsmarkt}

Sie sind Mitarbeiter/-in der \emph{Wirtschaftsförderungsgesellschaft Kreis Hollenbeck mbH}
und bereiten den jährlichen Arbeitsmarktbericht vor. Es liegen folgende Daten vor:

\begin{center}
\begin{tabular}{lrr}
\toprule
 & \textbf{Vorjahr} & \textbf{Berichtsjahr}\\
\midrule
Erwerbstätige & 236.800 & 234.000\\
registrierte Arbeitslose & 13.200 & 16.000\\
\bottomrule
\end{tabular}
\end{center}

Von den 16.000 Arbeitslosen des Berichtsjahres sind nach einer Auswertung der
Arbeitsagentur 7.400 Personen strukturell bedingt arbeitslos.

\begin{enumerate}[label=\alph*)]
  \item \textbf{Berechnen} Sie die Arbeitslosenquote für beide Jahre und die Veränderung
    in Prozentpunkten. \hfill (4~P.)
  \item \textbf{Ordnen} Sie die folgenden vier Fälle jeweils einer Form der
    Arbeitslosigkeit \textbf{zu} und \textbf{begründen} Sie kurz. \hfill (4~P.)
    \begin{enumerate}[label=(\arabic*)]
      \item Ein Maurer wird jedes Jahr von Dezember bis März nicht beschäftigt.
      \item Eine Industriemechanikerin eines Autozulieferers wird entlassen, weil die
        gesamtwirtschaftliche Nachfrage eingebrochen ist und die Aufträge fehlen.
      \item Ein Drucker verliert seine Stelle, weil seine Zeitungsdruckerei wegen der
        Digitalisierung dauerhaft schließt; sein Beruf wird in der Region kaum noch
        nachgefragt.
      \item Eine Bürokauffrau hat gekündigt, um in eine andere Stadt zu ziehen, und hat
        dort nach sechs Wochen Suche bereits eine Zusage für eine neue Stelle.
    \end{enumerate}
  \item \textbf{Berechnen} Sie den Anteil der strukturellen Arbeitslosigkeit an der
    gesamten Arbeitslosigkeit und \textbf{erläutern} Sie, warum ein hoher Anteil für den
    Landkreis besonders problematisch ist. \hfill (2~P.)
  \item \textbf{Schlagen} Sie je eine geeignete Maßnahme gegen die strukturelle und gegen
    die konjunkturelle Arbeitslosigkeit \textbf{vor} und \textbf{ordnen} Sie diese der
    aktiven bzw. passiven Arbeitsmarktpolitik \textbf{zu}. \hfill (3~P.)
\end{enumerate}

\subsection*{Aufgabe~3.5 --- Leitzinssenkung der EZB (12 Punkte)}
\addcontentsline{toc}{subsection}{Aufgabe~3.5 --- Leitzinssenkung der EZB}

Sie sind Mitarbeiter/-in in der Finanzierungsabteilung der \emph{Wohnbau Lindner GmbH},
eines Bauträgers für Eigentumswohnungen. Die Inflationsrate im Euroraum (HVPI) liegt seit
mehreren Monaten bei rund 1,3~\%; die Prognosen erwarten auch mittelfristig einen Wert
deutlich unter 2~\%. Die Wirtschaft im Euroraum stagniert. Der EZB-Rat beschließt, die drei
Leitzinsen um jeweils 0,25~Prozentpunkte zu senken.

\begin{enumerate}[label=\alph*)]
  \item \textbf{Erläutern} Sie, warum die EZB die Zinsen senkt, obwohl die Preise im
    Euroraum weiterhin steigen. \hfill (3~P.)
  \item \textbf{Stellen} Sie die Wirkungskette der Zinssenkung bis zur Auftragslage der
    Wohnbau Lindner GmbH \textbf{dar} (mindestens fünf Stufen). \hfill (5~P.)
  \item Ein Kollege meint: „Endlich hat die Bundesregierung die Zinsen gesenkt, um die
    Bauwirtschaft zu retten.“ \textbf{Nehmen} Sie zu dieser Aussage \textbf{Stellung}.
    \hfill (2~P.)
  \item \textbf{Erläutern} Sie zwei Gründe, warum die Zinssenkung möglicherweise nicht zu
    mehr Wohnungskäufen führt. \hfill (2~P.)
\end{enumerate}

\subsection*{Aufgabe~3.6 --- Konjunkturprogramm und Schuldenbremse (11 Punkte)}
\addcontentsline{toc}{subsection}{Aufgabe~3.6 --- Konjunkturprogramm und Schuldenbremse}

Sie sind Assistent/-in der Geschäftsführung der \emph{Brenner Tiefbau GmbH}. Deutschland
befindet sich in einer Rezession. In der Bundesregierung werden folgende Maßnahmen
diskutiert:

\begin{enumerate}[label=(\arabic*)]
  \item Bereits geplante Investitionen in Brücken und Schienenwege werden um zwei Jahre
    vorgezogen.
  \item Der Regelsteuersatz der Umsatzsteuer wird angehoben.
  \item Unternehmen erhalten befristet eine Sonderabschreibung auf neue Maschinen.
  \item Subventionen werden gekürzt; die Einsparungen werden einer Rücklage zugeführt.
\end{enumerate}

\begin{enumerate}[label=\alph*)]
  \item \textbf{Ordnen} Sie die vier Maßnahmen als expansiv oder restriktiv \textbf{zu}
    und \textbf{geben} Sie \textbf{an}, welche davon in der aktuellen Lage
    konjunkturgerecht sind. \hfill (4~P.)
  \item Das nominale BIP beträgt 4.400~Mrd.~€. Der Bund plant für das Konjunkturprogramm
    eine Nettokreditaufnahme von 41~Mrd.~€, davon sind 22~Mrd.~€ als konjunkturbedingt
    anerkannt (Konjunkturkomponente). Sonderregelungen (z.\,B. für Verteidigungsausgaben
    oder Sondervermögen) bleiben außer Betracht. \textbf{Prüfen} Sie rechnerisch, ob die
    Schuldenbremse des Grundgesetzes eingehalten wird. \hfill (4~P.)
  \item \textbf{Erläutern} Sie, warum das Investitionsprogramm (1) trotz guter Absicht
    prozyklisch wirken kann, und \textbf{erklären} Sie den Begriff Crowding-out.
    \hfill (3~P.)
\end{enumerate}

\subsection*{Aufgabe~3.7 --- Zahlungsbilanz und Handelskonflikt (14 Punkte)}
\addcontentsline{toc}{subsection}{Aufgabe~3.7 --- Zahlungsbilanz und Handelskonflikt}

Sie sind Mitarbeiter/-in der Exportabteilung der \emph{Weinmann Landtechnik GmbH}, die
Traktoren und Erntemaschinen weltweit verkauft. Für eine Marktanalyse liegen Ihnen
folgende (gerundete, angenommene) Zahlungsbilanzdaten eines Jahres für Deutschland vor
(in Mrd.~€):

\begin{center}
\begin{tabular}{lr}
\toprule
\textbf{Position} & \textbf{Mrd.~€}\\
\midrule
Warenexporte & 1.480\\
Warenimporte & 1.290\\
Einnahmen aus Dienstleistungen & 395\\
Ausgaben für Dienstleistungen & 455\\
Saldo der Primäreinkommen & $+115$\\
Saldo der Sekundäreinkommen & $-72$\\
\bottomrule
\end{tabular}
\end{center}

\begin{enumerate}[label=\alph*)]
  \item \textbf{Berechnen} Sie den Saldo der Handelsbilanz, der Dienstleistungsbilanz und
    der Leistungsbilanz. \hfill (4~P.)
  \item \textbf{Erläutern} Sie, in welcher Teilbilanz sich der Leistungsbilanzsaldo
    spiegelbildlich niederschlägt und warum die Zahlungsbilanz formal stets ausgeglichen
    ist. \hfill (2~P.)
  \item Der (fiktive) Staat \emph{Westland}, ein wichtiger Absatzmarkt der Weinmann
    Landtechnik GmbH, kündigt an: (1)~einen Einfuhrzoll von 20~\% auf Landmaschinen aus
    der EU, (2)~eine Höchstmenge von 500 importierten Traktoren pro Jahr, (3)~eine neue,
    aufwendige Sicherheitszertifizierung, die nur von Prüfstellen in Westland
    durchgeführt werden darf, (4)~Exportsubventionen für die eigenen
    Landmaschinenhersteller. \textbf{Unterscheiden} Sie tarifäre und nichttarifäre
    Handelshemmnisse und \textbf{ordnen} Sie die vier Maßnahmen \textbf{zu}. \hfill (4~P.)
  \item In Deutschland wird gefordert, die EU solle mit Gegenzöllen auf Waren aus Westland
    antworten. \textbf{Beurteilen} Sie diese Forderung, indem Sie je zwei Argumente dafür
    und dagegen anführen. \hfill (4~P.)
\end{enumerate}

\subsection*{Aufgabe~3.8 --- Emissionshandel im Unternehmen (11 Punkte)}
\addcontentsline{toc}{subsection}{Aufgabe~3.8 --- Emissionshandel im Unternehmen}

Sie sind Mitarbeiter/-in im Controlling der \emph{Hagedorn Glaswerke GmbH}. Das Werk
nimmt am Emissionshandel der EU teil. Für das kommende Jahr gelten folgende Planwerte:

\begin{itemize}
  \item kostenlos zugeteilte Emissionsberechtigungen: 40.000 (1 Berechtigung berechtigt
    zum Ausstoß von 1~t CO\textsubscript{2}),
  \item erwartete Emissionen ohne Gegenmaßnahmen: 46.000~t CO\textsubscript{2},
  \item angenommener Marktpreis: 80~€ je Berechtigung (konstant),
  \item Investitionsalternative: Eine Wärmerückgewinnungsanlage senkt die Emissionen um
    9.000~t CO\textsubscript{2} pro Jahr; sie verursacht jährliche Kosten (Abschreibung,
    Zinsen, Betrieb) von 540.000~€.
\end{itemize}

\begin{enumerate}[label=\alph*)]
  \item \textbf{Erläutern} Sie die Funktionsweise des Emissionshandels und \textbf{nennen}
    Sie das umweltpolitische Prinzip, auf dem er beruht. \hfill (3~P.)
  \item \textbf{Berechnen} Sie die jährliche Kostenbelastung aus dem Emissionshandel ohne
    und mit Wärmerückgewinnungsanlage und \textbf{entscheiden} Sie begründet, welche
    Alternative wirtschaftlich günstiger ist. \hfill (5~P.)
  \item \textbf{Beurteilen} Sie den Emissionshandel im Vergleich zu einer staatlichen
    Auflage (fester Emissionsgrenzwert je Werk), indem Sie zwei Vorteile und einen
    Nachteil des Emissionshandels darlegen. \hfill (3~P.)
\end{enumerate}



\clearpage
\addcontentsline{toc}{part}{Teil C\quad Lösungen}
\vspace*{5cm}
{\bfwlabelfont\color{bfwPrimary}TEIL C\par}
\vspace{0.5em}
{\Huge\bfseries\color{bfwPrimaryDark}Lösungen\par}
\vspace{0.6em}
{\color{bfwPrimary}\rule{\linewidth}{2pt}}\par
\vspace{0.6em}
{\large\color{bfwInkSoft}Musterlösungen mit Punktverteilung und Rechenwegen\par}
\clearpage

\clearpage
\section*{Lösungen zu Lektion 1: Grundbegriffe \& Preisbildung am Markt}
\addcontentsline{toc}{section}{Lösungen zu Lektion 1: Grundbegriffe \& Preisbildung am Markt}
% ===========================================================================
% Gesamtskript VWL -- Musterlösungen Lektion 1: Grundbegriffe & Preisbildung
% (Fragment, wird per \input eingebunden; Aufgaben in aufgaben-l1.tex)
% ===========================================================================

% ---------------------------------------------------------------------------
\begin{loesung}[title={Lösung Aufgabe~1.1 --- Wirtschaften und ökonomisches Prinzip}]
\textbf{a) Bedürfnis, Bedarf, Nachfrage} \hfill (3~P.: je Begriff 1~P.)
\begin{itemize}
  \item \emph{Bedürfnis:} Frau Jansen empfindet einen Mangel --- sie möchte ihre Kinder
    ohne Auto zur Kita bringen und wünscht sich dafür ein geeignetes Verkehrsmittel.
  \item \emph{Bedarf:} Da sie 4.500~€ angespart hat, ist ihr Bedürfnis mit
    \emph{Kaufkraft} ausgestattet; aus dem Bedürfnis wird Bedarf.
  \item \emph{Nachfrage:} Mit der Bestellung im Onlineshop wird der Bedarf am Markt
    \emph{wirksam} --- es entsteht Nachfrage nach einem Lastenrad.
\end{itemize}

\textbf{b) Zuordnung zum ökonomischen Prinzip} \hfill (6~P.: je Aussage 1~P. Zuordnung +
1~P. Begründung)
\begin{itemize}
  \item \emph{Aussage (1) --- Minimalprinzip:} Das Ziel ist vorgegeben (2.400 Rahmen); es
    soll mit möglichst geringem Mitteleinsatz (niedrigste Beschaffungskosten) erreicht
    werden.
  \item \emph{Aussage (2) --- Maximalprinzip:} Die Mittel sind vorgegeben (Messebudget
    18.000~€); mit ihnen soll ein möglichst großer Erfolg (Zahl der Händlerkontakte)
    erzielt werden.
  \item \emph{Aussage (3) --- kein ökonomisches Prinzip:} Mittel und Ziel werden
    \emph{gleichzeitig} optimiert („geringstmögliche Kosten`` \emph{und}
    „größtmögliche Zahl``). Das ist logisch nicht lösbar: Die geringstmöglichen Kosten
    wären null, dann würde nichts produziert. Eine der beiden Größen muss fest vorgegeben
    sein.
\end{itemize}

\textbf{c) Produktionsfaktoren} \hfill (2~P.: je Zuordnung 0,5~P.)

Montageteam $\rightarrow$ \emph{Arbeit}; Gewerbegrundstück $\rightarrow$ \emph{Boden}
(Standort); Lackierroboter $\rightarrow$ \emph{Kapital} (produziertes Produktionsmittel);
patentiertes Klebeverfahren $\rightarrow$ \emph{Wissen (Know-how)}.
\end{loesung}

% ---------------------------------------------------------------------------
\begin{loesung}[title={Lösung Aufgabe~1.2 --- Marktarten und Marktformen}]
\textbf{a) Marktarten} \hfill (4~P.: je Vorgang 1~P.)
\begin{enumerate}[label=(\arabic*)]
  \item Stellenausschreibung $\rightarrow$ \emph{Faktormarkt} (Arbeitsmarkt).
  \item Anmietung der Lagerhalle $\rightarrow$ \emph{Faktormarkt} (Immobilienmarkt,
    Faktor Boden).
  \item Investitionskredit $\rightarrow$ \emph{Geld- und Kapitalmarkt} (Finanzmarkt,
    Kreditmarkt).
  \item Verkauf von Kartoffeln $\rightarrow$ \emph{Gütermarkt} (Konsumgütermarkt).
\end{enumerate}

\textbf{b) Marktformen} \hfill (6~P.: je Situation 1~P. Marktform + 1~P. Begründung)
\begin{enumerate}[label=(\arabic*)]
  \item \emph{Nachfrageoligopol:} Viele Anbieter (rund 1.800 Landwirte) stehen
    \emph{wenigen} Nachfragern (zwei Zuckerfabriken) gegenüber. Die Marktmacht liegt bei
    den Nachfragern.
  \item \emph{Angebotsmonopol:} Ein einziger Anbieter (die Reederei) steht vielen
    Nachfragern (Fahrgästen) gegenüber; Ausweichmöglichkeiten fehlen weitgehend.
  \item \emph{Angebotsoligopol:} Wenige Anbieter (internationale Chemiekonzerne) mit
    großen Marktanteilen stehen vielen Nachfragern (Händler, Landwirte) gegenüber.
\end{enumerate}

\textbf{c) Beurteilung der Aussage} \hfill (2~P.)

Die Aussage ist \emph{falsch}. Viele Anbieter und viele Nachfrager (Polypol) sind nur
\emph{eine} von fünf Bedingungen des vollkommenen Marktes. Auf dem Wochenmarkt fehlen
weitere Bedingungen, z.\,B.: Die Güter sind nicht homogen (Bio- und konventionelle Ware,
unterschiedliche Sorten und Qualität), Kunden haben persönliche Präferenzen
(Stammhändler) und räumliche Präferenzen (Stand am Eingang), und es besteht keine
vollständige Markttransparenz. Es handelt sich daher um einen \emph{unvollkommenen} Markt
(unvollkommenes Polypol).
\end{loesung}

% ---------------------------------------------------------------------------
\begin{loesung}[title={Lösung Aufgabe~1.3 --- Preiselastizität und Umsatzwirkung}]
\textbf{a) Direkte Preiselastizität} \hfill (4~P.: je Produkt 2~P.)

Formel: $\displaystyle \varepsilon = \frac{\text{Mengenänderung in \%}}{\text{Preisänderung in \%}}$
\quad (Änderungen jeweils bezogen auf den Ausgangswert)

\smallskip
\emph{Espresso „Crema Oro``:}
\begin{itemize}
  \item Preisänderung: $\dfrac{15{,}40~\text{€} - 14{,}00~\text{€}}{14{,}00~\text{€}}
    = \dfrac{1{,}40~\text{€}}{14{,}00~\text{€}} = +10\,\%$
  \item Mengenänderung: $\dfrac{6.800 - 8.000}{8.000} = \dfrac{-1.200~\text{Pck.}}{8.000~\text{Pck.}} = -15\,\%$
  \item $\varepsilon = \dfrac{-15\,\%}{+10\,\%} = \mathbf{-1{,}5}$ \quad $\Rightarrow\ |\varepsilon| = 1{,}5$
\end{itemize}

\emph{Filterkaffee „Hausmischung``:}
\begin{itemize}
  \item Preisänderung: $\dfrac{6{,}60~\text{€} - 6{,}00~\text{€}}{6{,}00~\text{€}}
    = \dfrac{0{,}60~\text{€}}{6{,}00~\text{€}} = +10\,\%$
  \item Mengenänderung: $\dfrac{19.200 - 20.000}{20.000} = \dfrac{-800~\text{Pck.}}{20.000~\text{Pck.}} = -4\,\%$
  \item $\varepsilon = \dfrac{-4\,\%}{+10\,\%} = \mathbf{-0{,}4}$ \quad $\Rightarrow\ |\varepsilon| = 0{,}4$
\end{itemize}

\textbf{b) Interpretation} \hfill (4~P.: je Produkt 1~P. Einordnung + 1~P. Begründung)
\begin{itemize}
  \item \emph{Espresso:} $|\varepsilon| = 1{,}5 > 1$ $\Rightarrow$ \emph{elastische}
    Nachfrage. Die Menge sinkt prozentual stärker (15~\%) als der Preis steigt (10~\%).
    Begründung: Premium-Espresso ist kein Grundbedarf; es gibt viele Ausweichmöglichkeiten
    (andere Marken, Discounter-Espresso, Kapselsysteme) bzw. die Kunden können auf das
    Produkt verzichten.
  \item \emph{Filterkaffee:} $|\varepsilon| = 0{,}4 < 1$ $\Rightarrow$ \emph{unelastische}
    Nachfrage. Die Menge sinkt prozentual schwächer (4~\%) als der Preis steigt (10~\%).
    Begründung: Filterkaffee ist für viele Haushalte ein Alltagsgut mit festen
    Konsumgewohnheiten; die Kunden sind an „ihre`` Hausmarke gebunden und reagieren auf
    Preisänderungen kaum.
\end{itemize}
Das negative Vorzeichen ist bei der direkten Preiselastizität normal (Preis steigt,
Menge sinkt); maßgeblich für die Einordnung ist der Betrag.

\textbf{c) Umsatzwirkung} \hfill (4~P.: je Produkt 2~P.)

\begin{center}
\small
\begin{tabular}{l r r r r}
\toprule
\rowcolor{bfwTabHead}
\textbf{Produkt} & \textbf{Umsatz vorher} & \textbf{Umsatz nachher} & \textbf{Änderung in €} & \textbf{Änderung in \%}\\
\midrule
Espresso & $14{,}00 \cdot 8.000 = 112.000$~€ & $15{,}40 \cdot 6.800 = 104.720$~€ & $-7.280$~€ & $-6{,}5\,\%$\\
Filterkaffee & $6{,}00 \cdot 20.000 = 120.000$~€ & $6{,}60 \cdot 19.200 = 126.720$~€ & $+6.720$~€ & $+5{,}6\,\%$\\
\bottomrule
\end{tabular}
\end{center}
Prozentwerte: $-7.280 \div 112.000 = -0{,}065 = -6{,}5\,\%$;
$\ +6.720 \div 120.000 = +0{,}056 = +5{,}6\,\%$.\\
Probe: $1{,}10 \cdot 0{,}85 = 0{,}935$ ($-6{,}5\,\%$) bzw. $1{,}10 \cdot 0{,}96 = 1{,}056$
($+5{,}6\,\%$). Ergebnis: Beim elastischen Gut sinkt, beim unelastischen Gut steigt der
Umsatz.

\textbf{d) Beurteilung des Vorschlags} \hfill (4~P.)
\begin{itemize}
  \item \emph{Espresso:} Die Preiserhöhung ist aus Umsatzsicht \emph{nicht} zu empfehlen,
    da die elastische Nachfrage zu einem Umsatzrückgang von 7.280~€ je Quartal führt.
    Zudem droht der dauerhafte Verlust von Kunden an die Konkurrenz.
  \item \emph{Filterkaffee:} Die Preiserhöhung ist zu empfehlen; der Umsatz steigt um
    6.720~€ je Quartal.
  \item \emph{Gesamtbetrachtung:} Würden beide Erhöhungen umgesetzt, sänke der
    Gesamtumsatz um $7.280 - 6.720 = 560$~€ je Quartal. Sinnvoll ist daher eine
    \emph{differenzierte} Preispolitik (nur Filterkaffee erhöhen).
  \item \emph{Einschränkungen:} Umsatz ist nicht gleich Gewinn --- durch die geringere
    Absatzmenge sinken auch die variablen Kosten; die Entscheidung sollte daher auch über
    den Deckungsbeitrag geprüft werden. Außerdem gelten die Testmarktdaten nur
    \emph{ceteris paribus} und kurzfristig; bundesweit und langfristig (Reaktion der
    Wettbewerber, Gewöhnung der Kunden) können die Elastizitäten abweichen.
\end{itemize}
\end{loesung}

% ---------------------------------------------------------------------------
\begin{loesung}[title={Lösung Aufgabe~1.4 --- Marktgleichgewicht aus Angebots- und Nachfragefunktion}]
\textbf{a) Gleichgewicht} \hfill (4~P.: $p^*$ 2~P., $x^*$ 1~P., Umsatz 1~P.)

Im Gleichgewicht gilt: angebotene Menge = nachgefragte Menge, also $x_A = x_N$.
\[
60\,p - 300 = 1.200 - 40\,p
\quad\Rightarrow\quad 100\,p = 1.500
\quad\Rightarrow\quad p^* = \mathbf{15~\text{€ je Sack}}
\]
Einsetzen: $x_N = 1.200 - 40 \cdot 15 = 1.200 - 600 = \mathbf{600~\text{Säcke pro Tag}}$\\
Probe: $x_A = 60 \cdot 15 - 300 = 900 - 300 = 600$~Säcke (stimmt überein)\\
Umsatz im Gleichgewicht: $15~\text{€/Sack} \cdot 600~\text{Säcke} = \mathbf{9.000~\text{€ pro Tag}}$.

\textbf{b) Überhänge bei 12~€ und 20~€} \hfill (6~P.: je Preis 1~P. Mengen, 1~P. Überhang,
1~P. Marktsituation/Reaktion)

\begin{center}
\small
\begin{tabular}{c c c l}
\toprule
\rowcolor{bfwTabHead}
\textbf{Preis} & \textbf{Nachfrage $x_N$} & \textbf{Angebot $x_A$} & \textbf{Überhang}\\
\midrule
12~€ & $1.200 - 40 \cdot 12 = 720$ & $60 \cdot 12 - 300 = 420$ & Nachfrageüberhang $720 - 420 = 300$~Säcke\\
20~€ & $1.200 - 40 \cdot 20 = 400$ & $60 \cdot 20 - 300 = 900$ & Angebotsüberhang $900 - 400 = 500$~Säcke\\
\bottomrule
\end{tabular}
\end{center}
\begin{itemize}
  \item \emph{Preis 12~€ (unter $p^*$):} Nachfrageüberhang (Angebotsdefizit) von
    300~Säcken pro Tag $\Rightarrow$ \emph{Verkäufermarkt}. Nicht alle Kunden werden
    beliefert; die Anbieter können höhere Preise durchsetzen. Mit steigendem Preis
    weiten die Anbieter ihre Menge aus, die nachgefragte Menge sinkt --- der Preis
    steigt in Richtung 15~€.
  \item \emph{Preis 20~€ (über $p^*$):} Angebotsüberhang (Nachfragedefizit) von
    500~Säcken pro Tag $\Rightarrow$ \emph{Käufermarkt}. Die Anbieter bleiben auf Ware
    sitzen und senken die Preise, um abzusetzen. Mit sinkendem Preis steigt die
    nachgefragte Menge, die angebotene Menge sinkt --- der Preis fällt in Richtung 15~€.
\end{itemize}
Ergänzend: Abgesetzt wird jeweils nur die Menge der „kürzeren`` Marktseite (bei 12~€:
420~Säcke, bei 20~€: 400~Säcke). Nur im Gleichgewicht wird mit 600~Säcken die größtmögliche
Menge umgesetzt.

\textbf{c) Preis-Mengen-Diagramm} \hfill (4~P.: Achsen 1~P., beide Kurven 1~P.,
Gleichgewicht 1~P., Überhang 1~P.)

Hilfspunkte: Nachfrage $p = 30 - x/40$ (Prohibitivpreis 30~€ bei $x = 0$; Sättigungsmenge
1.200~Säcke bei $p = 0$). Angebot $p = 5 + x/60$ (Mindestpreis 5~€ bei $x = 0$; bei
$x = 1.200$ ist $p = 25$~€).

\begin{center}
\begin{tikzpicture}[x=1cm, y=1cm, >=Stealth, font=\small]
  % Renten (für Teilaufgabe d)
  \fill[bfwAccent!25] (0,6) -- (0,3) -- (5,3) -- cycle;
  \fill[bfwPrimary!20] (0,1) -- (0,3) -- (5,3) -- cycle;
  % Achsen
  \draw[->, thick, bfwInk] (0,0) -- (10.8,0) node[right] {$x$};
  \draw[->, thick, bfwInk] (0,0) -- (0,6.8) node[above] {Preis $p$ (€/Sack)};
  \node[below, font=\footnotesize] at (5,-0.5) {Menge $x$ (Säcke pro Tag)};
  \foreach \x in {200,400,600,800,1000,1200}
    \draw (\x/120,0.08) -- (\x/120,-0.08) node[below] {\footnotesize\x};
  \foreach \p in {5,10,15,20,25,30}
    \draw (0.08,\p/5) -- (-0.08,\p/5) node[left] {\footnotesize\p};
  % Kurven
  \draw[very thick, bfwPrimaryDark] (0,6) -- (10,0) node[pos=0.97, above right] {$N$};
  \draw[very thick, bfwWarn] (0,1) -- (10,5) node[right] {$A$};
  % Gleichgewicht
  \draw[dashed, bfwInkSoft] (0,3) -- (5,3) -- (5,0);
  \fill[bfwInk] (5,3) circle (2.5pt);
  \node[below right, bfwInk] at (5.1,2.95) {GG ($x^*=600$; $p^*=15$)};
  % Überhang bei 20 Euro
  \draw[dashed, bfwInkSoft] (0,4) -- (7.5,4);
  \fill[bfwPrimaryDark] (10/3,4) circle (2pt);
  \fill[bfwWarn] (7.5,4) circle (2pt);
  \draw[<->, thick, bfwAccent!85!black] (10/3,4.35) -- (7.5,4.35)
    node[midway, above] {Angebotsüberhang 500};
  % Beschriftung Renten
  \node[font=\footnotesize] at (1.4,3.7) {KR};
  \node[font=\footnotesize] at (1.4,2.35) {PR};
\end{tikzpicture}
\end{center}

\textbf{d) Konsumentenrente} \hfill (3~P.: Rechnung 2~P., Erläuterung 1~P.)

Die Konsumentenrente ist das Dreieck zwischen Nachfragekurve und Gleichgewichtspreis
(Fläche KR im Diagramm):
\[
KR = \tfrac{1}{2} \cdot x^* \cdot (\text{Prohibitivpreis} - p^*)
   = \tfrac{1}{2} \cdot 600~\text{Säcke} \cdot (30~\text{€} - 15~\text{€})
   = \mathbf{4.500~\text{€ pro Tag}}
\]
\emph{Erläuterung:} Viele Käufer wären bereit gewesen, mehr als 15~€ je Sack zu zahlen.
Da sie nur den einheitlichen Marktpreis von 15~€ zahlen, entsteht ihnen ein
Nutzengewinn (eine „Ersparnis``) von insgesamt 4.500~€ pro Tag.
(Zum Vergleich: Produzentenrente $= \tfrac{1}{2} \cdot 600 \cdot (15 - 5) = 3.000$~€ pro Tag.)
\end{loesung}

% ---------------------------------------------------------------------------
\begin{loesung}[title={Lösung Aufgabe~1.5 --- Kurvenverschiebungen und Kreuzpreiselastizität}]
\textbf{a) Bewegung oder Verschiebung} \hfill (8~P.: je Entwicklung 2~P.)
\begin{enumerate}[label=(\arabic*)]
  \item \emph{Hitzesommer:} Die Bedürfnisstruktur bzw. Zukunftserwartung ändert sich,
    nicht der Preis $\Rightarrow$ \emph{Verschiebung der Nachfragekurve nach rechts}
    (Nachfrageerhöhung): Bei jedem Preis wird mehr Limonade nachgefragt.
  \item \emph{Zuckerpreis $+40\,\%$:} Die Kosten eines Produktionsfaktors (Rohstoff)
    steigen $\Rightarrow$ \emph{Verschiebung der Angebotskurve nach links}
    (Angebotssenkung): Bei jedem Preis wird weniger angeboten.
  \item \emph{Preis der Limonade sinkt:} Es ändert sich der Preis des Gutes selbst
    $\Rightarrow$ \emph{keine Verschiebung}, sondern \emph{Bewegung auf den Kurven}:
    auf der Nachfragekurve nach rechts unten (nachgefragte Menge steigt), auf der
    Angebotskurve nach links unten (angebotene Menge sinkt).
  \item \emph{Neue Abfülltechnik:} Technischer Fortschritt erhöht die Produktivität und
    senkt die Stückkosten $\Rightarrow$ \emph{Verschiebung der Angebotskurve nach rechts}
    (Angebotserhöhung).
\end{enumerate}

\textbf{b) Gleichzeitiger Eintritt von (1) und (2)} \hfill (4~P.)

Die Nachfragekurve verschiebt sich nach rechts, die Angebotskurve nach links. Beide
Verschiebungen wirken \emph{preissteigernd} $\Rightarrow$ Der Gleichgewichtspreis
\emph{steigt eindeutig}. Für die Gleichgewichtsmenge wirken die Effekte
\emph{gegenläufig}: Die Nachfrageerhöhung erhöht, die Angebotssenkung verringert die
Menge. Die Wirkung auf die Gleichgewichtsmenge ist daher \emph{unbestimmt}; sie hängt
davon ab, welche Verschiebung stärker ausfällt (überwiegt die Nachfrageerhöhung, steigt
die Menge; überwiegt die Angebotssenkung, sinkt sie).

\textbf{c) Kreuzpreiselastizität} \hfill (4~P.: Mengenänderung 1~P., Elastizität 2~P.,
Güterbeziehung 1~P.)

Mengenänderung Limonade: $\dfrac{11.750 - 12.500}{12.500} = \dfrac{-750~\text{Kisten}}{12.500~\text{Kisten}} = -6\,\%$
\[
\varepsilon_K = \frac{\text{Mengenänderung Limonade in \%}}{\text{Preisänderung Eistee in \%}}
= \frac{-6\,\%}{-20\,\%} = \mathbf{+0{,}3}
\]
Das \emph{positive} Vorzeichen zeigt: Preis des Eistees und Menge der Limonade bewegen
sich gleichgerichtet (Eistee billiger $\rightarrow$ weniger Limonade). Eistee und Limonade
sind \emph{Substitutionsgüter}. Da $|\varepsilon_K| = 0{,}3$ relativ klein ist, ist die
Substitutionsbeziehung nur schwach ausgeprägt.
\end{loesung}

% ---------------------------------------------------------------------------
\begin{loesung}[title={Lösung Aufgabe~1.6 --- Funktionen des Marktpreises}]
\textbf{a) Funktionen des Marktpreises} \hfill (6~P.: je Funktion 1~P. Erläuterung +
1~P. Situationsbezug; drei der folgenden)
\begin{itemize}
  \item \emph{Signalfunktion (Informationsfunktion):} Der Preis zeigt die Knappheit eines
    Gutes an. Der Anstieg von 420~€ auf 680~€ je m³ signalisiert allen Marktteilnehmern,
    dass Konstruktionsvollholz knapp geworden ist.
  \item \emph{Allokationsfunktion (Lenkungsfunktion):} Hohe Preise lenken
    Produktionsfaktoren in die Herstellung des knappen Gutes. Die Sägewerke setzen mehr
    Arbeit ein (Sonderschichten), ein Investor lenkt Kapital in ein neues Sägewerk.
  \item \emph{Selektionsfunktion (Sanktionsfunktion):} Marktteilnehmer, die den Preis
    nicht zahlen bzw. nicht weitergeben können, scheiden aus. Kleine Zimmereien, die die
    höheren Einkaufspreise nicht an ihre Kunden weitergeben können, geben auf.
  \item \emph{Ausgleichsfunktion:} Der höhere Preis bringt Angebots- und Nachfragemenge
    wieder zum Ausgleich: Die Nachfrager schränken ihre Käufe ein (z.\,B. Ausweichen auf
    andere Baustoffe), die Anbieter weiten ihre Menge aus, bis der Markt geräumt ist.
\end{itemize}

\textbf{b) Marktsituation und Preismechanismus} \hfill (4~P.)

Durch den Nachfrageanstieg hat sich die Nachfragekurve nach rechts verschoben. Beim alten
Preis von 420~€ je m³ übersteigt die nachgefragte Menge die angebotene Menge: Es besteht
ein \emph{Nachfrageüberhang} (Angebotsdefizit), also ein \emph{Verkäufermarkt} (1~P.).
Nicht alle Kaufwünsche können erfüllt werden; die Nachfrager überbieten sich, die
Anbieter setzen höhere Preise durch (1~P.). Mit steigendem Preis weiten die Sägewerke
ihre angebotene Menge aus (Bewegung auf der Angebotskurve) und ein Teil der Nachfrager
verzichtet oder weicht aus (Bewegung auf der neuen Nachfragekurve) (1~P.). Der Prozess
endet im neuen Gleichgewicht mit \emph{höherem} Preis und \emph{höherer} Menge als vor dem
Nachfrageanstieg (1~P.).
\end{loesung}

% ---------------------------------------------------------------------------
\begin{loesung}[title={Lösung Aufgabe~1.7 --- Monopolistischer Absatzbereich im unvollkommenen Polypol}]
\textbf{a) Marktform und Preisspielraum} \hfill (4~P.)

Es liegt ein \emph{unvollkommenes Polypol} vor: viele Anbieter, viele Nachfrager
(Polypol, 1~P.), aber ein unvollkommener Markt, weil die Güter \emph{heterogen} sind
(unterschiedliche Marken, Rezepturen, Verpackungen) und die Kunden \emph{Präferenzen}
haben (1~P.). Durch \emph{Produktdifferenzierung} --- hier das Steingut-Glas, das
traditionelle Etikett, die Marke „Altmühltaler`` und die besondere Rezeptur --- bindet
die Brinkhaus GmbH Stammkunden an sich (1~P.). Daraus entsteht ein
\emph{monopolistischer Absatzbereich} (zwischen 2,20~€ und 3,10~€): Innerhalb dieses
Bereichs kann das Unternehmen seinen Preis wie ein Monopolist weitgehend frei festlegen,
ohne größere Nachfrageeinbußen befürchten zu müssen (1~P.).

\textbf{b) Beurteilung der Preiserhöhungen} \hfill (4~P.: je Fall 2~P.)
\begin{enumerate}[label=(\arabic*)]
  \item \emph{2,95~€:} Der Preis liegt noch \emph{innerhalb} des monopolistischen
    Absatzbereichs (unter 3,10~€). Die Stammkunden bleiben weitgehend treu; der Absatz
    geht nur geringfügig zurück. Die Nachfrage reagiert hier unelastisch, der Umsatz
    dürfte steigen.
  \item \emph{3,40~€:} Der Preis liegt \emph{oberhalb} der Obergrenze von 3,10~€. Die
    Präferenzen reichen nicht mehr aus; viele Kunden wandern zu Konkurrenzmarken ab. Die
    Nachfrage reagiert hier stark elastisch, Absatz und Umsatz brechen deutlich ein.
    Die Erhöhung ist nicht zu empfehlen.
\end{enumerate}

\textbf{c) Maßnahmen zur Vergrößerung des Spielraums} \hfill (2~P.: je Maßnahme 1~P.)

Zum Beispiel: Ausbau der Markenbekanntheit durch Werbung; Qualitätsverbesserung bzw.
Betonung der Herkunft (Regionalität, Bio-Siegel); weitere Differenzierung durch
Sondereditionen oder Geschenkverpackungen; Kundenbindungsmaßnahmen und Service;
Vertrieb über exklusive Kanäle (Feinkosthandel).
\end{loesung}

% ---------------------------------------------------------------------------
\begin{loesung}[title={Lösung Aufgabe~1.8 --- Preisbildung im Oligopol und im Monopol}]
\textbf{a) Marktform} \hfill (2~P.)

\emph{Angebotsoligopol} (1~P.): Wenige Anbieter (drei Zementwerke) mit großen
Marktanteilen stehen vielen Nachfragern (mehrere Hundert Abnehmer) gegenüber (1~P.).

\textbf{b) Typische Verhaltensweisen} \hfill (4~P.: je Verhaltensweise 2~P.; zwei der
folgenden)
\begin{itemize}
  \item \emph{Parallelverhalten / Preisruhe:} Die Anbieter wissen, dass jede eigene
    Preisänderung die Konkurrenten unmittelbar trifft und eine Reaktion auslöst. Sie
    verhalten sich deshalb ohne ausdrückliche Absprache gleichgerichtet und lassen die
    Preise über längere Zeit unverändert (Preisstarre); Änderungen erfolgen im
    Gleichschritt.
  \item \emph{Preisführerschaft:} Ein (meist großer) Anbieter erhöht als Erster die
    Preise, z.\,B. wegen gestiegener Energiekosten; die übrigen Werke ziehen nach.
  \item \emph{Verlagerung des Wettbewerbs:} Statt über den Preis konkurrieren die Werke
    über Qualität, Lieferzuverlässigkeit, Service und Konditionen.
  \item \emph{Marktverdrängungspolitik:} Ein Anbieter versucht, mit aggressiven
    Preissenkungen Marktanteile zu gewinnen bzw. Konkurrenten zu verdrängen (selten).
\end{itemize}

\textbf{c) Beurteilung der Preissenkung um 10\,\%} \hfill (3~P.)

Die Erfolgsaussichten sind gering. Im Oligopol besteht eine enge
\emph{Reaktionsverbundenheit}: Eine Preissenkung von 10~\% zieht spürbar Kunden der
beiden Konkurrenten ab, die deshalb mit hoher Wahrscheinlichkeit sofort nachziehen
(1~P.). Die Marktanteile bleiben dann annähernd gleich, alle drei Werke erzielen jedoch
niedrigere Erlöse (1~P.). Im schlimmsten Fall entsteht ein \emph{ruinöser Preiskampf},
bei dem alle Anbieter Verluste erleiden. Sinnvoller ist ein Wettbewerb über Qualität und
Service (1~P.).

\textbf{d) Übernahme --- Angebotsmonopol} \hfill (3~P.: Erläuterung 1~P., je Grenze 1~P.)

Nach der Übernahme ist das Werk alleiniger Anbieter in der Region
(\emph{Angebotsmonopol}). Es muss keine Reaktionen von Konkurrenten mehr berücksichtigen
und kann den Preis als \emph{Preisfixierer} weitgehend selbst festsetzen. Grenzen der
Preissetzungsmacht sind z.\,B.:
\begin{itemize}
  \item die \emph{Nachfrage}: Mit steigendem Preis sinkt die Absatzmenge; Abnehmer
    verzichten auf Bauvorhaben oder weichen auf \emph{Substitute} aus (z.\,B. Holz-,
    Stahl- oder Ziegelbau);
  \item der Bezug aus \emph{entfernteren Werken bzw. Importe}, der sich bei stark
    steigenden Preisen trotz Transportkosten wieder lohnt;
  \item der drohende \emph{Markteintritt} neuer Anbieter, die von den hohen Preisen
    angelockt werden;
  \item die staatliche \emph{Wettbewerbsaufsicht} (Fusionskontrolle,
    Missbrauchsaufsicht durch das Bundeskartellamt --- vgl. Lektion~2).
\end{itemize}
\end{loesung}

\clearpage
\section*{Lösungen zu Lektion 2: Wettbewerbspolitik, Staatseingriffe \& VGR}
\addcontentsline{toc}{section}{Lösungen zu Lektion 2: Wettbewerbspolitik, Staatseingriffe \& VGR}
% ===========================================================================
% Gesamtskript VWL --- Musterlösungen Lektion 2
% Wettbewerb, Staatseingriffe & VGR (Rahmenplan DIHK 1.1)
% ===========================================================================

\begin{loesung}[title={Lösung Aufgabe~2.1 --- UWG oder GWB?}]
\textbf{a) Schutzzweck (4~P.: je Gesetz 2~P.)}\\
\emph{UWG:} schützt die \emph{Lauterkeit} des Wettbewerbs, also die \emph{Art und Weise},
wie Wettbewerb geführt wird. Geschützt werden Mitbewerber, Verbraucher und sonstige
Marktteilnehmer vor unlauteren geschäftlichen Handlungen (z.\,B. irreführende oder
aggressive Werbung, Herabsetzung von Mitbewerbern, unzumutbare Belästigung) sowie das
Interesse der Allgemeinheit an einem unverfälschten Wettbewerb (\S\,1 UWG).\\
\emph{GWB („Kartellgesetz``):} schützt den \emph{Bestand} bzw. die \emph{Freiheit} des
Wettbewerbs als Institution vor Beschränkungen --- durch Kartellverbot,
Missbrauchsaufsicht über marktbeherrschende Unternehmen und Fusionskontrolle. Über die
Einhaltung wacht das Bundeskartellamt.

\textbf{b) Zuordnung (6~P.: je Vorgang 2~P.)}
\begin{itemize}
  \item \textbf{(1) UWG} --- \emph{irreführende geschäftliche Handlung} (\S\,5 UWG): Die
    Angabe „Testsieger`` täuscht über die Qualität bzw. die Eigenschaften der Ware; die
    Werbung beeinflusst die Kaufentscheidung mit unwahren Angaben. Es geht um die
    \emph{Art} des Wettbewerbs, nicht um seine Beschränkung.
  \item \textbf{(2) UWG} --- \emph{unzumutbare Belästigung} (\S\,7 UWG): Telefonwerbung
    gegenüber Verbrauchern ist ohne deren vorherige ausdrückliche Einwilligung stets
    unzulässig.
  \item \textbf{(3) GWB} --- \emph{Preiskartell}, verboten nach \S\,1 GWB (europarechtlich
    Art.~101 AEUV): Die Absprache einer gemeinsamen Preiserhöhung schaltet den
    Preiswettbewerb zwischen den Anbietern aus. Dass die Vereinbarung nur mündlich auf
    einem Branchentreffen zustande käme, ist unerheblich --- auch aufeinander abgestimmte
    Verhaltensweisen sind verboten. Rheinquell sollte sich daran \emph{nicht} beteiligen
    und sich von dem Vorschlag erkennbar distanzieren.
\end{itemize}

\textbf{c) Rechtsfolgen (2~P.: je Gesetz 1~P., je zwei Nennungen)}\\
\emph{UWG:} Ansprüche auf Beseitigung und Unterlassung (\S\,8), Schadensersatz (\S\,9),
Gewinnabschöpfung (\S\,10); bei unerlaubter Telefonwerbung zusätzlich Bußgeld bis
300.000~€ (\S\,20). Klagebefugt sind u.\,a. Mitbewerber, Wirtschaftsverbände, IHK und
Verbraucherverbände.\\
\emph{GWB:} Bußgeld durch das Bundeskartellamt (bis zu 10~\% des weltweiten
Gesamtumsatzes des Vorjahres, \S\,81c GWB), Schadensersatzansprüche der Geschädigten
(\S\,33a GWB), Vorteilsabschöpfung, Unwirksamkeit der Absprache.
\end{loesung}

\begin{loesung}[title={Lösung Aufgabe~2.2 --- Kartellverbot und Ausnahmen}]
\textbf{a) Abgrenzung (2~P.)}\\
\emph{Kartell:} Die Unternehmen bleiben \emph{rechtlich selbstständig} und geben durch
die (auch mündliche) Vereinbarung nur einen \emph{Teil ihrer wirtschaftlichen}
Selbstständigkeit auf (z.\,B. bei der Preis- oder Gebietspolitik).
\emph{Fusion:} Die Unternehmen verschmelzen zu einer Einheit und verlieren ihre
\emph{rechtliche und wirtschaftliche} Selbstständigkeit vollständig. Für Fusionen gilt
nicht das Kartellverbot, sondern die Fusionskontrolle (\S\,35 GWB).

\textbf{b) Beurteilung der Vorhaben (8~P.: je Vorhaben 2~P. --- 1~P. Einordnung, 1~P. Begründung)}
\begin{itemize}
  \item \textbf{(1) Submissionskartell --- verboten.} Bei der Ausschreibung wird
    abgesprochen, welches Unternehmen zum Zug kommt; der öffentliche Auftraggeber zahlt
    überhöhte Preise. Verstoß gegen \S\,1 GWB; zusätzlich \emph{Straftat} nach \S\,298
    StGB (wettbewerbsbeschränkende Absprachen bei Ausschreibungen).
  \item \textbf{(2) Mittelstandskartell --- zulässig} (\S\,3 GWB). Kleine und mittlere
    Unternehmen rationalisieren durch zwischenbetriebliche Zusammenarbeit (gemeinsamer
    Einkauf), um Größennachteile gegenüber Großkonzernen auszugleichen. Voraussetzung:
    Der Wettbewerb wird nicht wesentlich beeinträchtigt und die Wettbewerbsfähigkeit der
    KMU verbessert --- hier erfüllt, da Preise und Angebote weiter selbstständig gestaltet
    werden.
  \item \textbf{(3) Forschungs- und Entwicklungskooperation --- freigestellt und damit
    zulässig.} Die Zusammenarbeit fördert den technischen Fortschritt, ohne den
    Wettbewerb zwischen den Beteiligten auszuschalten (getrennte Preissetzung, geringer
    gemeinsamer Marktanteil). Sie fällt unter die \emph{EU-Gruppenfreistellungsverordnung
    für F\&E-Vereinbarungen} bzw. die Legalausnahme nach \S\,2 GWB / Art.~101 Abs.~3 AEUV.
    Die Unternehmen müssen die Zulässigkeit selbst prüfen (Selbstprüfungssystem).
  \item \textbf{(4) Gebietskartell --- verboten.} Die Absatzmärkte werden regional
    aufgeteilt; in jedem Gebiet entfällt ein Wettbewerber, die Auftraggeber verlieren
    Auswahlmöglichkeiten. Verstoß gegen \S\,1 GWB.
\end{itemize}
\end{loesung}

\begin{loesung}[title={Lösung Aufgabe~2.3 --- Fusionskontrolle und Missbrauchsaufsicht}]
\textbf{a) Anmeldepflicht nach \S\,35 Abs.~1 GWB (5~P.)}\\
Anmeldepflichtig ist ein Zusammenschluss, wenn im letzten Geschäftsjahr \emph{beide}
Schwellen überschritten sind (1~P.):
\begin{enumerate}[label=(\arabic*)]
  \item weltweiter Umsatz aller beteiligten Unternehmen zusammen \emph{mehr als
    500~Mio.~€}: $430 + 85 = 515$~Mio.~€ $> 500$~Mio.~€ (erfüllt) (1~P.)
  \item im Inland mindestens ein Unternehmen \emph{mehr als 50~Mio.~€} (Nordlicht:
    390~Mio.~€, erfüllt) \emph{und} ein anderes \emph{mehr als 17,5~Mio.~€}
    (Hansemilch: 22~Mio.~€, erfüllt) (1~P.)
\end{enumerate}
$\Rightarrow$ Der Zusammenschluss ist \textbf{anmeldepflichtig} (1~P.).\\
\emph{Variante:} Bei nur 15~Mio.~€ Inlandsumsatz liegt Hansemilch unter 17,5~Mio.~€;
die zweite Inlandsschwelle ist nicht erfüllt $\Rightarrow$ \textbf{keine Anmeldepflicht}
nach den Umsatzschwellen (1~P.). (Ergänzend: Eine Anmeldepflicht könnte dann nur noch
über die Transaktionswertschwelle nach \S\,35 Abs.~1a GWB --- Gegenleistung über
400~Mio.~€ und erhebliche Inlandstätigkeit --- entstehen; das ist hier ausgeschlossen.)

\textbf{b) Verfahren und Untersagungsvoraussetzung (3~P.)}\\
Bis zur Freigabe gilt ein \emph{Vollzugsverbot} (\S\,41 GWB). Das Bundeskartellamt prüft
zunächst innerhalb eines Monats (Vorprüfverfahren); bei Bedenken leitet es das
Hauptprüfverfahren ein (Entscheidung grundsätzlich innerhalb von fünf Monaten ab
Anmeldung) (1~P.). Untersagt wird ein Zusammenschluss, wenn er \emph{wirksamen Wettbewerb
erheblich behindern} würde, insbesondere weil eine \emph{marktbeherrschende Stellung}
entsteht oder verstärkt wird (\S\,36 Abs.~1 GWB) (1~P.). Ein Unternehmen mit mindestens
40~\% Marktanteil gilt nach \S\,18 Abs.~4 GWB als marktbeherrschend (Vermutung) ---
mit 43~\% wäre diese Schwelle überschritten, eine Untersagung ist daher wahrscheinlich
(1~P.).

\textbf{c) Ministererlaubnis (2~P.)}\\
Nordlicht kann beim Bundeswirtschaftsminister eine \emph{Ministererlaubnis} (\S\,42 GWB)
beantragen. Er kann den untersagten Zusammenschluss erlauben, wenn die
gesamtwirtschaftlichen Vorteile die Wettbewerbsbeschränkung aufwiegen oder er durch ein
überragendes Interesse der Allgemeinheit gerechtfertigt ist und die marktwirtschaftliche
Ordnung nicht gefährdet wird. Die Erlaubnis ist eine seltene Ausnahme.

\textbf{d) Missbrauchsform (2~P.)}\\
\emph{Ausbeutungsmissbrauch} (\S\S\,19\,ff. GWB): Nordlicht nutzt ihre (Nachfrage-)Macht
gegenüber wirtschaftlich abhängigen Lieferanten aus, um unangemessen niedrige
Einkaufspreise durchzusetzen. Die Landwirte haben kaum Ausweichmöglichkeiten. Ein
\emph{Behinderungsmissbrauch} läge dagegen vor, wenn Konkurrenten gezielt verdrängt oder
vom Markt ferngehalten würden.
\end{loesung}

\begin{loesung}[title={Lösung Aufgabe~2.4 --- Höchstpreis auf dem Pelletmarkt}]
\textbf{a) Gleichgewicht (3~P.)}\\
Gleichsetzen: $900 - 1{,}5\,p = -300 + 2{,}5\,p \;\Rightarrow\; 1.200 = 4\,p
\;\Rightarrow\; p_0 = \textbf{300~€/t}$ (2~P.)\\
Einsetzen: $x_0 = 900 - 1{,}5 \cdot 300 = 900 - 450 = \textbf{450} \text{ (Tsd.~t)}$;
Probe: $x_A = -300 + 2{,}5 \cdot 300 = 450$ (stimmt) (1~P.)

\textbf{b) Höchstpreis 240~€/t (4~P.)}\\
$x_N = 900 - 1{,}5 \cdot 240 = 900 - 360 = \textbf{540 Tsd.~t}$ (1~P.)\\
$x_A = -300 + 2{,}5 \cdot 240 = -300 + 600 = \textbf{300 Tsd.~t}$ (1~P.)\\
\emph{Nachfrageüberhang} (Angebotslücke) von $540 - 300 = \textbf{240 Tsd.~t}$
($x_A < x_0 < x_N$) (1~P.).\\
Ein Höchstpreis von 320~€/t liegt \emph{über} dem Gleichgewichtspreis von 300~€/t. Der
Markt würde sich weiterhin beim zulässigen Gleichgewichtspreis von 300~€/t räumen --- der
Höchstpreis wäre \emph{wirkungslos} (1~P.).

\textbf{c) Schwarzmarktpreis und flankierende Maßnahmen (3~P.)}\\
Angeboten werden nur 300~Tsd.~t. Zahlungsbereitschaft für diese Menge aus der
Nachfragefunktion: $300 = 900 - 1{,}5\,p \;\Rightarrow\; 1{,}5\,p = 600
\;\Rightarrow\; p = \textbf{400~€/t}$. Auf dem Schwarzmarkt könnten also bis zu 400~€/t
erzielt werden --- mehr als der Gleichgewichtspreis (1~P.).\\
Flankierende Maßnahmen (2~P., je 1~P.):
\emph{Rationierung/Kontingentierung} der knappen Menge (z.\,B. Bezugsscheine je
Haushalt), damit die Verteilung nicht nach Zufall oder Beziehungen erfolgt;
\emph{Subventionen an die Anbieter}, damit trotz niedrigen Preises genügend produziert
wird; ggf. \emph{Kontrollen und Sanktionen} gegen den Schwarzhandel bzw. Förderung von
Importen.

\textbf{d) Umsatzveränderung der Anbieter (2~P.)}\\
Vorher: $300~\text{€/t} \cdot 450.000~\text{t} = 135~\text{Mio.~€}$;
nachher: $240~\text{€/t} \cdot 300.000~\text{t} = 72~\text{Mio.~€}$ (1~P.)\\
Veränderung: $72 - 135 = \textbf{$-$63~Mio.~€}$; prozentual
$-63 \div 135 \cdot 100 = \textbf{$-$46,67~\%}$ (1~P.)
\end{loesung}

\begin{loesung}[title={Lösung Aufgabe~2.5 --- Marktkonforme und marktkonträre Eingriffe}]
\textbf{a) Zuordnung (4~P.: je Vorschlag 1~P.)}
\begin{itemize}
  \item \textbf{(1) marktkonträr} --- staatlich festgesetzter \emph{Höchstpreis}; der
    Preis wird unmittelbar vorgegeben, der Preismechanismus außer Kraft gesetzt.
  \item \textbf{(2) marktkonform} --- Steuersenkung; sie verändert die Kosten der
    Anbieter, der Preis bildet sich weiterhin über Angebot und Nachfrage.
  \item \textbf{(3) marktkonform} --- \emph{Transferleistung} an Nachfrager; sie stärkt
    deren Kaufkraft, ohne den Preis festzulegen.
  \item \textbf{(4) marktkonform} --- \emph{Subvention} an Anbieter; sie beeinflusst das
    Angebot, der Preismechanismus bleibt erhalten.
\end{itemize}

\textbf{b) Wirkung der Stromsteuersenkung (3~P.)}\\
Die Stromsteuer ist ein Kostenbestandteil der Versorger. Sinkt sie, können die Anbieter
zu jedem Preis mehr bzw. jede Menge billiger anbieten: Die \emph{Angebotskurve verschiebt
sich nach rechts (unten)} (1~P.). Bei unveränderter Nachfrage \emph{sinkt der
Gleichgewichtspreis} (1~P.) und die \emph{Gleichgewichtsmenge steigt} (1~P.). Der Staat
nimmt weniger Stromsteuer ein; ein Teil der Entlastung kann bei den Anbietern verbleiben,
wenn sie die Senkung nicht voll weitergeben.

\textbf{c) Folgen der Preisobergrenze (3~P.: je Folge 1~P.)}
\begin{itemize}
  \item \emph{Nachfrageüberhang:} Zum niedrigen Preis wird mehr Strom nachgefragt, die
    Versorger bieten weniger an.
  \item \emph{Verluste und Rückzug von Anbietern:} Liegen die Beschaffungskosten der
    Stadtwerke über 30~Cent/kWh, machen sie mit jeder kWh Verlust; Anbieter nehmen keine
    Neukunden mehr auf oder ziehen sich aus dem Markt zurück, Investitionen (z.\,B. in
    Netze und erneuerbare Erzeugung) unterbleiben.
  \item \emph{Ausfall der Preisfunktionen:} Der Preis signalisiert die Knappheit nicht
    mehr (Signalfunktion), Sparanreize entfallen; der Staat müsste rationieren oder doch
    subventionieren (Folgekosten).
\end{itemize}

\textbf{d) Beurteilung (2~P.)}\\
Vorschlag (3) ist vorzuziehen: Der Zuschuss ist \emph{marktkonform} und
\emph{zielgenau} --- er entlastet nur bedürftige Haushalte, während der Strompreis
weiterhin die Knappheit anzeigt und Sparanreize sowie Investitionsanreize erhalten
bleiben. Der Höchstpreis (1) entlastet dagegen alle Haushalte unabhängig vom Einkommen,
verursacht Versorgungsengpässe und Folgeeingriffe. Nachteil von (3): Finanzierung durch
die Steuerzahler und Verwaltungsaufwand.
\end{loesung}

\begin{loesung}[title={Lösung Aufgabe~2.6 --- VGR: Vom BIP zum Volkseinkommen}]
\textbf{a) Außenbeitrag und BIP (4~P.)}\\
Außenbeitrag $=$ Exporte $-$ Importe $= 1.950 - 1.730 = \textbf{220~Mrd.~€}$ (1~P.)\\
BIP $=$ Private Konsumausgaben $+$ Konsumausgaben des Staates $+$ Bruttoinvestitionen
$+$ Außenbeitrag (1~P.)\\
BIP $= 2.180 + 860 + 740 + 220 = \textbf{4.000~Mrd.~€}$ (2~P.)

\textbf{b) BNE (3~P.)}\\
BNE $=$ BIP $-$ Primäreinkommen an die übrige Welt $+$ Primäreinkommen aus der übrigen Welt\\
BNE $= 4.000 - 145 + 210 = \textbf{4.065~Mrd.~€}$ (2~P.)\\
Das BNE ist um 65~Mrd.~€ größer, weil Inländer aus dem Ausland mehr Einkommen beziehen
(210~Mrd.~€) als Nichtinländer im Inland erzielen (145~Mrd.~€) --- der Saldo der
Primäreinkommen aus der übrigen Welt ist positiv (Inländerkonzept) (1~P.).

\textbf{c) Nettonationaleinkommen und Volkseinkommen (4~P.)}\\
Nettonationaleinkommen $=$ BNE $-$ Abschreibungen $= 4.065 - 690 =
\textbf{3.375~Mrd.~€}$ (2~P.)\\
Volkseinkommen $=$ Nettonationaleinkommen $-$ (Produktions- und Importabgaben $-$
Subventionen)\\
Volkseinkommen $= 3.375 - (420 - 70) = 3.375 - 350 = \textbf{3.025~Mrd.~€}$ (2~P.)

\textbf{d) Unternehmens- und Vermögenseinkommen, Lohn- und Gewinnquote (3~P.)}\\
Unternehmens- und Vermögenseinkommen $=$ Volkseinkommen $-$ Arbeitnehmerentgelt
$= 3.025 - 1.936 = \textbf{1.089~Mrd.~€}$ (1~P.)\\
Lohnquote $= 1.936 \div 3.025 \cdot 100 = \textbf{64~\%}$ (1~P.)\\
Gewinnquote $= 1.089 \div 3.025 \cdot 100 = \textbf{36~\%}$
(Kontrolle: $64~\% + 36~\% = 100~\%$) (1~P.)

\textbf{e) Steigende Lohnquote bei konstanten Löhnen (2~P.)}\\
Die Lohnquote ist ein \emph{Anteil}: Steigt sie bei unveränderten Arbeitnehmerentgelten,
muss das Volkseinkommen gesunken sein, d.\,h. die \emph{Unternehmens- und
Vermögenseinkommen sind zurückgegangen} --- typisch für einen konjunkturellen Abschwung,
in dem Gewinne stärker schwanken als Löhne (1~P.). Rechnerisch: $1.936 \div 0{,}66
\approx 2.933$~Mrd.~€ Volkseinkommen, also rund 92~Mrd.~€ weniger Unternehmens- und
Vermögenseinkommen. Eine steigende Lohnquote bedeutet also nicht automatisch, dass es
den Arbeitnehmern besser geht (1~P.).
\end{loesung}

\begin{loesung}[title={Lösung Aufgabe~2.7 --- Nominales und reales Wachstum}]
\textbf{a) Reales BIP (3~P.)}\\
Formel: reales BIP $=$ nominales BIP $\div$ Deflator $\cdot$ 100 (1~P.)\\
Jahr~2: $4.202{,}4 \div 103 \cdot 100 = \textbf{4.080,0~Mrd.~€}$ (1~P.)\\
Jahr~3: $4.263{,}0 \div 105 \cdot 100 = \textbf{4.060,0~Mrd.~€}$ (1~P.)\\
(Jahr~1 als Basisjahr: real $=$ nominal $= 4.000{,}0$~Mrd.~€)

\textbf{b) Wachstumsraten (4~P.: je Wert 1~P.)}\\
Formel: Wachstumsrate $=$ (Wert Berichtsjahr $-$ Wert Vorjahr) $\div$ Wert Vorjahr $\cdot$ 100
\begin{itemize}
  \item Jahr~2 nominal: $(4.202{,}4 - 4.000{,}0) \div 4.000{,}0 \cdot 100 = \textbf{+5,06~\%}$
  \item Jahr~2 real: $(4.080{,}0 - 4.000{,}0) \div 4.000{,}0 \cdot 100 = \textbf{+2,00~\%}$
  \item Jahr~3 nominal: $(4.263{,}0 - 4.202{,}4) \div 4.202{,}4 \cdot 100 = \textbf{+1,44~\%}$
  \item Jahr~3 real: $(4.060{,}0 - 4.080{,}0) \div 4.080{,}0 \cdot 100 = \textbf{$-$0,49~\%}$
\end{itemize}

\textbf{c) Preissteigerung und Erläuterung (3~P.)}\\
Preissteigerung Jahr~2 $\rightarrow$ Jahr~3: $(105 - 103) \div 103 \cdot 100 =
\textbf{1,94~\%}$ (1~P.). (Kontrolle: $1{,}0144 \approx 0{,}9951 \cdot 1{,}0194$.)\\
Erläuterung (2~P.): Das \emph{nominale} BIP steigt in Jahr~3 nur deshalb um 1,44~\%, weil
die Preise um 1,94~\% gestiegen sind. Die produzierte Gütermenge --- gemessen am
\emph{realen} BIP --- ist dagegen um 0,49~\% \emph{gesunken}. Wirtschaftswachstum wird
stets am realen BIP gemessen; die Medien sprechen daher zu Recht von einer schrumpfenden
Wirtschaft. Das Wachstumsziel des magischen Vierecks (rund 2~\% reales Wachstum) ist in
Jahr~3 verfehlt, in Jahr~2 dagegen erreicht.

\textbf{d) BIP-Deflator vs. VPI (2~P.)}\\
Der \emph{BIP-Deflator} misst die Preisentwicklung \emph{aller im Inland produzierten}
Güter und Dienstleistungen (Konsum, Investitionen, Staatsleistungen, Exporte);
Importgüter sind nicht enthalten (1~P.). Der \emph{Verbraucherpreisindex} misst die
Preisentwicklung eines repräsentativen \emph{Warenkorbs der privaten Haushalte} ---
einschließlich importierter Konsumgüter (z.\,B. Energie), aber ohne Investitionsgüter.
Er dient zur Messung der Inflationsrate (Ziel Preisniveaustabilität) (1~P.).
\end{loesung}

\begin{loesung}[title={Lösung Aufgabe~2.8 --- Magisches Viereck: Lage beurteilen}]
\textbf{a) Zielerreichung (8~P.: je Ziel 2~P. --- Zielwert und Beurteilung)}
\begin{itemize}
  \item \emph{Preisniveaustabilität} --- Zielwert: Inflationsrate von 2~\% (HVPI,
    symmetrisches EZB-Ziel). Mit 4,1~\% deutlich überschritten $\Rightarrow$
    \textbf{nicht erreicht}; die Kaufkraft sinkt spürbar.
  \item \emph{Hoher Beschäftigungsstand} --- Zielwert: Arbeitslosenquote von ca.
    3--4~\%. Mit 3,4~\% \textbf{erreicht}; es besteht nur noch Sockelarbeitslosigkeit.
  \item \emph{Stetiges und angemessenes Wirtschaftswachstum} --- Zielwert: reales
    BIP-Wachstum von rund 2~\% p.\,a. Mit 0,2~\% nahezu Stagnation $\Rightarrow$
    \textbf{nicht erreicht}.
  \item \emph{Außenwirtschaftliches Gleichgewicht} --- Zielwert: positiver Außenbeitrag
    von 1--2~\% des nominalen BIP. Mit 5,5~\% deutlich darüber $\Rightarrow$
    \textbf{nicht erreicht}; auch ein zu hoher Überschuss ist ein Ungleichgewicht.
\end{itemize}
Fazit: Nur ein Ziel ist erreicht --- die Lage zeigt typische Züge einer
\emph{Stagflation} (hohe Inflation bei stagnierendem Wachstum).

\textbf{b) Zielkonflikt (3~P.)}\\
Kürzungen der Staatsausgaben senken die gesamtwirtschaftliche Nachfrage (1~P.). Dies
dämpft den Preisauftrieb und dient der \emph{Preisniveaustabilität} (1~P.). Zugleich
erhalten Unternehmen weniger Aufträge, investieren und produzieren weniger: Das ohnehin
schwache \emph{Wachstum} sinkt weiter, und der \emph{hohe Beschäftigungsstand} ist
gefährdet, weil Arbeitskräfte abgebaut werden. Es besteht also ein Zielkonflikt zwischen
Preisniveaustabilität einerseits und Wachstum/Beschäftigung andererseits (1~P.).

\textbf{c) Beurteilung der Aussage (3~P.)}\\
Teilweise richtig: Ein positiver Außenbeitrag zeigt eine hohe internationale
Wettbewerbsfähigkeit der Exportwirtschaft (1~P.). Dauerhaft hohe Überschüsse sind aber
ein \emph{außenwirtschaftliches Ungleichgewicht}: Das Land produziert mehr, als es
selbst verbraucht und investiert (schwache Binnennachfrage); es entsteht eine hohe
\emph{Abhängigkeit vom Ausland} (Konjunktur- und Handelskonflikte schlagen direkt durch)
(1~P.). Den Überschüssen stehen Defizite und Verschuldung der Handelspartner gegenüber;
das kann zu Gegenmaßnahmen wie Zöllen sowie zu Aufwertungsdruck führen. Die Aussage „je
höher, desto besser`` ist daher abzulehnen --- Ziel ist ein moderater Überschuss
(1~P.).
\end{loesung}

\clearpage
\section*{Lösungen zu Lektion 3: Konjunktur, Wirtschaftspolitik \& Außenwirtschaft}
\addcontentsline{toc}{section}{Lösungen zu Lektion 3: Konjunktur, Wirtschaftspolitik \& Außenwirtschaft}
% ============================================================
% Gesamtskript VWL --- Musterlösungen Lektion 3
% Konjunktur, Wirtschaftspolitik & Außenwirtschaft (Rahmenplan DIHK 1.1)
% Punktangaben in eckigen Klammern = Bewertungsvorschlag
% ============================================================

\begin{loesung}[title={Lösung Aufgabe~3.1 --- Konjunkturlage richtig deuten}]
\textbf{a) Zuordnung} \hfill [5~P.: je richtige Zuordnung 1~P.]

\begin{tabular}{p{7.6cm} p{5.6cm}}
\toprule
\textbf{Datum} & \textbf{Indikatorart}\\
\midrule
(1) ifo-Geschäftsklimaindex & Frühindikator\\
(2) Auftragseingänge & Frühindikator\\
(3) Kapazitätsauslastung & Präsenzindikator\\
(4) Arbeitslosenquote & Spätindikator\\
(5) Unternehmensinsolvenzen & Spätindikator\\
\bottomrule
\end{tabular}

\medskip
\textbf{b) Konjunkturphase} \hfill [3~P.]

Die Volkswirtschaft befindet sich am \emph{Ende der Tiefphase (Depression/Talsohle)} bzw.
am \emph{Übergang zum Aufschwung} (unterer Wendepunkt). [1~P.]

Begründung: Die \emph{Präsenz}- und \emph{Spätindikatoren} zeigen noch eine schwache Lage
(unterdurchschnittliche Auslastung, steigende Arbeitslosigkeit, viele Insolvenzen). Die
\emph{Frühindikatoren} (Geschäftsklima, Auftragseingänge) steigen jedoch bereits und
kündigen eine Belebung an. Da Spätindikatoren der Entwicklung hinterherlaufen, ist ihr
negativer Verlauf mit einer beginnenden Erholung vereinbar. [2~P.]

\medskip
\textbf{c) Schlussfolgerungen} \hfill [3~P.: zwei Schlussfolgerungen, schlüssig erläutert]

Erwartet werden z.\,B. (zwei genügen):
\begin{itemize}
  \item \emph{Personal:} Fachkräfte halten statt entlassen (ggf. Kurzarbeit nutzen und
    schrittweise zurückfahren), Einstellungen für steigende Auftragszahlen vorbereiten
    --- im Aufschwung wird der Arbeitsmarkt schnell enger.
  \item \emph{Produktion/Kapazität:} Kapazitäten so planen, dass steigende Aufträge
    bedient werden können (freie Kapazitäten zunächst auslasten, Erweiterungen vorbereiten).
  \item \emph{Investitionen/Finanzierung:} Investitionen jetzt planen, solange Zinsen und
    Preise für Investitionsgüter in der Tiefphase noch vergleichsweise niedrig sind.
  \item \emph{Einkauf/Lager:} Material frühzeitig disponieren, da im Aufschwung Preise
    und Lieferzeiten steigen.
\end{itemize}
\end{loesung}

\begin{loesung}[title={Lösung Aufgabe~3.2 --- Inflation, Realzins und Lohnabschluss}]
\textbf{a) Inflationsrate} \hfill [3~P.]
\[
\text{Inflationsrate} = \frac{127{,}1 - 124{,}0}{124{,}0} \times 100
= \frac{3{,}1}{124{,}0} \times 100 = 0{,}025 \times 100 = \textbf{2{,}5~\%}
\]
Probe: $124{,}0 \times 1{,}025 = 127{,}1$. \emph{Typischer Fehler:} Punktdifferenz
(3,1~Punkte) mit der Inflationsrate gleichsetzen --- es muss durch den alten Indexstand
geteilt werden.

\medskip
\textbf{b) Zinsertrag und Realzins} \hfill [4~P.]

Zinsertrag: $250.000~\text{€} \times 0{,}034 = \textbf{8.500~€}$ [1~P.]

Realzins (Näherung): Nominalzins $-$ Inflationsrate $= 3{,}4~\% - 2{,}5~\% =
\textbf{0{,}9~\%}$ [2~P.]\\
{\small (Exakt: $1{,}034 \div 1{,}025 - 1 \approx 0{,}0088 = 0{,}88~\%$ --- beide Wege
werden anerkannt.)}

Beurteilung: Der reale Vermögenszuwachs ist gering. Von den 8.500~€ Zinsen gleicht der
größte Teil nur den Kaufkraftverlust aus ($250.000 \times 0{,}025 = 6.250$~€); nur rund
0,9~\% sind echter Kaufkraftgewinn (vor Steuern). Nach Abzug von Steuern kann der
Realzins sogar negativ werden. [1~P.]

\medskip
\textbf{c) Nominal- und Reallohn} \hfill [3~P.]

Nominale Steigerung: $\dfrac{3.264 - 3.200}{3.200} \times 100 = \dfrac{64}{3.200} \times 100
= \textbf{2{,}0~\%}$ [1~P.]

Reallohnänderung (Näherung): $2{,}0~\% - 2{,}5~\% = \textbf{$-$0{,}5~\%}$ [1~P.]\\
{\small (Exakt: $1{,}020 \div 1{,}025 - 1 \approx -0{,}49~\%$.)}

Interpretation: Obwohl der Disponent nominal mehr verdient, kann er sich real
\emph{weniger} kaufen, weil die Preise stärker gestiegen sind als sein Gehalt
(Kaufkraftverlust). [1~P.]

\medskip
\textbf{d) Verteilungsspielraum und Beurteilung} \hfill [4~P.]

Verteilungsspielraum $=$ Produktivitätsfortschritt $+$ Zielinflation
$= 1{,}3~\% + 2{,}0~\% = \textbf{3{,}3~\%}$ [1~P.]\\
Die Forderung liegt $5{,}5 - 3{,}3 = \textbf{2{,}2~Prozentpunkte}$ über dem Spielraum.
[1~P.]

\emph{Sicht der Arbeitgeber} (produktivitätsorientierte Lohnpolitik): Ein Abschluss über
3,3~\% erhöht die Lohnstückkosten. Folgen: Preiserhöhungen (Lohn-Preis-Spirale,
Inflationsdruck), sinkende Wettbewerbsfähigkeit, Rationalisierung oder Verlagerung
--- die Forderung ist zu hoch. [1~P.]

\emph{Sicht der Gewerkschaft} (Kaufkrafttheorie): Die tatsächliche Inflation (2,5~\%)
liegt über dem Ziel; die Beschäftigten haben real verloren (vgl.~c). Höhere Löhne
sichern die Kaufkraft, stärken die Nachfrage und damit Produktion und Beschäftigung
--- die Forderung ist berechtigt. [1~P.]
\end{loesung}

\begin{loesung}[title={Lösung Aufgabe~3.3 --- Wechselkursänderung im Exportgeschäft}]
\emph{Rechenregel Mengennotierung (1~€ $=$ x~USD):} Europreis $\times$ Kurs $=$
Dollarpreis; Dollarpreis $\div$ Kurs $=$ Europreis.

\medskip
\textbf{a) Umrechnung zum Kurs 1,08} \hfill [2~P.: je 1~P.]

Messgerät: $45.000~\text{€} \times 1{,}08 = \textbf{48.600~USD}$\\
Sensoren: $32.400~\text{USD} \div 1{,}08 = \textbf{30.000~€}$

\medskip
\textbf{b) Kurs 1,20} \hfill [5~P.]

Für 1~€ erhält man nun mehr US-Dollar (1,20 statt 1,08) $\Rightarrow$ der Euro hat
\textbf{aufgewertet} (der Dollar abgewertet). [1~P.]

Neuer Dollarpreis: $45.000~\text{€} \times 1{,}20 = \textbf{54.000~USD}$ [1~P.]\\
Veränderung: $\dfrac{54.000 - 48.600}{48.600} \times 100 = \dfrac{5.400}{48.600} \times 100
\approx \textbf{+11,11~\%}$ [1~P.]

Sensoren: $32.400~\text{USD} \div 1{,}20 = \textbf{27.000~€}$, also
$3.000$~€ bzw. $\dfrac{3.000}{30.000} \times 100 = 10~\%$ günstiger. [2~P.]

Fazit: Die Aufwertung verteuert die Exporte für US-Kunden und verbilligt die Importe
für Nordlicht Optik.

\medskip
\textbf{c) Unveränderter Dollarpreis} \hfill [2~P.]

Euro-Erlös: $48.600~\text{USD} \div 1{,}20 = \textbf{40.500~€}$ [1~P.]\\
Erlösminderung: $45.000 - 40.500 = \textbf{4.500~€}$, das sind
$\dfrac{4.500}{45.000} \times 100 = \textbf{10~\%}$ je Gerät. [1~P.]

\medskip
\textbf{d) Absicherung} \hfill [2~P.: je Nennung 1~P.]

Z.\,B.: Devisentermingeschäft (Kurs wird heute für einen späteren Zahlungstermin
festgeschrieben); Devisenoption; Fakturierung in Euro (Risiko trägt der Kunde);
Kurssicherungsklauseln im Vertrag; „natürliche Absicherung“, indem Einkäufe (z.\,B.
Sensoren) und Verkäufe in derselben Währung abgewickelt werden.

\medskip
\textbf{e) Lieferungen nach Österreich} \hfill [3~P.]

Österreich gehört zur Eurozone; Lieferungen werden in derselben Währung (Euro) abgerechnet
--- es gibt keinen Wechselkurs und damit weder Kursrisiko noch Umtauschkosten. [1~P.]

Weitere Vorteile (zwei genügen, je 1~P.): keine Zölle und keine mengenmäßigen
Beschränkungen im Binnenmarkt (freier Warenverkehr); harmonisierte technische Normen
(z.\,B. CE-Kennzeichnung), dadurch keine zusätzlichen Zulassungen; Preistransparenz und
einfache Kalkulation; schneller, günstiger Zahlungsverkehr (SEPA); Monteure können
ohne Grenzkontrollen und Arbeitserlaubnis eingesetzt werden (freier Personen- und
Dienstleistungsverkehr).
\end{loesung}

\begin{loesung}[title={Lösung Aufgabe~3.4 --- Regionaler Arbeitsmarkt}]
\textbf{a) Arbeitslosenquoten} \hfill [4~P.]

Formel: Arbeitslosenquote $= \dfrac{\text{registrierte Arbeitslose}}{\text{zivile
Erwerbspersonen}} \times 100$ mit Erwerbspersonen $=$ Erwerbstätige $+$ Arbeitslose.

Vorjahr: Erwerbspersonen $= 236.800 + 13.200 = 250.000$;\quad
$\dfrac{13.200}{250.000} \times 100 = \textbf{5,28~\%}$ [1,5~P.]

Berichtsjahr: Erwerbspersonen $= 234.000 + 16.000 = 250.000$;\quad
$\dfrac{16.000}{250.000} \times 100 = \textbf{6,40~\%}$ [1,5~P.]

Veränderung: $6{,}40 - 5{,}28 = \textbf{+1,12~Prozentpunkte}$ [1~P.]\\
{\small \emph{Achtung:} Prozentpunkte, nicht Prozent! (Relativ wäre der Anstieg der
Quote $1{,}12 \div 5{,}28 \approx 21{,}2~\%$.) Häufiger Fehler: nur durch die
Erwerbstätigen statt durch die Erwerbspersonen teilen.}

\medskip
\textbf{b) Formen der Arbeitslosigkeit} \hfill [4~P.: je Fall 1~P.]

\begin{tabular}{p{1.1cm} p{2.6cm} p{9.4cm}}
\toprule
\textbf{Fall} & \textbf{Form} & \textbf{Begründung}\\
\midrule
(1) & saisonal & jahreszeitlich bedingt (Witterung im Baugewerbe), kehrt regelmäßig wieder\\
(2) & konjunkturell & Ursache ist ein allgemeiner Nachfragerückgang (Rezession)\\
(3) & strukturell & dauerhafter Wandel (technischer Fortschritt); Qualifikation passt nicht mehr zur Nachfrage\\
(4) & friktionell & kurze Such- bzw. Übergangsphase beim Arbeitsplatzwechsel\\
\bottomrule
\end{tabular}

\medskip
\textbf{c) Anteil struktureller Arbeitslosigkeit} \hfill [2~P.]

$\dfrac{7.400}{16.000} \times 100 = \textbf{46,25~\%}$ [1~P.]

Fast die Hälfte der Arbeitslosigkeit ist strukturell bedingt. Sie verschwindet auch im
Aufschwung nicht von selbst, weil Qualifikation, Branche oder Region nicht zu den offenen
Stellen passen (Sockelarbeitslosigkeit); sie ist nur langfristig durch Qualifizierung und
Strukturwandel abzubauen. Es droht Langzeitarbeitslosigkeit mit Kaufkraft- und
Steuerausfällen für die Region. [1~P.]

\medskip
\textbf{d) Maßnahmen} \hfill [3~P.]

\emph{Strukturell:} Umschulung bzw. Weiterbildung (z.\,B. zum Mediengestalter),
Mobilitätshilfen, Ansiedlung neuer Branchen durch Strukturpolitik $\Rightarrow$
\textbf{aktive} Arbeitsmarktpolitik. [1,5~P.]

\emph{Konjunkturell:} Kurzarbeitergeld (Arbeitsplätze werden über die Krise gehalten)
$\Rightarrow$ \textbf{passive} Arbeitsmarktpolitik; ergänzend wirken expansive Geld- und
Fiskalpolitik. [1,5~P.]\\
{\small (Andere sachgerechte Maßnahmen mit korrekter Zuordnung werden anerkannt.)}
\end{loesung}

\begin{loesung}[title={Lösung Aufgabe~3.5 --- Leitzinssenkung der EZB}]
\textbf{a) Begründung der Zinssenkung} \hfill [3~P.]

Vorrangiges Ziel der EZB ist Preisniveaustabilität, definiert als Anstieg des HVPI von
\textbf{2~\% mittelfristig}. Das Ziel ist \textbf{symmetrisch}: Ein dauerhaftes
\emph{Unterschreiten} ist ebenso unerwünscht wie ein Überschreiten. [1~P.]

Bei 1,3~\% und Prognosen deutlich unter 2~\% verfehlt die EZB ihr Ziel nach unten. Es
droht ein Abgleiten in sehr niedrige Inflation oder Deflation: Verbraucher und
Unternehmen schieben Käufe und Investitionen auf, Nachfrage und Beschäftigung sinken
(Abwärtsspirale). [1~P.]

Zudem stagniert die Wirtschaft --- eine expansive Geldpolitik soll die Nachfrage beleben
und die Inflation mittelfristig zum Zielwert zurückführen. [1~P.]

\medskip
\textbf{b) Wirkungskette (expansive Geldpolitik)} \hfill [5~P.: je Stufe 1~P.]

Leitzinssenkung $\Rightarrow$
(1)~Refinanzierung der Geschäftsbanken bei der EZB wird billiger $\Rightarrow$
(2)~Banken geben niedrigere Zinsen weiter, Bauzinsen/Kreditzinsen sinken (Sparen wird
unattraktiver) $\Rightarrow$
(3)~Baufinanzierungen werden günstiger, monatliche Kreditraten sinken; mehr Haushalte
können sich Wohneigentum leisten $\Rightarrow$
(4)~Kreditnachfrage und Nachfrage nach Eigentumswohnungen steigen (auch Kapitalanleger
investieren eher in Immobilien) $\Rightarrow$
(5)~Auftragslage und Umsatz der Wohnbau Lindner GmbH verbessern sich; zugleich sinken
ihre eigenen Finanzierungskosten für Bauprojekte $\Rightarrow$ gesamtwirtschaftlich
steigen Produktion und Beschäftigung, die Inflationsrate nähert sich wieder 2~\%.

\medskip
\textbf{c) Stellungnahme} \hfill [2~P.]

Die Aussage ist falsch. Die Leitzinsen legt der \textbf{EZB-Rat} fest; die EZB ist
\textbf{unabhängig} und an Weisungen der Regierungen nicht gebunden (Geldpolitik). Die
Bundesregierung betreibt Fiskalpolitik (Steuern, Staatsausgaben). [1~P.]\\
Ziel der EZB ist auch nicht die Rettung einzelner Branchen, sondern die
Preisniveaustabilität im gesamten Euroraum; die Bauwirtschaft profitiert nur als
Nebenwirkung. [1~P.]

\medskip
\textbf{d) Grenzen der Geldpolitik} \hfill [2~P.: je Grund 1~P.]

Z.\,B.: Banken geben die Zinssenkung nicht oder nur verzögert weiter (z.\,B. wegen
höherer Risikoaufschläge); pessimistische Erwartungen (Angst vor Arbeitsplatzverlust,
unsichere Konjunktur) --- selbst günstige Kredite werden nicht nachgefragt; hohe Bau- und
Grundstückspreise belasten die Erschwinglichkeit stärker, als die kleine Zinssenkung
entlastet; Wirkungsverzögerungen (Timelags) der Geldpolitik.
\end{loesung}

\begin{loesung}[title={Lösung Aufgabe~3.6 --- Konjunkturprogramm und Schuldenbremse}]
\textbf{a) Zuordnung} \hfill [4~P.: je Maßnahme 0,5~P. Zuordnung, 2~P. Gesamturteil]

\begin{tabular}{p{6.4cm} p{2.4cm} p{4.4cm}}
\toprule
\textbf{Maßnahme} & \textbf{Wirkung} & \textbf{in der Rezession geeignet?}\\
\midrule
(1) Investitionen vorziehen & expansiv & ja (Staatsnachfrage steigt)\\
(2) Umsatzsteuer erhöhen & restriktiv & nein (Kaufkraft sinkt)\\
(3) Sonderabschreibung & expansiv & ja (Steuerentlastung, Investitionsanreiz)\\
(4) Subventionen kürzen, Rücklage & restriktiv & nein (Boom-Maßnahme)\\
\bottomrule
\end{tabular}

\smallskip
Konjunkturgerecht im Sinne der antizyklischen Fiskalpolitik sind in der Rezession nur die
expansiven Maßnahmen \textbf{(1) und (3)}.

\medskip
\textbf{b) Schuldenbremse (Art.~109, 115~GG)} \hfill [4~P.]

Die Schuldenbremse begrenzt die \emph{strukturelle} (konjunkturunabhängige)
Neuverschuldung des Bundes auf 0,35~\% des BIP; die konjunkturbedingte Kreditaufnahme
ist zusätzlich zulässig.

Zulässige strukturelle Neuverschuldung: $4.400~\text{Mrd.~€} \times 0{,}0035 =
\textbf{15,4~Mrd.~€}$ [1~P.]\\
Geplante strukturelle Neuverschuldung: $41 - 22 = \textbf{19~Mrd.~€}$ [1~P.]\\
$\Rightarrow$ $19 - 15{,}4 = \textbf{3,6~Mrd.~€}$ über der Grenze (strukturell
$19 \div 4.400 \times 100 \approx 0{,}43~\%$ des BIP). [1~P.]

Ergebnis: Die Schuldenbremse wird \textbf{nicht eingehalten}. Zulässig wäre die Planung
nur bei einer vom Bundestag festgestellten außergewöhnlichen Notsituation (mit
Tilgungsplan); sonst muss das strukturelle Defizit um 3,6~Mrd.~€ gesenkt werden. [1~P.]

\medskip
\textbf{c) Prozyklische Wirkung und Crowding-out} \hfill [3~P.]

\emph{Prozyklische Wirkung (Timelags):} Zwischen Erkennen der Rezession, Beschluss,
Ausschreibung und Bauausführung vergeht viel Zeit (Erkennungs-, Entscheidungs- und
Wirkungsverzögerung). Setzt die zusätzliche Nachfrage erst ein, wenn die Wirtschaft
schon wieder im Aufschwung ist, verstärkt sie den Boom (Kapazitätsengpässe,
steigende Baupreise) --- die Maßnahme wirkt dann prozyklisch. [2~P.]

\emph{Crowding-out:} Die kreditfinanzierte Staatsnachfrage erhöht die Nachfrage am
Kapitalmarkt und damit die Zinsen; private Investitionen werden dadurch verdrängt, sodass
der gesamtwirtschaftliche Impuls ganz oder teilweise verpufft. [1~P.]
\end{loesung}

\begin{loesung}[title={Lösung Aufgabe~3.7 --- Zahlungsbilanz und Handelskonflikt}]
\textbf{a) Teilbilanzen (in Mrd.~€)} \hfill [4~P.]

Handelsbilanz: $1.480 - 1.290 = \textbf{+190}$ (Überschuss) [1~P.]\\
Dienstleistungsbilanz: $395 - 455 = \textbf{$-$60}$ (Defizit) [1~P.]\\
Leistungsbilanz: $190 + (-60) + 115 + (-72) = \textbf{+173}$ (Überschuss) [2~P.]

\medskip
\textbf{b) Ausgleich der Zahlungsbilanz} \hfill [2~P.]

Jede Transaktion wird doppelt verbucht (Leistung und Gegenleistung). Dem
Leistungsbilanzüberschuss von 173~Mrd.~€ steht daher spiegelbildlich ein
\textbf{Nettokapitalexport} in der \textbf{Kapitalbilanz} in gleicher Höhe gegenüber
(abgesehen von Vermögensänderungsbilanz und Restposten): Deutschland erwirbt Forderungen
bzw. Vermögen im Ausland. Ungleichgewichte zeigen sich deshalb nur in den
\emph{Teilbilanzen}, die Zahlungsbilanz insgesamt ist formal stets ausgeglichen.

\medskip
\textbf{c) Handelshemmnisse} \hfill [4~P.]

\emph{Tarifäre Handelshemmnisse} wirken über den \emph{Preis} der gehandelten Güter
(preispolitische Instrumente). \emph{Nichttarifäre Handelshemmnisse} wirken über
\emph{Mengen} oder \emph{administrative/technische Hürden}. [2~P.]

Zuordnung [je 0,5~P.]: (1)~Einfuhrzoll $\Rightarrow$ tarifär; (2)~Einfuhrkontingent
(Höchstmenge) $\Rightarrow$ nichttarifär; (3)~Sicherheitszertifizierung (technische Norm,
Verwaltungsverfahren) $\Rightarrow$ nichttarifär; (4)~Exportsubvention $\Rightarrow$
tarifär.

\medskip
\textbf{d) Beurteilung von Gegenzöllen} \hfill [4~P.: je Argument 1~P.]

\emph{Dafür} (z.\,B.):
\begin{itemize}
  \item Druckmittel, um Westland an den Verhandlungstisch zu bringen und die Zölle
    zurückzunehmen.
  \item Schutz europäischer Hersteller vor subventionierten Billigimporten aus
    Westland (Herstellung fairer Wettbewerbsbedingungen).
\end{itemize}
\emph{Dagegen} (z.\,B.):
\begin{itemize}
  \item Gefahr einer Eskalation (Handelskrieg) mit weiteren Zöllen auf EU-Waren ---
    gerade für ein exportorientiertes Unternehmen wie Weinmann ein großes Risiko.
  \item Importe aus Westland (Vorprodukte, Konsumgüter) verteuern sich; Kosten für
    Unternehmen und Preise für Verbraucher steigen, Wohlfahrtsgewinne der
    internationalen Arbeitsteilung gehen verloren.
\end{itemize}
Hinweis: Über Zölle entscheidet nicht Deutschland, sondern die \textbf{EU} (gemeinsame
Handelspolitik, gemeinsamer Außenzoll); Maßnahmen müssen mit den WTO-Regeln vereinbar
sein. Ein begründetes eigenes Fazit wird erwartet.
\end{loesung}

\begin{loesung}[title={Lösung Aufgabe~3.8 --- Emissionshandel im Unternehmen}]
\textbf{a) Funktionsweise und Prinzip} \hfill [3~P.]

Der Staat bzw. die EU legt eine \textbf{Obergrenze} für die Gesamtemissionen der
erfassten Anlagen fest (Cap), die schrittweise sinkt. In dieser Höhe werden
Emissionsberechtigungen (Zertifikate) ausgegeben bzw. versteigert. Ein Unternehmen darf nur
so viel CO\textsubscript{2} ausstoßen, wie es Berechtigungen abgeben kann. Fehlende
Berechtigungen müssen am Markt gekauft werden, nicht benötigte können verkauft werden
(Trade). So entsteht ein Marktpreis für CO\textsubscript{2}, der zur Vermeidung dort
anreizt, wo sie am günstigsten ist. [2~P.]

Prinzip: \textbf{Verursacherprinzip}; externe Kosten werden internalisiert. [1~P.]

\medskip
\textbf{b) Kostenvergleich} \hfill [5~P.]

\emph{Ohne Anlage:} Fehlmenge $46.000 - 40.000 = 6.000$ Berechtigungen;\\
Zukaufkosten $6.000 \times 80~\text{€} = \textbf{480.000~€}$ [1,5~P.]

\emph{Mit Anlage:} Emissionen $46.000 - 9.000 = 37.000$~t; Überschuss
$40.000 - 37.000 = 3.000$ Berechtigungen;\\
Verkaufserlös $3.000 \times 80~\text{€} = 240.000~\text{€}$;\\
Nettobelastung $540.000 - 240.000 = \textbf{300.000~€}$ [2~P.]

\emph{Entscheidung:} Die Wärmerückgewinnungsanlage ist um $480.000 - 300.000 =
\textbf{180.000~€}$ pro Jahr günstiger und sollte realisiert werden. [1~P.]\\
Kontrolle über die Vermeidungskosten: $540.000 \div 9.000 = 60$~€/t $<$ 80~€/t
Zertifikatspreis; Vorteil $9.000 \times (80 - 60) = 180.000$~€. [0,5~P.]

\medskip
\textbf{c) Beurteilung Emissionshandel vs. Auflage} \hfill [3~P.: je Argument 1~P.]

\emph{Vorteile des Emissionshandels} (zwei genügen):
\begin{itemize}
  \item \emph{Ökologische Treffsicherheit:} Die Gesamtmenge der Emissionen ist durch das
    Cap sicher begrenzt.
  \item \emph{Kosteneffizienz:} Vermieden wird dort, wo es am billigsten ist (wie bei
    Hagedorn mit 60~€/t); das Umweltziel wird volkswirtschaftlich zu geringsten Kosten
    erreicht.
  \item \emph{Dauerhafter Innovationsanreiz:} Jede eingesparte Tonne bringt Geld; bei
    einer Auflage fehlt der Anreiz, den Grenzwert zu unterschreiten.
\end{itemize}
\emph{Nachteil} (einer genügt): schwankende Zertifikatspreise erschweren die Planung;
Gefahr der Abwanderung in Länder ohne CO\textsubscript{2}-Preis (Carbon Leakage);
hoher Verwaltungs- und Kontrollaufwand (Messung, Berichtspflichten). Eine Auflage wirkt
dagegen direkt und sicher beim einzelnen Werk, ist aber unflexibel.
\end{loesung}


\end{document}
